\section{Correspondences, tautological classes, and extension problems}
\label{sec:part-c}
\providecommand{\Lie}{\operatorname{Lie}}
\providecommand{\tr}{\operatorname{tr}}
\providecommand{\Alt}{\operatorname{Alt}}
\providecommand{\sgn}{\operatorname{sgn}}
\providecommand{\Gr}{\operatorname{Gr}}
\providecommand{\id}{\operatorname{id}}

\subsection{C.1: a finite presentation for every cohomological tautological ring}

We use a Weil cohomology with coefficient field $L$ of characteristic zero, the usual Koszul commutativity constraint, and cohomological grading. Let $X$ be smooth, projective, geometrically connected, and of dimension $n>0$ over an algebraically closed field. Put $D=2n$, $A=H^*(X)$, and $\kappa(a)=\int_Xa$. We choose the top class $y$ with $\kappa(y)=1$. All the generators under consideration have degree divisible by $D$; we write $T_k[d]=T_k\cap H^{Dd}(X^k)$. For $n=0$, $X$ is a point and $T_k=L$ in degree zero, which we henceforth set aside.

\begin{definition}[Pointed partitions]
\label{c:def-pointed}
\status{Definition, with degrees verified}
Let $\mathcal P_k$ be the set of partitions of $\{0,1,\ldots,k\}$. The block containing $0$ is the distinguished block. For $P\in\mathcal P_k$, put $J=P(0)\setminus\{0\}$ and define
\[
 \varepsilon_P=\prod_{i\in J}y_i\prod_{B\in P,\,0\notin B}\delta_B,
 \qquad d(P)=k-|P|+1.
\]
Here $\delta_B$ is the pullback of the small diagonal in $X^B$ and $\delta_{\{i\}}=1$. Thus $\varepsilon_P\in T_k[d(P)]$. This describes the pointed and unpointed partial diagonals of manuscript 11:55 without duplicate notation.
\end{definition}

\begin{theorem}[C.1(b): the complete graph product rule]
\label{c:graph-rule}
\status{Proved}
Choose a tree on each block of $P$ and of $Q$, using the edges $\{0,i\}$ in a distinguished block, and take the union as a multigraph on $\{0,\ldots,k\}$. Write $b(C)=|E(C)|-|V(C)|+1$ for the cycle rank of a connected component. A component not containing $0$ contributes
\[
 \begin{cases}
 \delta_C,&b(C)=0,\\
 \chi(X)\prod_{i\in C}y_i,&b(C)=1,\\
 0,&b(C)\geq2,
 \end{cases}
\]
whereas the component containing $0$ contributes $\prod_{i\in C\setminus\{0\}}y_i$ if $b(C)=0$, and zero otherwise. Multiplying these contributions gives $\varepsilon_P\varepsilon_Q$, independently of the chosen trees. In particular, the $\varepsilon_P$ span $T_k$.
\end{theorem}

\begin{proof}
A tree of pairwise equalities imposes a small diagonal, since its edges impose independent diagonal conditions. In the algebra $A^{\otimes k}$ this also follows by composing the coevaluation for the perfect pairing $\kappa(ab)$. After the tree conditions have been imposed, each further edge in the same component is the self-intersection of the diagonal. Its pullback is the Euler element
\[
 e_A=\mu(\Delta_A)=\chi(X)y,
\]
where $\Delta_A\in A\otimes A$ is the coevaluation and $\mu$ is multiplication. To verify the last equality directly, choose homogeneous dual bases; the coefficient of $y$ is the supertrace of the identity, namely $\sum_r(-1)^r\dim H^r(X)$. Since $n>0$, $y^2=0$. Consequently a second excess edge gives zero. If the component already contains a point condition, an excess edge gives $y e_A=0$; the tree itself makes every vertex in that component equal to the point. The coevaluation identities show that changing the tree leaves the same diagonal class. Every product of generators is therefore either zero or a scalar times a pointed-partition class. These arguments establish the previously unchecked products of 11:55--57, including the cycle and point cases.
\end{proof}

We next give an actual finite system of generators for all relations, valid even when the restricted intersection pairing on $T_k$ is degenerate. Using only that restricted pairing would not suffice.

\begin{theorem}[C.1(a): all relations from the multilinear forms]
\label{c:matrix-presentation}
\status{Proved}
Choose a homogeneous basis $(e_a)$ of $A$, containing $e_0=1$. For fixed $d$, form a matrix $C_{k,d}$ whose columns are the partitions $P$ with $d(P)=d$ and whose rows are the basis tensors $e_{a_1}\otimes\cdots\otimes e_{a_k}$ of total degree $D(k-d)$. Its entry is
\begin{equation}
\label{c:contraction-matrix}
 (C_{k,d})_{\mathbf a,P}
 =\epsilon(\mathbf a,P)
 \prod_{B\in P,\,0\notin B}\kappa\Bigl(\prod_{i\in B}e_{a_i}\Bigr)
 \prod_{i\in P(0)\setminus\{0\}}[a_i=0].
\end{equation}
In each block the indices are in increasing order; $\epsilon$ is the Koszul sign of collecting the factors into those blocks. Then
\[
 \dim_L T_k[d]=\operatorname{rank}_L C_{k,d},\qquad
 \sum_Pc_P\varepsilon_P=0\ \Longleftrightarrow\ (c_P)_P\in\ker C_{k,d}.
\]
A finite presentation of $T_k$ consists of one variable $z_P$ for each $P\in\mathcal P_k$, the unit relation, all product relations of Theorem~\ref{c:graph-rule}, and one linear relation for each vector in a basis of each $\ker C_{k,d}$. Thus Gaussian elimination gives generators for the entire ideal of relations. The input is exactly the graded vector space, its unit and trace, and the multilinear forms $\kappa(u_1\cdots u_m)$ on positive-degree elements.
\end{theorem}

\begin{proof}
The projection formula for each diagonal gives
\[
 \int_{X^k}\varepsilon_P(e_{a_1}\otimes\cdots\otimes e_{a_k})
 =(C_{k,d})_{\mathbf a,P}.
\]
For a pinned coordinate, $\kappa(ye_a)=[a=0]$; for an unpinned block it is the integral of the product of its inputs. These are all factors, with precisely the displayed Koszul permutation sign. The full Poincar\'e pairing on $A^{\otimes k}$ is perfect, so its contraction with all rows detects the zero vector. This proves both the kernel and rank assertions. The graph relations reduce any polynomial to a linear combination of $z_P$; the matrix kernel is exactly the kernel remaining after this reduction. No nondegeneracy of the restricted pairing on the span of the $\varepsilon_P$ is used. Finally, multilinear forms involving a unit are obtained by deleting that unit, so the stated positive-degree data suffice.
\end{proof}

\begin{proposition}[The modified diagonal and the correction to R8]
\label{c:modified-diagonal}
\status{Proved; proposed generation by this relation alone refuted}
The class
\[
 M=\delta_{123}-\delta_{12}y_3-\delta_{23}y_1-\delta_{13}y_2
       +y_1y_2+y_2y_3+y_1y_3
 \in H^{4n}(X^3)
\]
vanishes if and only if $\kappa(uvw)=0$ for every three positive-degree classes $u,v,w$. The degree $H^{2n}$ stated in R8 is incorrect when $n$ denotes the geometric dimension. Neither $M$ nor its higher-arity analogues, without further dimension relations, generate all relations in every $T_k$.
\end{proposition}

\begin{proof}
Pair $M$ with $u\otimes v\otimes w$. If all three inputs have positive degree, all terms except the small diagonal vanish and the pairing is $\kappa(uvw)$. If at least one input is $1$, the terms cancel by inclusion and exclusion; multilinearity reduces every input to these cases. Poincar\'e duality proves the equivalence. For the generation assertion take $X=\mathbb P^1$ and $k=2$. There is the linear relation $\delta_{12}=y_1+y_2$; relations with at least three distinct indices are absent in this ring, while the universal self-intersection and point relations have polynomial degree at least two. Their homogeneous ideal cannot contain this nonzero linear polynomial. More generally, the dimension relations in Theorem~\ref{c:pair-ranks} below are additional indispensable relations. We distinguish this correction from the valid triple-product criterion in R8.
\end{proof}

\begin{theorem}[Dimensions for curves and for surfaces with $b_1=0$]
\label{c:pair-ranks}
\status{Proved}
For a curve of genus $g$, put $b=-2g$; for a smooth projective surface with $b_1=0$, put $b=b_2$. Define $q_{ij}=\delta_{ij}-y_i-y_j$. All relations are generated by
\begin{align}
 y_i^2&=0,&y_iq_{ij}&=0,&q_{ij}^2&=b\,y_iy_j,&q_{ij}q_{ik}&=y_iq_{jk},\label{c:pair-local}
\end{align}
and the following dimension relations, with all displayed indices distinct. For a surface they are
\begin{equation}
\label{c:orth-det}
 \sum_{\sigma\in S_{b+1}}\sgn(\sigma)
   \prod_{a=1}^{b+1}q_{i_a,j_{\sigma(a)}}=0;
\end{equation}
for a curve they are
\begin{equation}
\label{c:sympl-pf}
 \sum_{\text{perfect matchings }\mathcal M\text{ on }2g+2\text{ slots}}
       \prod_{\{i,j\}\in\mathcal M}q_{ij}=0.
\end{equation}
There is no sign in \eqref{c:sympl-pf}, because the symplectic signs are already included by the odd degree of $H^1$ and the Koszul constraint.

For a partition $\lambda$, write $f^\lambda$ for the number of standard Young tableaux of shape $\lambda$, explicitly $|\lambda|!$ divided by the product of the hook lengths. Put
\[
 r_b(m)=
 \begin{cases}
 \displaystyle\sum_{\lambda\vdash m,\,\ell(\lambda)\leq b} f^{2\lambda},&b\geq0\quad\text{(surface)},\\[4pt]
 \displaystyle\sum_{\lambda\vdash m,\,\lambda_1\leq g} f^{2\lambda},&b=-2g\quad\text{(curve)},
 \end{cases}
 \qquad r_b(0)=1.
\]
Then the answer in every degree is
\begin{equation}
\label{c:dimension-pair}
 \dim T_k[d]=
 \sum_{\substack{m\geq0,\ a=d-m\geq0\\a+2m\leq k}}
 \frac{k!}{a!(2m)!(k-a-2m)!}\,r_b(m).
\end{equation}
The curve formula includes $g=0$, when $r_0(m)=0$ for $m>0$ and the answer is $\binom{k}{d}$.
\end{theorem}

\begin{proof}
In both cases $A=L1\oplus W\oplus Ly$, with $W=H^1$ for a curve and $W=H^2$ for a surface. The product on $W$ is a perfect alternating, respectively symmetric, pairing with values in $Ly$. Direct contraction gives \eqref{c:pair-local}; the superdimension of $W$ gives $b$. In particular the positive triple products vanish. Repeated use of those relations reduces a monomial to $y_I\prod_{\{i,j\}\in\mathcal M}q_{ij}$, with disjoint indices. Factorwise cohomological degrees separate different point sets $I$ and different sets of paired vertices. For a fixed set of $2m$ paired vertices, the remaining problem is exactly the span of the pairing tensors in $W^{\otimes2m}$.

For completeness, the matching representation of $S_{2m}$ decomposes without multiplicity into the Specht modules indexed by the even-row shapes $2\lambda$. Its contraction matrix has nonzero eigenvalue exactly for $\ell(\lambda)\leq b$ in the orthogonal case and $\lambda_1\leq g$ in the symplectic case. The orthogonal calculation and the argument that every kernel Specht module is generated by the first determinant are stated and proved in \citet[Proposition 3.5 and \S3.6, pp.~508--510]{c_yin2015}; the symplectic calculation, including the all-plus convention and generation by \eqref{c:sympl-pf}, is worked out in \citet[\S6, equation (2), and the argument preceding Corollary 6.4]{c_tavakol2018}. These calculations concern a matrix of pairings and therefore apply to the present $W$ independently of whether it comes from a K3 surface or a hyperelliptic curve. To see the generation step, a kernel Young diagram has either a column of length $b+1$ or a row of length $2g+2$; antisymmetrizing that column, respectively symmetrizing that row with the Koszul signs, supplies \eqref{c:orth-det}, respectively \eqref{c:sympl-pf}. Decomposing the larger stabilizer into cosets reduces every larger kernel generator to its pullbacks and products. The hook formula gives $r_b(m)$. Finally, choose the $a$ point slots and the $2m$ pairing slots; the number of choices is the multinomial coefficient in \eqref{c:dimension-pair}. Summing gives the formula and proves completeness of the relations.
\end{proof}

\begin{proposition}[Explicit dimension algorithms for projective spaces and abelian varieties]
\label{c:other-dimensions}
\status{Proved}
The following finite integer matrices determine every dimension by ordinary exact rational rank; in particular they give a uniform computable answer for every $n,g,k,d$, with no geometric unknown in their entries.

For $X=\mathbb P^n$, the columns are $P\in\mathcal P_k$ with $d(P)=d$, and the rows are $(a_1,\ldots,a_k)\in\{0,\ldots,n\}^k$ with $\sum a_i=nd$. The entry is $1$ precisely when
\[
 a_i=n\quad(i\in P(0)\setminus\{0\}),\qquad
 \sum_{i\in B}a_i=n(|B|-1)\quad(0\notin B),
\]
and is zero otherwise. Its rank is $\dim T_k[d]$.

For an abelian variety of dimension $g>0$, use the exterior algebra on the $2gk$ symbols $e_{i,a}$, ordered lexicographically. In this algebra put
\[
 y_i=\prod_{a=1}^{2g}e_{i,a},\qquad
 \delta_{ij}=\prod_{a=1}^{2g}(e_{i,a}-e_{j,a}).
\]
For each $P$ multiply these expressions along any of its trees and its pinned vertices. The coefficient matrix in the exterior-monomial basis of degree $2gd$ has entries in $\Z$ and rank $\dim T_k[d]$. These ranks depend only on $g,k,d$.
\end{proposition}

\begin{proof}
For projective space, $H^*(X)=L[h]/(h^{n+1})$, $\int h^n=1$, and the small diagonal on a block $B$ is the sum of the monomials satisfying the displayed condition, each with coefficient one. This follows either by pairing with monomials or by induction along a tree. For an abelian variety, $H^*(X)=\bigwedge H^1(X)$. Choose a volume-normalized basis; the diagonal formula follows from the difference map $(x,y)\mapsto x-y$, whose pullback on $H^1$ is $e\mapsto e\otimes1-1\otimes e$, and whose inverse image of the origin is the diagonal. Expansion in the basis is therefore an explicit expression for every $\varepsilon_P$. In both cases the coefficient matrix has precisely these classes as columns, so its rank is their dimension. These matrices are distinct from an unevaluated Chow-theoretic condition: their entries require only finite integer arithmetic.
\end{proof}

The exact-arithmetic program \texttt{code/part\_c/check\_tautological.py} implements the Frobenius algebra construction independently of the display formulas. It verified 2,286 graph products and the curve/surface rank formula in the finite ranges recorded in its JSON output. For example, the dimensions for $k=4$, in normalized degrees $d=0,\ldots,4$, are
\[
\begin{array}{c|ccccc}
 X&0&1&2&3&4\\\hline
 \mathbb P^1&1&4&6&4&1\\
 \mathbb P^2&1&10&19&10&1\\
 \mathbb P^3&1&10&25&14&1\\
 \text{curve }g=1&1&10&20&10&1\\
 \text{curve }g=2&1&10&21&10&1
\end{array}
\]
The non-palindromic row for $\mathbb P^3$ is one reason not to replace the full pairing in Theorem~\ref{c:matrix-presentation} by an assumed perfect tautological pairing.

\subsection{C.1(c): Chow lifting, and the exact K3 obstruction}

Write $T_k^{\CH}(X,o)$ for the rational Chow subring generated by the same diagonals and a degree-one zero-cycle $o$. The requested lifting of rational cohomological relations is exactly the injectivity of the cycle class map on this subring. For a general Weil coefficient field, this rational statement is distinguished from injectivity after extending Chow coefficients to that field; the latter coefficient convention is addressed separately in Corollary~\ref{c:finite-weil-criterion}. In Chow groups the graph proof has Euler class $c_n(T_X)$ in place of $\chi(X)o$; this replacement itself is a condition to verify, not a consequence of the cohomological computation.

\begin{theorem}[The abelian and hyperelliptic cases]
\label{c:chow-positive}
\status{Known, with the required implication proved}
For an abelian variety with $o$ its origin, every cohomological relation in the rings of C.1(a) holds in $T_k^{\CH}(X,o)$ for every $k$. The same conclusion holds for a hyperelliptic curve with $o$ a Weierstrass point (and for an elliptic curve with its origin and for $\mathbb P^1$).
\end{theorem}

\begin{proof}
For an abelian variety, \citet[Theorem 6.2.5 and Corollary 6.2.6]{c_osullivan2011} identifies the algebra of symmetrically distinguished cycles with the numerical Chow algebra, compatibly with products and with pullbacks and pushforwards by homomorphisms. The origin is the pushforward of $1$ by the zero homomorphism from the zero-dimensional abelian variety. Each diagonal is the pushforward of $1$ by a diagonal homomorphism. Their products and pullbacks are consequently symmetrically distinguished. A homologically zero cycle is numerically zero; the stated isomorphism then forces a distinguished such cycle to vanish in rational Chow groups. This establishes the assertion for precisely the requested generators, without assuming the Hodge or Tate conjecture.

For a hyperelliptic curve of genus $g\geq2$, $K_X=(2g-2)o$ in $\Pic(X)\otimes\Q$. The modified diagonal vanishes by \citet[Proposition 4.8 and its relation with the cycle on $X^3$ in \S4]{c_grossschoen1995}; equivalently one may use the explicit divisor calculation in that proposition followed by the finite symmetric-power push--pull. The remaining relations are precisely \eqref{c:pair-local} and \eqref{c:sympl-pf}. \citet[\S6, equation (2), and Corollary 6.4]{c_tavakol2018} proves them and the perfectness of every tautological intersection pairing on the fixed curve. A homologically zero tautological cycle pairs to zero with every complementary tautological cycle, hence is zero by that perfectness. The identification of Tavakol's class $a_i=K_i/(2g-2)$ with $o_i$ verifies the change of generators. Genus one is covered both by the abelian argument and by Tavakol's Example 6.2. Genus zero is the polynomial Chow presentation of projective space.
\end{proof}

\begin{theorem}[Classification of the allowed zero-cycles on an abelian variety]
\label{c:abelian-point-classification}
\status{Proved over every algebraically closed field}
Let $A/F$ be an abelian variety over an algebraically closed field of any characteristic, and let $o\in\CH_0(A)_\Q$ have degree one. The following are equivalent:
\begin{enumerate}[label=(\roman*)]
\item all the cohomological point-and-diagonal relations hold in $T_k^{\CH}(A,o)$ for every $k$;
\item $\Delta_*o=o\boxtimes o$ in $\CH_0(A\times A)_\Q$;
\item $o=[a]$ in $\CH_0(A)_\Q$ for a point $a\in A(F)$.
\end{enumerate}
The point is unique modulo torsion in $A(F)$. In particular, specifying the origin is sufficient but is not necessary for the requested lifting.
\end{theorem}

\begin{proof}
The equality in (ii) is the cohomological relation $\delta_{12}y_1=y_1y_2$, since both classes are degree-one zero-cycles on $A^2$. Thus (i) implies (ii). If (iii) holds, simultaneous translation by $a$ on every factor preserves all partial diagonals and carries the origin class to $[a]$. Theorem~\ref{c:chow-positive} proves (i).

We prove the remaining implication algebraically, so that no characteristic-zero restriction is hidden. Put $g=\dim A$ and write $\widehat A$ for the dual abelian variety. For the normalized Poincar\'e bundle $\mathcal P$ on $A\times\widehat A$, define the rational Chow Fourier transform by
\[
 \mathcal F_A(z)=(p_2)_*\bigl(p_1^*z\,\exp(c_1(\mathcal P))\bigr).
\]
It is an isomorphism and $\mathcal F_{\widehat A}\mathcal F_A=(-1)^g[-1]^*$. These identities hold over the arbitrary algebraically closed field of \citet[Theorem 2.2, p.~156, with the field convention on p.~154]{c_mukai1981}: apply Grothendieck--Riemann--Roch to his Fourier--Mukai identity, using the trivial tangent bundle of an abelian variety. The Chern character identifies rational $K$-theory with rational Chow groups. We will also use the directly checked Poincar\'e identities below; over $\C$ they are formulas F1--F3 and Proposition 1 of \citet[pp.~647--648]{c_beauville1986}.

Set $z=\mathcal F_A(o)$. Its codimension-zero part is $\deg(o)=1$. Fourier transformation of (ii), using the defining Poincar\'e bundle and the projection formula, gives
\[
 m_{\widehat A}^*z=p_1^*z\,p_2^*z.
\]
Indeed Fourier transformation changes $\Delta_*$ into pullback by the dual homomorphism $m_{\widehat A}$, and changes an exterior product into the exterior product of the transforms. The ideal of positive-codimension classes is nilpotent, so the finite series $D=\log z$ is defined and satisfies
\[
 m_{\widehat A}^*D=p_1^*D+p_2^*D,
 \qquad [n]^*D=nD\quad(n\geq1).
\]

Write $D=\sum_{r=1}^gD_r$ by codimension. We verify that $D_r=0$ for $r\geq2$ without assuming the conjectural vanishing of any negative Beauville piece. Choose an integer $n\geq2$ prime to the characteristic. If $w_q$ is the codimension-$q$ part of $\mathcal F_{\widehat A}(D_r)$, its defining formula uses the power $c_1(\mathcal P)^{g-r+q}$. The identity $(1\times[n])^*\mathcal P=\mathcal P^{\otimes n}$ and flat base change therefore give
\[
 [n]^*w_q=n^{g-r+q}w_q.
\]
Fourier functoriality for the isogeny $[n]$, applied to $[n]^*D_r=nD_r$, gives $[n]_*w_q=nw_q$. Since $[n]_*[n]^*=n^{2g}$, a nonzero $w_q$ requires
\[
 n^{2g}=n^{g-r+q+1},\qquad q=g+r-1.
\]
The inequality $q\leq g$ forces $r\leq1$. Fourier inversion now gives $D=D_1$. On the N\'eron--Severi group $[n]^*$ acts by $n^2$, whereas it acts on $D_1$ by $n$. Hence $D_1\in\Pic^0(\widehat A)\otimes\Q$.

Biduality identifies $\Pic^0(\widehat A)$ with $A(F)$. This group is divisible: multiplication by every positive integer is a surjective isogeny and is surjective on points over an algebraically closed field, also when the integer is divisible by the characteristic. Thus there is $a\in A(F)$ with $D_1=c_1(\mathcal P_a)$ in rational Chow groups. Directly from its definition, $\mathcal F_A([a])=\exp(c_1(\mathcal P_a))=z$. Fourier inversion proves $o=[a]$. Finally $[a]=[b]$ in rational Chow groups if and only if their Albanese classes differ by torsion: the forward implication follows from Albanese, and the reverse from the identical Fourier exponentials of torsion-equivalent points. This proves the uniqueness assertion and the theorem.
\end{proof}

\begin{proposition}[A finite criterion for curves]
\label{c:curve-chow-criterion}
\status{Proved over every algebraically closed field}
Let $C$ be a smooth projective connected curve and $o\in\CH_0(C)_\Q$ have degree one. All point-and-diagonal cohomological relations hold in all its powers if and only if $M_o=0$ in $\CH^2(C^3)_\Q$. If $g\geq2$, this forces
\[
 o=\frac{K_C}{2g-2}\quad\text{in }\CH^1(C)_\Q.
\]
For $g=0$ or $1$ every degree-one rational zero-cycle has the required property. Over $\C$, the genus-at-least-two criterion is also the multiplicativity criterion in \citet[Propositions 3.1 and 3.5]{c_fulaterveervial2021}.
\end{proposition}

\begin{proof}
Necessity follows because $M_o$ is always cohomologically zero for a curve. Suppose it vanishes. Intersecting with $p_1^*o$ and pushing to the last two factors gives
\[
 \Delta_*o=o\boxtimes o.
\]
Intersecting instead with $\delta_{12}$ and pushing to the third factor gives
\[
 c_1(T_C)=(2-2g)o.
\]
In the second calculation, the small-diagonal term contributes $c_1(T_C)$, the term $\delta_{12}o_3$ contributes $(2-2g)o$, and the four nonzero remaining point terms cancel in pairs. These identities and the expansion of $M_o=0$ give all the local relations \eqref{c:pair-local} for $q_{ij}=\delta_{ij}-o_i-o_j$, with parameter $b=-2g$.

We check the remaining matching relation, including for a rational degree-one zero-cycle. Choose a base point $e$ and the Abel--Jacobi map $C\to J(C)$. The Jacobian's principal autoduality and the theorem of the square identify the pullback of its normalized symmetric Poincar\'e class, up to a common sign, with $\delta_{12}-e_1-e_2$. Translation of the Abel--Jacobi map by $z\in J(C)$ adds the corresponding algebraically trivial divisor to each factor. Since $J(C)(k)$ is divisible, we may choose $z$ so that the new pullback is $q_{12}$ in rational Chow groups, up to the same sign. Thus every matching monomial in the $q_{ij}$ is the pullback of a matching monomial in Poincar\'e divisors on a power of $J(C)$.

Let $\theta$ be a symmetric principal polarization on $J(C)$ and use integer parameters $n_1,\ldots,n_{2g+2}$. Pull back $\theta^{g+1}=0$ along the homomorphism $(x_i)\mapsto\sum n_ix_i$. The theorem of the square gives a polynomial identity whose quadratic form is
\[
 \sum_i n_i^2\theta_i+\sum_{i<j}n_in_jP_{ij}.
\]
The coefficient of $n_1\cdots n_{2g+2}$ in its $(g+1)$st power is $(g+1)!$ times the sum over all perfect matchings of the products of $P_{ij}$. A rational polynomial vanishing at every integer tuple has every coefficient zero. Pulling back to $C^{2g+2}$ proves precisely \eqref{c:sympl-pf}. Theorem~\ref{c:pair-ranks} now proves sufficiency in every power. These are algebraic arguments valid also in characteristic two.

For genus zero, $\CH_0(C)_\Q=\Q$. For genus one, choose an origin. The degree-zero Chow group is its Jacobian group, so divisibility shows that every degree-one rational zero-cycle equals a point class. Apply Theorem~\ref{c:abelian-point-classification}. The formula for $o$ when $g\geq2$ follows from the second displayed identity.
\end{proof}

\begin{proposition}[The exact additional obstruction for regular surfaces]
\label{c:surface-chow-criterion}
\status{Proved equivalence; does not assume finite-dimensionality}
Let $S/\C$ be a smooth projective surface with $b_1(S)=0$, and let $o$ have degree one. All the point-and-diagonal cohomological relations lift in all powers if and only if $M_o=0$ and $h(S)$ is Kimura finite-dimensional. If they hold, $o$ is uniquely determined: for every ample divisor $H$,
\[
 o=\frac{H^2}{\deg H^2}.
\]
\end{proposition}

\begin{proof}
The projector $\pi_2=\Delta-o\times S-S\times o$ has realization $H^2(S)$; the other two summands are Tate. Write $b=b_2(S)$. The alternating cross-matching determinant on $2(b+1)$ factors is $(b+1)!$ times the alternating projector on $h^2(S)^{\otimes(b+1)}$. Therefore the lifting of all relations implies $\bigwedge^{b+1}h^2(S)=0$, hence finite-dimensionality. It also implies $M_o=0$, since three positive-degree classes have zero product on such a surface.

Conversely, let $M_o=0$. Acting on $o$, on the diagonal, and on $H\boxtimes H$ gives respectively
\[
 \Delta_*o=o\boxtimes o,\qquad c_2(T_S)=(b+2)o,
 \qquad H^2=(\deg H^2)o.
\]
The same expansions as for the curve give all local relations \eqref{c:pair-local}, now with $b=b_2(S)$. Finite-dimensionality supplies the remaining alternating determinant: the realization is purely even of dimension $b$, so its Kimura exterior dimension is $b$. Theorem~\ref{c:pair-ranks} says that the local identities and these determinants generate every cohomological relation. This proves sufficiency and uniqueness. This formal argument is consistent with the larger divisor-generated criterion in \citet[Proposition 4.5]{c_fulaterveervial2021}; it explicitly applies to the narrower ring asked for here.
\end{proof}

\begin{theorem}[A finite cycle criterion when the cohomology is even]
\label{c:finite-even-criterion}
\status{Proved in the pure-even case}
Let $X/\C$ be smooth projective and connected, with $H^{\rm odd}(X,\Q)=0$,
let $n=\dim X>0$, and let $o\in\CH^n(X)_\Q$ have degree $1$.
Put $d=\dim_\Q H^*(X,\Q)=\chi(X)$.  There is a finite list of rational
tensor expressions $P_1,\ldots,P_s$, constructed by the procedure below
from the rational cup-product table and trace of $X$, for which the
following criterion is necessary and sufficient for all the relations
of Theorem~\ref{c:matrix-presentation} to hold in rational Chow groups
on every power:
\begin{enumerate}[label=(\roman*)]
\item $c_n(T_X)=d\,o$ and $\Delta_*o=o\boxtimes o$;
\item $\bigwedge^{d+1}h(X)=0$ in rational Chow motives;
\item $P_a(h(X),\mu,1,\operatorname{tr},o,\beta^{-1})=0$ for
      $1\leq a\leq s$.
\end{enumerate}
Here $\mu$ is the small-diagonal multiplication,
$\operatorname{tr}:h(X)\to\Q(-n)$ is the trace,
$\beta=\operatorname{tr}\circ\mu$, and
$\beta^{-1}:\Q(-n)\to h(X)\otimes h(X)$ is its inverse pairing.
Every condition is the vanishing of a specified finite linear
combination of correspondence diagrams.  Subject to (i), each can
equivalently be written as a finite list of cycles in the requested
point-and-diagonal rings.  The construction uses finite-dimensional
polynomial algebra, so is a finite criterion, rather than one separate
test for each power of $X$.
\end{theorem}

\begin{proof}
Write $A=H^*(X,\Q)$ as an ordinary vector space; the pure-even
hypothesis makes its tensor symmetry the ordinary one.  Retain the
Tate degree by a one-dimensional auxiliary representation $L$, and
put $G=\operatorname{GL}(A)\times\mathbb G_m$, with the second factor
acting on $L$.  The data
\[
 t=(\mu,1,\operatorname{tr},o,\beta^{-1})
\]
form one rational point of the finite-dimensional rational
$G$-representation
\[
 W=(A\otimes A^\vee\otimes A^\vee)\oplus A
   \oplus(A^\vee\otimes L^n)
   \oplus(A\otimes L^{-n})
   \oplus(A\otimes A\otimes L^{-n}).
\]
In the point $t$, the coordinate denoted by $o$ is the cohomology class
of the degree-one zero-cycle.  Let $I\subset\Q[W]$ be the ideal of
the Zariski closure of $G\cdot t$.  We describe all steps producing
the finite list.

First compute generators of $I$ by elimination: write the orbit map
in matrix entries of $g\in\operatorname{GL}(A)$ and a nonzero scalar
$z$, introduce inverses with equations $u\det(g)=1$ and $vz=1$, and
eliminate $g,u,z,v$ from the graph ideal.  The resulting elimination
ideal is the ideal of the image closure.  Hilbert's basis theorem
makes the resulting list finite.  Enlarge its span to a finite-dimensional
$G$-stable subspace $E\subset I$ generating $I$ as an ideal.
This enlargement is finite because the action on the polynomial ring
is locally finite and preserves the finite degree bound of the chosen
generators.  It can be computed by closing their span under the matrix
units of the Lie algebra and the second torus action; $G$ is connected
in characteristic zero.

The inclusion $E\hookrightarrow\Q[W]$ is equivalently a polynomial
covariant
\[
 p_E\in(E^\vee\otimes\Q[W])^G.
\]
Choose a split $G$-embedding of $E^\vee$ into a finite direct sum $T$
of mixed tensor powers of $A,A^\vee,L,L^\vee$.  Such a split embedding
exists because rational representations of $G$ are semisimple and
are generated by its defining representations and their duals.
For an explicit terminating search, enumerate these finite direct
sums and compute the rational intertwiner spaces.  Enumerate their
rational elements until one has full rank on $E^\vee$, and then
solve the linear intertwiner equations for a retraction.  The
existence of a split embedding proves termination.  The components of the image of
$p_E$ in $T\otimes\Q[W]$ are the covariants we denote by $P_a$.

We recall the two representation-theoretic facts needed here.
Over $\C$, averaging a Hermitian form over
$U(d)\times U(1)$ gives invariant complements.  This compact group
is Zariski dense in $G_\C$, so the complements are $G_\C$-stable.
The intertwiner and splitting equations have rational coefficients;
Gaussian elimination descends their solvability to $\Q$.
Every finite-dimensional rational representation embeds, by its
matrix coefficients, into finitely many copies of $\Q[G]$.
The relevant coefficient functions have bounded polynomial degree
and bounded powers of $\det^{-1}$ and the second torus coordinate.
They therefore occur in mixed tensor powers, using symmetric and
exterior projectors.  Complete reducibility makes the embedding
split.  This proves the existence statements just used.  It also
proves exactness of invariants for the locally finite rational
modules below: place a preimage and its translates in a
finite-dimensional submodule and split the resulting surjection.
The exact references are \citet[Theorem~4.14, pp.~88--89;
Theorem~22.42, pp.~478--479; Remark~12.53(a), p.~248,
and the Reynolds operator on p.~249]{c_milneag2022}.

We explain why these covariants are tensor expressions of the promised
kind.  Polarize a homogeneous polynomial covariant in the finitely
many tensor coordinates of $W$.  Its coefficients are invariants in
mixed tensor powers of $A$ and $A^\vee$.  The first fundamental
theorem says that these are sums of complete contractions, with any
uncontracted indices supplying the target tensor.  The only linear
relations among such contractions are generated by the
$(d+1)$-fold antisymmetrizer and by the scalar loop value $d$.
These are precisely \citet[Theorems~1.2--1.5, pp.~4--5]{c_procesi2021tensor};
polarization and restitution apply because the characteristic is zero.
This use is of the classical tensor theorems, and does not apply a
matrix-only statement to arbitrary tensor variables without that step.
The extra $\mathbb G_m$ factor merely selects the matching Tate
degrees.  Thus each $P_a$ is a finite rational linear combination of
the prescribed diagrams, and every universal identity in their
ordinary $d$-dimensional interpretation follows from the stated
antisymmetrizer and loop relation.

Now let $F$ be any polynomial covariant with target a mixed tensor
space $T'$ that vanishes at $t$.  Equivariance implies that it
vanishes on $G\cdot t$ and its closure.  Hence
\[
 F\in(T'\otimes I)^G.
\]
The multiplication map $E\otimes\Q[W]\twoheadrightarrow I$ is
$G$-equivariant.  Since $G$ is reductive over $\Q$, taking invariants
is exact, and $F$ lifts to
\[
 \widetilde F\in(T'\otimes E\otimes\Q[W])^G.
\]
Contracting its $E$ index against $p_E$ gives $F$ exactly.  Extend
that contraction through the chosen split embedding into $T$.
The first fundamental theorem again writes $\widetilde F$ as a
diagram.  Consequently the tensor ideal of all relations at $t$
is generated by the finite covariants $P_a$, together with the
$(d+1)$-fold antisymmetrizer and the loop relation.  No reductivity
of the stabilizer of $t$ has been assumed or used.

Evaluate this presentation in Chow motives.  The scalar loop value
is $\chi(X)=d$ by the categorical trace formula.  Condition (ii)
annihilates the dimension relation, and condition (iii) annihilates
the finite orbit-ideal generators.  Therefore every cohomological
diagram relation vanishes in Chow, in particular every relation in
the requested rings.  This proves sufficiency.
The additional identities for the auxiliary line,
$\bigwedge^2 L=0$ and $\operatorname{tr}(\id_L)=1$, already hold
for the fixed even invertible Tate object $L=\Q(-1)$; they are
not additional conditions on $(X,o)$.

Conversely suppose all the requested relations hold in Chow.
The cohomological relations
$\Delta\cdot o_1=o_1o_2$ and
$\Delta^2=d\,o_1o_2$ hold there.  The first is the second assertion
of (i); pushing the second to $X$ gives its first assertion.
The alternating projector in (ii) is a rational linear combination
of permutation graphs on $X^{d+1}$, each a product of pairwise
diagonals in $X^{2(d+1)}$.  Its cohomology is zero, so (ii) follows.

Finally the correspondence diagrams used above do not enlarge the
vanishing test after (i).  The Frobenius identities reduce each
connected diagram to a small diagonal with some point insertions
and some handles.  A handle inserts $c_n(T_X)$; a point insertion
inserts $o$.  By (i) both reduce to scalar multiples of $o$, two
such insertions vanish for dimension reasons, and a single insertion
on a small diagonal is the external product of the point classes.
Closed components give their degrees.  Thus every diagram reduces
to the pointed partitions of Definition~\ref{c:def-pointed}.
This is the same graph reduction as Theorem~\ref{c:graph-rule},
now including contractions and using the two verified Chow
identities.  Each $P_a$ has zero cohomology and reduces to a relation
of the requested kind; it consequently vanishes in Chow.  This
proves necessity and the finite criterion.
\end{proof}

\begin{remark}[Scope of the finite criterion]
\label{c:finite-criterion-scope}
\status{Proved scope clarification; the arbitrary-pair classification remains unresolved}
Theorem~\ref{c:finite-even-criterion} applies to arbitrary cup-product
algebras with pure-even cohomology, not only to surfaces.  It gives
a finite algebraic construction of the cycle conditions, without
claiming to decide their rational equivalence by computation.
For mixed cohomology the direct use of
$\operatorname{GL}_{a|b}$ would require exactness of its invariant
functor, which fails.  Theorem~\ref{c:finite-super-criterion}
replaces that step by projective lifting and an external algebraic
parity correspondence.  Corollary~\ref{c:finite-weil-criterion}
also separates finite existence over an arbitrary Weil coefficient
field from effective construction with specified coefficient data.
Theorem~\ref{c:finite-unconditional-criterion} subsequently removes
the sign hypothesis by first passing to scalar forms in the
covariant polynomial category.  These finite criteria do not
decide the resulting Chow vanishings.
The unrestricted classification of all $(X,o)$ therefore remains
distinct from the specific K3 obstruction proved below.
\end{remark}


\begin{proposition}[Products and the failure of the tautological parity converse]
\label{c:product-parity}
\status{Product assertion proved; the proposed parity converse refuted}
Let $(X,o_X)$ and $(Y,o_Y)$ be pointed smooth projective complex
varieties for which all the requested point-and-diagonal relations
hold in rational Chow groups on every power.  Then the same is
true of $(X\times Y,o_X\boxtimes o_Y)$.
Nevertheless, this property does not imply that cohomological
parity is represented by a correspondence in the tensor category
generated by the marked Frobenius diagrams.  A counterexample is
\[
 X=E\times\mathbb P^1,\qquad
 E:\ y^2z=x^3-xz^2,\qquad
 o=[(0:1:0)]\boxtimes[\infty].
\]
Its parity correspondence is algebraic, but lies outside that
diagram category.
\end{proposition}

\begin{proof}
Under $(X\times Y)^k=X^k\times Y^k$, each marked diagonal is
the external product of the two corresponding marked diagonals.
Let $A_k$ and $B_k$ be the rational Chow tautological rings of the
two factors.  The two injective cycle maps and the rational Betti
K\"unneth isomorphism make
\[
 A_k\otimes_\Q B_k\longrightarrow
 H^*(X^k,\Q)\otimes_\Q H^*(Y^k,\Q)
\]
injective.  Its factorization through the external-product map
to $\CH^*(X^k\times Y^k)_\Q$ shows that the cycle map is
injective on the image of that map.  The requested product ring
is contained in this image, proving the first assertion.

For the displayed pair, the elliptic and projective-line cases
were proved in Theorem~\ref{c:chow-positive}.
Also $c_2(T_X)=0=\chi(X)o$, because $T_E$ is trivial, and
$\Delta_*o=o\boxtimes o$ because $o$ is an actual point.
The Frobenius graph reduction therefore applies to every
contracted diagram, and its degree-zero endomorphisms of $h(X)$
are contained in
\[
 T^{\CH}_2(X,o)\cap\CH^2(X^2)_\Q
   =\operatorname{span}_\Q\{\Delta_X,o\times X,X\times o\}.
\]
The latter two correspondences act only on $H^0$ and $H^4$,
respectively.  Thus any correspondence in this span acts by the
same scalar on all of $H^1,H^2,H^3$.  Both $H^1$ and $H^2$ are
nonzero, whereas parity must act as $-1$ on the former and $+1$
on the latter.  This is impossible.  Explicitly, the three diagonal
matrices of this span have rank three as vectors of endomorphism
coefficients, and adjoining parity raises the rank to four; the
exact computation is recorded with the accompanying code.

On the other hand, the graph of
$[-1]_E\times\id_{\mathbb P^1}$ acts as parity on cohomology,
so parity is algebraic in this example.  The counterexample
refutes only the proposed implication that parity belongs to the
\emph{narrow diagram category}.  It neither refutes algebraicity
of parity nor rules out a different finite mixed-parity criterion.
\end{proof}


\begin{lemma}[Rational projective mixed tensors]
\label{c:super-projective-tensors}
\status{Proved}
Let $A=\Q^{a|b}$ be a finite supervector space, let $L$ be an even
line, and put $G=\operatorname{GL}(A)\times\mathbb G_m$.
We use even morphisms and the category $\mathcal R$ of rational
$G$-supermodules on which
$\epsilon=(\operatorname{diag}(1_a,-1_b),1)$ acts as intrinsic parity.
Every finite object $E$ of $\mathcal R$ is a quotient of a finite
projective-injective object $P$ of $\mathcal R$.  Such $P$, its dual,
and a fixed nonzero projective object $P_0$ are direct summands of
finite sums of mixed tensor powers of $A,A^\vee,L,L^\vee$.
The modules, quotient maps, mixed tensor embeddings and retractions
can be constructed over $\Q$ by terminating finite procedures.
Projectivity also applies to surjections of locally finite rational
modules, including the polynomial modules used below.
\end{lemma}

\begin{proof}
Write $G_0=\operatorname{GL}_a\times\operatorname{GL}_b\times
\mathbb G_m$, $\mathfrak g=\operatorname{Lie}G$, and
$B=U(\mathfrak g_0)$.  We take
\[
 P=U(\mathfrak g)\otimes_B E|_{G_0},\qquad
 P\longrightarrow E,\quad v\otimes e\longmapsto ve.
\]
The odd part of $\mathfrak g$ is
$\operatorname{Hom}(\Q^a,\Q^b)\oplus
\operatorname{Hom}(\Q^b,\Q^a)$.
The PBW theorem makes $P$ finite, of total dimension
$2^{2ab}\dim E$.  Its natural rational $G_0$-action, together with
the compatible $\mathfrak g$-action, integrates to an algebraic
$G$-action over $\Q$.  Here we use
\citet[Corollary~3.6, p.~9; Proposition~5.4 and its proof,
pp.~23--24; Theorem~5.8 and Corollary~5.10, p.~26]{c_masuokashibata2017super}.
The identity $\operatorname{hy}(G)=U(\mathfrak g)$ in characteristic
zero is stated on p.~22 there.  In the proof of Proposition~5.4,
PBW splits off the finite exterior factor, the rational even action
supplies its coefficient functions, and injectivity of the pairing
with the hyperalgebra proves the coassociativity of the resulting
coaction.  Thus the statement applies over $\Q$, not only after
extension to $\C$.  The induced module still satisfies the
$\epsilon$-parity condition.

Induction is left adjoint to restriction to $G_0$.  The latter has
semisimple finite rational representations, so $P$ is projective.
In this case induction and coinduction coincide, with no character
or parity twist.  We give the reason to check this point over $\Q$:
\[
 \det(\mathfrak g_1)
 =\det(\Q^a)^{-b}\det(\Q^b)^a
  \det(\Q^a)^b\det(\Q^b)^{-a}=1,
 \qquad \dim\mathfrak g_1=2ab.
\]
Choose an ordered odd PBW basis.  Projection to its top exterior
coefficient gives a $B$-bimodule map $U(\mathfrak g)\to B$:
the displayed determinant identity is precisely the vanishing of
the possible even-character twist, and the top degree is even.
Multiplication followed by this map is a nondegenerate relative
pairing.  Complementary exterior monomials pair to $\pm1$, while
the commutator correction terms are triangular in exterior degree.
It identifies $U(\mathfrak g)$ with its $B$-linear dual with the
induction/coinduction bimodule structures.  Tensoring with $E$
gives $\operatorname{Ind}(E)\simeq\operatorname{Coind}(E)$.
Right adjunction and semisimplicity over $G_0$ now make $P$
injective as well.  We may take $P_0=\operatorname{Ind}(\Q)$.

For completeness, the tensor-generation assertion can be checked
without assuming $a\geq b$.  Matrix coefficients and the coaction
embedding express every finite rational superrepresentation as a
subquotient of finite sums of mixed tensors and their parity shifts.
Indeed $\Q[G]$ is generated by the entries of the standard
supermatrix, its inverse, and the two powers of the last torus
coordinate.  Its usual localization by the determinants of the two
even blocks has the same generators: an inverse diagonal block is
recovered from the corresponding block of the full inverse by a
finite geometric series in the odd, hence nilpotent, correction.
Every finite set of coefficient functions therefore comes from
bounded mixed tensor degrees.  To retain the correct category,
the operator consisting of intrinsic parity times $\rho(\epsilon)$
splits the full superrepresentation category into $\mathcal R$ and
its parity shift.  It is $+1$ on the mixed tensors in question and
$-1$ on their parity shifts.  Taking its $+1$ summand shows that an
object of $\mathcal R$ is already a subquotient of the unshifted
mixed tensors.  This also explains the category convention in
\citet[section~2, pp.~3--4]{c_heidersdorf2017mixed};
his Theorem~1.1, p.~2, includes the assertion that every projective
is a mixed tensor.

Write $P=U/V$ with $U\subset T$, where $T$ is a finite sum of mixed
tensors.  Projectivity splits $U\twoheadrightarrow P$, so $P$
embeds in $T$; injectivity retracts that embedding.  The dual is
then a summand of $T^\vee$.  These operations are rational.  For an
explicit terminating search, we enumerate finite mixed tensor sums,
compute their rational intertwiner spaces using the even and odd
Lie generators and the torus weights, find an intertwiner of full
rank, and solve the linear equations for its retraction.  The
existence just proved ensures termination.  All the linear
equations have rational coefficients.

Finally $\operatorname{Hom}_{G}(P_0,-)=
\operatorname{Hom}_{G_0}(\Q,-)$ is exact on locally finite rational
modules.  More generally the same assertion holds for the induced
$P$ with $E$ in place of $\Q$.  Exactness follows by placing
finitely many desired preimages in a finite $G_0$-submodule and
using its semisimplicity.  Alternatively, a finite-dimensional
image in a locally finite $G$-module has preimages in a finite
$G$-submodule, so finite projectivity suffices.  This proves the
last assertion without treating an infinite polynomial module as
a finite representation.
\end{proof}

\begin{theorem}[A finite criterion under the sign conjecture]
\label{c:finite-super-criterion}
\status{Proved under the explicitly stated hypothesis $C^+(X)$}
Let $X/\C$ be smooth projective and connected of dimension $n>0$,
let $o\in\CH^n(X)_\Q$ have degree $1$, and put
\[
 a=\dim H^{\rm even}(X,\Q),\qquad
 b=\dim H^{\rm odd}(X,\Q),\qquad
 N=(a+1)(b+1).
\]
Assume $C^+(X)$: a rational algebraic correspondence on $X$ realizes
the even cohomological projector.  Equivalently, some
$u\in\operatorname{End}(h(X))$ realizes cohomological parity
$(-1)^i$ on $H^i(X,\Q)$.  Neither idempotence of its Chow lift nor
membership in the point-and-diagonal ring is assumed.
There is a finite list of rational marked-Frobenius tensor
expressions $P_1,\ldots,P_s$, constructed from the rational cup
table by the procedure below, such that all the cohomological
relations in the requested rings on every power of $X$ hold in
rational Chow if and only if the following finite conditions hold:
\begin{enumerate}[label=(\roman*)]
\item $c_n(T_X)=(a-b)o$ and $\Delta_*o=o\boxtimes o$;
\item $y_\lambda(h(X)^{\otimes N})=0$, where
      $\lambda=((b+1)^{a+1})$ and $y_\lambda$ is its rectangular
      Young symmetrizer;
\item $P_s(h(X),\mu,1,\operatorname{tr},o,\beta^{-1})=0$ for
      every member of the constructed finite list.
\end{enumerate}
The Young symmetrizer is the row symmetrizer followed by the column
antisymmetrizer, with the tensor symmetry of Chow motives; no
choice of its nonzero normalization affects (ii).
Subject to (i), conditions (ii)--(iii) are finite lists of cycles
in the specified point-and-diagonal rings.  No assertion that
rational equivalence itself is algorithmically decidable is made.
\end{theorem}

\begin{proof}
We spell out the superalgebra construction, because exactness of
the invariant functor is unavailable here.  The hypothesis $C^+$
has exactly the usual meaning in
\citet[section~1, p.~1389]{c_tabuada2018sign}; if $v$ realizes the
even projector, $u=2v-\id$ realizes parity.
Put $A=H^*(X,\Q)$ with its parity, $L=\Q(-1)$ as an even line,
and $G=\operatorname{GL}(a|b)\times\mathbb G_m$.
All tensor representations used in the covariant construction,
including $W$, $R$ and $I$ below, belong to the category $\mathcal R$
of Lemma~\ref{c:super-projective-tensors}.  The entire regular
representation $\mathcal O(G)$ need not belong to that category.
The same marked tensors as in
Theorem~\ref{c:finite-even-criterion} give an even rational point
$t=(\mu,1,\operatorname{tr},o,\beta^{-1})$ in
\[
 W=(A\otimes A^\vee\otimes A^\vee)\oplus A
   \oplus(A^\vee\otimes L^n)
   \oplus(A\otimes L^{-n})
   \oplus(A\otimes A\otimes L^{-n}).
\]
Write $R=\mathcal O(W)$ and define
\[
 I=\ker\bigl(R\longrightarrow\mathcal O(G)\bigr),
 \qquad f\longmapsto\bigl(g\longmapsto f(g\cdot t)\bigr).
\]
Thus $I$ is the ideal of the schematic super-orbit closure.  We
retain its odd directions and nilpotents; the ideal of the orbit
of ordinary complex points would not suffice.

The supercommutative algebra
$R=\operatorname{Sym}(W_{\bar0}^\vee)\otimes
\bigwedge(W_{\bar1}^\vee)$ is finite over its ordinary even
polynomial subalgebra and is noetherian.  We can compute generators
of $I$ by orbit-graph elimination.  Coordinates on $G$ consist of
the two even matrix blocks, the odd off-diagonal blocks, and the
inverses of the two even determinants and of the last torus
coordinate.  Introduce ordinary variables for those inverses and
their defining equations.  Since the odd variables have a finite
exterior basis, elimination is a finite-module polynomial
elimination with explicitly known signed multiplication matrices.
The usual noetherian termination of module Gr\"obner elimination
therefore applies.  Enlarge the resulting finite generating span
to a finite $G$-stable submodule $E\subset I$.  The action preserves
polynomial degree, so this is a finite linear computation; one can
close the span under the rational coaction or the even and odd Lie
generators together with the torus weights.

Use Lemma~\ref{c:super-projective-tensors} to construct a finite
projective-injective $P\twoheadrightarrow E$ and a split mixed
tensor embedding of $P^\vee$.  The composite $P\to E\hookrightarrow R$
corresponds to the polynomial covariant
\[
 p_P\in(P^\vee\otimes R)^G.
\]
Its components in the chosen mixed tensor sum are our finite
expressions $P_1,\ldots,P_s$.  We also fix the nonzero projective
$P_0=\operatorname{Ind}_{G_0}^{G}(\Q)$ and a mixed tensor projector
whose image is $P_0$.

We verify both the diagram interpretation and its exact relations.
\citet[Theorem~2.1, Corollary~2.2, Theorem~2.3 and
Corollary~2.5]{c_lehrerzhang2021super} state the first and second
fundamental theorems for the general linear supergroup in this
scope: mixed invariants are contractions, and their kernel is the
tensor ideal generated by the Young symmetrizer with $a+1$ rows
and $b+1$ columns.  Although the article's title concerns an
orthosymplectic supergroup, these numbered statements are explicitly
its preliminary general linear theorems.  To pass from the
endomorphism version to all mixed tensors we bend dual indices;
the central scalar in $\operatorname{GL}(A)$ forces the two
remaining tensor degrees to be equal.  The stated kernel in
$\Q[\mathfrak S_r]$ is the ideal generated by the rectangular
symmetrizer in the first $N$ tensor positions.  Bending back gives
exactly its tensor ideal in all mixed degrees.  The scalar loop
has value $a-b$.

For polynomial covariants, the supersymmetric power in each finite
degree is the image of the averaging idempotent for the signed
symmetric-group action on the corresponding tensor power of
$W^\vee$.  That idempotent provides a split inclusion before taking
invariants.  Polarization and the preceding tensor theorem
therefore express the covariant as a finite contraction diagram;
no exactness of supergroup invariants is needed.  This is also the
superpolarization proved in Appendix~A, Proposition~A.1 and
Corollary~A.2 of the same source.  The source states these theorems
over $\C$, but all invariant spaces, maps and symmetrizers here
are defined by rational linear equations.  The equality of their
images and kernels descends along $\Q\subset\C$.
The extra torus just enforces the Tate degree.  The fixed Tate
line already satisfies $\bigwedge^2L=0$ and
$\operatorname{tr}(\id_L)=1$.

We now replace the invalid reductive invariant argument.  Let
$F\in(T'\otimes R)^G$ be any polynomial diagram covariant that
vanishes at $t$, where $T'$ is a mixed tensor target.  Equivariance
on superalgebra-valued points implies that its coefficients lie
in $I$, so $F\in(T'\otimes I)^G$.  Multiplication gives an
equivariant surjection $P\otimes R\twoheadrightarrow I$.
Regard $F\otimes\id_{P_0}$ as a $G$-map
\[
 P_0\longrightarrow T'\otimes I\otimes P_0.
\]
Projectivity lifts it through
\[
 T'\otimes P\otimes R\otimes P_0
     \twoheadrightarrow T'\otimes I\otimes P_0.
\]
The lift has bounded polynomial degree: its domain is finite and
the target is locally finite.  The preceding invariant theorem
writes it as a finite diagram between the chosen mixed tensor
summands.  Composing with $p_P$ and multiplying polynomial factors
recovers $F\otimes\id_{P_0}$.  Consequently every relation,
after tensoring with this one fixed $P_0$, lies in the tensor ideal
generated by the finite list $P_s$ and the rectangular universal
relation.  This is a projective lifting argument, not an assertion
that $(\cdot)^G$ is exact.

Evaluate in rational Chow motives, assuming (i)--(iii).  The loop
value is already $\chi(X)=a-b$.  Condition (ii) kills all universal
mixed tensor relations.  In particular, a diagram lift $e$ of the
projector defining $P_0$ is now an actual Chow idempotent: the
universal identity $e^2-e$ is among those relations.  Thus
$P_0(h(X))=\operatorname{im}e$ is defined.  Condition (iii) gives
\[
 F(h(X))\otimes\id_{P_0(h(X))}=0.
\]
It remains to cancel the projective factor, whose categorical
dimension can itself be zero.

Let $u$ be the algebraic correspondence realizing parity supplied
by $C^+(X)$.  On the mixed tensor sum $T(h(X))$ containing $P_0(h(X))$
we tensor $u$ and its dual, and use the identity on the even Tate
line; denote the resulting endomorphism by $u_T$.
Set $u_0=e u_T e$ on $P_0(h(X))$.  Its Betti realization is
parity on $P_0(A)$, because the even projector $e$ commutes with
parity in supervector spaces.  Hence
\[
 \operatorname{tr}(u_0)
  =\operatorname{str}\bigl((-1)^{\rm parity}|P_0(A)\bigr)
  =\dim P_0(A)_{\bar0}+\dim P_0(A)_{\bar1}>0.
\]
The first equality is an equality of rational scalars: Betti
realization is injective on $\operatorname{End}(\mathbf1)=\Q$.
Compose the preceding vanishing with $\id\otimes u_0$ and take
the partial categorical trace.  We obtain
$\operatorname{tr}(u_0)F(h(X))=0$, whence $F(h(X))=0$.
No equality $u^2=1$ in Chow, and no commutation of $u_T$ with $e$
in Chow, was used.  This proves sufficiency for every diagram
relation and therefore for every requested relation.

Conversely, all requested relations imply (i) by
$\Delta\cdot o_1=o_1o_2$ and
$\Delta^2=(a-b)o_1o_2$ and their pushforwards.  The Young
symmetrizer in (ii) is a finite rational combination of permutation
graphs on $X^N$, each a product of the allowed diagonals on
$X^{2N}$, and its Betti class vanishes.  Thus (ii) follows as well.
After (i), the Frobenius graph normalization in the proof of
Theorem~\ref{c:finite-even-criterion} reduces every dualized or
contracted marked diagram to a rational combination of pointed
partitions: each handle inserts $c_n(T_X)$, each marking inserts
$o$, and the two identities in (i) remove these insertions or turn
them into external point classes.  Closed components give their
degrees.  Every $P_s$ vanishes at $t$, so its normalized cycle is
one of the requested cohomological relations.  It therefore
vanishes in Chow.  This proves necessity and the finite iff.
\end{proof}

\begin{corollary}[Finite existence over an arbitrary Weil coefficient field]
\label{c:finite-weil-criterion}
\status{Proved under $C^+_H(X)$; effectiveness requires the stated coefficient data}
Let $X/k$ be smooth projective and connected over an arbitrary
algebraically closed field and of positive dimension, let $H$ be a
Weil cohomology with coefficient field $K$ of characteristic zero,
and let $o$ be a
rational zero-cycle of degree one.  Suppose $C^+_H(X)$, meaning
that a rational algebraic correspondence realizes parity on
$H^*(X)$.  The finite criterion (i)--(iii) of
Theorem~\ref{c:finite-super-criterion} still exists, with a finite
list of \emph{rational} marked diagrams, and is equivalent to
injectivity of
\[
 T_k^{\CH}(X,o)\longrightarrow H^*(X^k)
 \quad\hbox{for every }k.
\]
Here injectivity refers to rational linear relations, as required
for cycles in $\CH_\Q$.  No rational form of the whole cohomology
algebra is assumed.  The finite list is effectively constructible
if a finite presentation of the $\Q$-algebra generated by its
chosen structure constants is supplied.  For arbitrary elements
of an unspecified coefficient field, only finite existence is
asserted.  In dimension zero, $X=\operatorname{Spec}(k)$ and
all the requested relations hold immediately.
\end{corollary}

\begin{proof}
Choose a homogeneous $K$-basis of $H^*(X)$ and let
$A_\Q=\Q^{a|b}$ be a formal supervector space of the same
dimensions.  It is not asserted to be a rational form of the cup
algebra.  Build $G_\Q$, $W_\Q$ and $R_\Q=\mathcal O(W_\Q)$
from this formal space exactly as above.  The actual structure
tensors give a $K$-point $t$ of $W_\Q$.  Replace the orbit ideal by
\[
 I_\Q=\ker\left(R_\Q\longrightarrow
        \mathcal O(G_\Q)\otimes_\Q K\right),
 \qquad f\longmapsto\bigl(g\longmapsto f(g\cdot t)\bigr).
\]
The map is a $G_\Q$-comodule map, with the regular action on the
first target factor and trivial action on the coefficient field.
Its kernel is consequently a rational $G_\Q$-stable ideal.
Equivalently, the kernel assertion follows by tensoring its
underlying exact sequence with $\mathcal O(G_\Q)$ and using the
coaction identity.  A rational covariant $F$ satisfies $F(t)=0$
if and only if its coordinates lie in $I_\Q$: in one direction
equivariance gives $F(g\cdot t)=gF(t)$ on all superalgebra-valued
points, and in the other we evaluate at the identity.

Noetherianity of $R_\Q$ gives finite ideal generators; closing
their span in a bounded polynomial degree gives a finite rational
stable $E\subset I_\Q$.  Lemma~\ref{c:super-projective-tensors}
and every subsequent projective lift take place over $\Q$.
Evaluation now takes place in $K$-cohomology.  The trace cancellation
still gives the positive integer $\dim_\Q P_0$, since
$\operatorname{End}(\mathbf1)=\Q$ embeds in $K$.
The same necessity argument uses rational permutation graphs and
the Frobenius normalization in Chow.  Thus the finite iff survives
unchanged.

For the effectiveness assertion, let $c_1,\ldots,c_m\in K$ be
all the coordinates of $t$, including the inverse pairing, and
write $B=\Q[c_1,\ldots,c_m]$.  If a finite presentation
$B=\Q[z_1,\ldots,z_m]/J$ is given, the orbit map lands already
in $\mathcal O(G_\Q)\otimes B$, which injects into
$\mathcal O(G_\Q)\otimes K$.  Orbit-graph elimination including
the equations of $J$ therefore computes $I_\Q$.  Without a
presentation of $B$, no such computational claim follows merely
from the abstract knowledge of the elements $c_i$.

If instead one asks for every $K$-linear relation to vanish in
$\CH_\Q\otimes_\Q K$, the same proof is performed entirely over
$K$, using the full $K$-orbit ideal.  That is a separate
coefficient convention; it is not deduced by tensoring an
injective rational cycle map with $K$ and identifying its target
with the original cohomology space.
\end{proof}

\begin{corollary}[The sign hypothesis over an algebraic closure of a finite field]
\label{c:finite-field-sign}
\status{Proved for the stated cohomology theories}
For $X/\overline{\F}_p$, the sign hypothesis in
Corollary~\ref{c:finite-weil-criterion} holds for $\ell$-adic
cohomology, $\ell\ne p$, and crystalline cohomology with
characteristic-zero coefficients.  More generally it holds for
every Weil theory satisfying the Poincar\'e duality, weak
Lefschetz and zeta-function hypotheses in section~II of
\citet{c_katzmessing1974}.  Thus in this setting the finite
criterion has no additional sign-conjecture hypothesis.
\end{corollary}

\begin{proof}
Descend $X$ to a smooth projective geometrically connected variety
over a finite field $\F_q$.  The finitely many points and rational
coefficients representing $o$ also descend after enlarging
$\F_q$; the parity construction itself does not depend on this
choice of $o$.  Put
$Q_i(T)=\det(T-F\mid H^i(X))$.
\citet[Theorem~1, pp.~74--75; Theorem~2(1), pp.~76--77]{c_katzmessing1974}
proves that $Q_i\in\Q[T]$ and that all its roots have absolute
value $q^{i/2}$ in every complex embedding, with the same
polynomials for the cohomology theories just specified.
Consequently the nonconstant $Q_i$ are pairwise coprime.
By the Chinese remainder theorem there are polynomials
$r_i\in\Q[T]$ with $r_i\equiv1\pmod{Q_i}$ and
$r_i\equiv0\pmod{Q_j}$ for $j\ne i$.
Cayley--Hamilton shows that $r_i(F)$ acts as the projector onto
$H^i(X)$.  Each is a rational combination of graphs of iterates
of the Frobenius morphism, and hence an algebraic correspondence.
Their sum over even $i$ gives the required sign projector after
base extension to $\overline{\F}_p$.  This is the explicit
polynomial argument in Theorem~2(1), so it does not require
Frobenius semisimplicity.
\end{proof}


\begin{remark}[What the additional finite criterion settles]
\label{c:super-criterion-scope}
\status{Proved scope clarification; unrestricted classification remains unresolved}
Theorem~\ref{c:finite-super-criterion} supplies a finite cycle
criterion for mixed cohomology whenever $C^+(X)$ is known.  The
parity correspondence may be completely external to the requested
ring, as it is for the positive pair of
Proposition~\ref{c:product-parity}.  Thus that example's obstruction
to a tautological parity projector does not obstruct this proof.
We do not infer $C^+(X)$ from the point-and-diagonal relations.
Theorem~\ref{c:finite-unconditional-criterion} below removes that
hypothesis from the finite criterion by a different scalar-form
construction, which needs no cancellation step.
Corollary~\ref{c:finite-weil-criterion} removes the rational Betti
restriction for finite existence; its effectiveness qualification
is essential.
The criterion does not prove the finite Chow vanishing conditions
for an arbitrary input variety.  The general classification in
C.1(c), and in particular the K3 vanishing below, remain open
items of this task rather than consequences of finite generation.

Condition (ii) retains a precise finite-dimensionality obstruction.
Under $C^+_H(X)$ it is equivalent to Kimura finite-dimensionality,
by \citet[Theorem~3.2, p.~286]{c_delpadronemazza2005}.
To check the input, their Proposition~3.1 applies the partial-trace
identity of Theorem~2.1 to an algebraic even projector: all its
power traces equal $a$, which lies outside the forbidden contents
for the minimal $(a+1)$-by-$(b+1)$ rectangle.  The numerical ideal
is therefore nilpotent.  The same rectangle identity propagates
to tensor powers by the super tensor relations, so their
numerical ideals are nilpotent as well.  The even and odd
cohomological projectors lift across these nilpotent ideals, and
the cohomologically zero exterior and symmetric projectors on
their tensor powers are nilpotent idempotents and hence zero.
This gives the Kimura decomposition.  Conversely its even and
odd dimensions give the rectangular Schur vanishing.  In
particular, arbitrary Schur-finiteness with an unspecified larger
partition is not being substituted for this exact condition.
\end{remark}


\begin{lemma}[Exact lifting for covariant polynomial tensors]
\label{c:covariant-polynomial-lifting}
\status{Proved}
Let $V_0=\Q^{u|v}$, let $L_0$ be an even line, and set
$G=\operatorname{GL}(V_0)\times\mathbb G_m$.  Fix $q\geq1$ and put
\[
 W=\bigoplus_{m=1}^{q}(V_0^\vee)^{\otimes m}\otimes L_0,
 \qquad R=\mathcal O(W)
 =\operatorname{Sym}_{\mathrm{super}}
       \left(\bigoplus_{m=1}^{q}V_0^{\otimes m}\otimes L_0^{-1}\right).
\]
Grade $R$ by total $V_0$-degree $r$ and number $d$ of form factors.
Every $R_{r,d}$ is finite and is a direct summand of a finite sum
of copies of $V_0^{\otimes r}\otimes L_0^{-d}$.
If $I\subset R$ is a $G$-stable ideal and $E\subset I$ is a finite
homogeneous $G$-submodule generating $I$, then
\begin{equation}
 \Hom_G\bigl(V_0^{\otimes r}L_0^{-d},(E\otimes R)_{r,d}\bigr)
 \longrightarrow
 \Hom_G\bigl(V_0^{\otimes r}L_0^{-d},I_{r,d}\bigr)
 \label{eq:c-polynomial-lift}
\end{equation}
is surjective.  The same conclusion holds after a finite direct
sum in the source.  All the splittings and lifts can be found by
rational linear algebra.
\end{lemma}

\begin{proof}
The degree-$r$ polynomial representations of
$\operatorname{GL}(u|v)$ are modules for its Schur superalgebra.
The coefficient-coalgebra description and the explicit
coefficient-basis proof in
\citet[section~5, (5.2), Lemma~5.1 and Theorem~5.2,
author pp.~13--15]{c_brundankujawa2003}
identify that algebra, over $\Q$, with
\[
 S(u|v,r)=\End_{\Q[\mathfrak S_r]}(V_0^{\otimes r}),
\]
where the symmetric group acts by superpermutations.
Maschke's theorem makes this a semisimple algebra: decompose the
finite $\Q[\mathfrak S_r]$-module into its simple isotypical
parts and take the product of their matrix endomorphism algebras.
This is also the characteristic-zero argument in
\citet[Theorem~1.3.1(a)(i), proof of Theorem~4.2.6(a),
author pp.~3 and~12]{c_hemmerkujawanakano2006}.
An ungraded splitting can be made even by replacing a map $f$
by $(f+\Pi f\Pi)/2$, where $\Pi$ denotes intrinsic parity on the
source and target.  Conjugation by $\Pi$ preserves the module
intertwiner equations, so this averaging still gives a splitting.
Thus the category with even morphisms is semisimple as well.
Tensoring by a fixed character of the independent last torus does
not change this conclusion.

The degree-$d$ supersymmetric power is the image of the averaging
idempotent on the $d$-fold tensor power of its generators.  After
fixing total degree $r$, the latter is a finite sum of copies of
$V_0^{\otimes r}L_0^{-d}$.  Finiteness follows also directly from
$d\leq r\leq qd$, except for $R_{0,0}=\Q$.
The two central scalar actions split every stable ideal into
these homogeneous components.  The multiplication surjection
$(E\otimes R)_{r,d}\twoheadrightarrow I_{r,d}$ is therefore a
surjection of finite polynomial representations of the same
degree and Tate character.  Its source, target and the source of
the displayed Hom functor belong to the semisimple category just
proved.  It splits, proving the assertion.  Intertwining and
splitting impose finite rational linear equations.  We have used
exactness in the covariant polynomial category, not exactness of
invariants on arbitrary mixed rational superrepresentations.
\end{proof}

\begin{theorem}[A finite criterion without a sign hypothesis]
\label{c:finite-unconditional-criterion}
\status{Proved; finite existence is unconditional, and effectiveness has the stated input condition}
Let $X/k$ be smooth projective and connected, of dimension $n>0$
over an algebraically closed field, let $H$ be a Weil cohomology
with coefficient field $K$, and let
$o\in\CH^n(X)_\Q$ have degree one.  Set
\[
 a=\dim_K H^{\mathrm{even}}(X),\qquad
 b=\dim_K H^{\mathrm{odd}}(X),\qquad
 N=a(b+1),\qquad \lambda=((b+1)^a).
\]
In Chow motives put $A=h(X)$, $L=\Q(-n)$,
\[
 \epsilon(x)=\operatorname{tr}(o\cdot x),\qquad
 \pi=\id_A-\eta\epsilon,\qquad V=(A,\pi),
 \qquad
 \kappa_m=\operatorname{tr}\circ\mu^{(m)}\circ\pi^{\otimes m}
       :V^{\otimes m}\longrightarrow L.
\]
Here $\eta$ is the unit and $\mu^{(m)}$ is iterated cup product;
the displayed notation for $\epsilon$ includes its usual Tate
cancellation.  Since $\epsilon\eta=1$, $\pi$ is an actual Chow
projector, independently of any standard conjecture.  There is
a finite list of rational expressions $P_1,\ldots,P_s$ in the
forms $\kappa_1,\ldots,\kappa_{2n}$, using tensor products,
permutations and rational linear combinations, such that the
following are equivalent:
\begin{enumerate}[label=(\Alph*)]
\item All rational cohomological relations among the specified
      point and partial-diagonal generators on every power of $X$
      hold in rational Chow.
\item The following finite list of Chow identities holds:
      \begin{enumerate}[label=(\roman*)]
      \item $c_n(T_X)=(a-b)o$ and $\Delta_*o=o\boxtimes o$;
      \item $\kappa_{2n+1}=0$;
      \item the rectangular Young symmetrizer
            $y_\lambda:V^{\otimes N}\to V^{\otimes N}$ is zero;
      \item $P_\nu=0$ for $1\leq\nu\leq s$.
      \end{enumerate}
\end{enumerate}
Each identity in (ii)--(iv), after Poincar\'e-duality bending,
is a finite rational combination of cycles in the requested
point-and-diagonal rings.  The list can be constructed effectively
from a rational cohomology table, and more generally from a finite
presentation of the $\Q$-algebra generated by its structure
constants.  With an arbitrary unspecified Weil coefficient
field, finite existence is asserted, without an algorithm for
testing its constants.  No decidability assertion for rational
equivalence is made.  No algebraic parity projector is assumed.
More explicitly, the bent form $\kappa_m$ is the modified diagonal
\begin{equation}
 \Gamma_m(X,o)=
 \sum_{\varnothing\ne J\subseteq\{1,\ldots,m\}}
       (-1)^{m-|J|}\Delta_J\prod_{i\notin J}p_i^*o
       \ \in\CH^{n(m-1)}(X^m)_\Q,
 \label{eq:c-explicit-higher-modified-diagonal}
\end{equation}
where a singleton diagonal means the unit on that coordinate.
Thus condition (ii) is the explicit identity
$\Gamma_{2n+1}(X,o)=0$.
\end{theorem}

\begin{proof}
We first normalize the Frobenius algebra.  Poincar\'e duality
identifies the equality $\Delta_*o=o\boxtimes o$ with
\[
 \epsilon\mu=\epsilon\otimes\epsilon.
\]
Indeed the two sides pair with $x\otimes y$ as, respectively,
$\operatorname{tr}(oxy)$ and
$\operatorname{tr}(ox)\operatorname{tr}(oy)$.  This calculation
is an identity of Chow correspondences, since the pairing and
its inverse are the diagonal evaluation and coevaluation.
Consequently $V=\ker\epsilon$ is an algebra ideal, with a
restricted multiplication $\mu_V$.  For $m\geq1$ we have
\begin{equation}
 \kappa_{m+1}
   =\kappa_m\circ(\mu_V\otimes\id_V^{\otimes(m-1)}).
 \label{eq:c-high-form-recursion}
\end{equation}
Thus (ii) kills all $\kappa_m$ with $m>2n$ in Chow.  In
cohomology $H(V)=H^{>0}(X)$, so these same forms vanish by
degree for every $m>2n$.  This explains why a single additional
modified-diagonal identity suffices to truncate the forms.

Here is the finite construction.  Let $V_0=\Q^{a-1|b}$ and let
$L_0$ be a formal even line.  We do not assert that $V_0$ is a
rational form of the given cohomology algebra.  Apply
Lemma~\ref{c:covariant-polynomial-lifting} with $q=2n$ and write
$G=\operatorname{GL}(V_0)\times\mathbb G_m$.
The cohomological forms give an even $K$-point
$t=(\kappa_1,\ldots,\kappa_{2n})$ of $W$.  Define
\begin{equation}
 I_\Q=\ker\left(R\longrightarrow\mathcal O(G)\otimes_\Q K\right),
 \qquad f\longmapsto\bigl(g\longmapsto f(g\cdot t)\bigr).
 \label{eq:c-pure-form-orbit}
\end{equation}
This is the kernel of a rational comodule map, with the regular
action on $\mathcal O(G)$ and trivial action on $K$.  It is
therefore a $G$-stable ideal.  We use the schematic super-orbit:
odd directions and nilpotents are retained.  Taking only the
ordinary-point orbit would be a different and insufficient
construction.  The two central scalar actions make $I_\Q$
homogeneous in $(r,d)$.

The supercommutative algebra $R$ is a finite module over the
polynomial algebra on its even generators, and hence is
noetherian.  Choose finite ideal generators of $I_\Q$ and close
their homogeneous span under $G$.  This yields a finite
homogeneous stable submodule $E\subset I_\Q$ that still generates
the ideal.  Each $E_{r,d}$ lies in the finite semisimple module
$R_{r,d}$ and hence is a split summand of a finite sum of copies
of $V_0^{\otimes r}L_0^{-d}$.  Choose a finite such sum $T$ and
maps
\[
 s:E\longrightarrow T,\qquad p:T\longrightarrow E,
 \qquad ps=\id_E.
\]
For each tensor summand $T_\nu=V_0^{\otimes r_\nu}L_0^{-d_\nu}$,
the composite
\begin{equation}
 T_\nu\longrightarrow T\xrightarrow{p}E\hookrightarrow R
 \label{eq:c-finite-form-relations}
\end{equation}
is a polynomial multilinear form expression.  After evaluation
at $t$ and cancellation of the indicated line twist it is a map
$V_0^{\otimes r_\nu}\otimes K\to L_0^{d_\nu}\otimes K$.
Its Chow expression is $P_\nu$.

We justify both the expression and every equality among such
expressions.  In each bidegree the symmetric-power averaging
idempotents split $R_{r,d}$ into a finite sum of pure tensor
powers with its fixed line character.  The first fundamental
theorem in \citet[Theorem~2.1]{c_lehrerzhang2021super} says that
the intertwiners between equal pure powers are generated by
permutations.  Hence \eqref{eq:c-finite-form-relations} is a rational combination of products
of the universal forms, evaluated on disjoint groups of inputs.
The second fundamental theorem, their Theorem~2.3, says that all
relations between those permutation maps are generated by the
Young symmetrizer with $(a-1)+1=a$ rows and $b+1$ columns.
The central scalar excludes unequal pure degrees.  Inserting
the averaging idempotents merely gives direct summands of this
same description.  These statements descend from $\C$ to $\Q$
because the maps and their kernel equations are rational.
The fixed Tate line is even and invertible, and its usual
one-dimensional identities already hold in Chow motives.
Thus (iii) makes every representation-theoretic equality used
in this construction a valid Chow tensor identity.

Now let a rational universal form
\[
 F:V_0^{\otimes r}L_0^{-d}\longrightarrow R
\]
be equivariant and satisfy $F(t)=0$.  Equivariance on all
superalgebra-valued points gives $F(g\cdot t)=gF(t)=0$;
equivalently its image is contained in $I_\Q$.  Conversely,
membership in $I_\Q$ implies zero evaluation at the identity.
The multiplication map $E\otimes R\twoheadrightarrow I_\Q$
and Lemma~\ref{c:covariant-polynomial-lifting} give a lift
\[
 \widetilde F:V_0^{\otimes r}L_0^{-d}
                 \longrightarrow(E\otimes R)_{r,d}.
\]
Compose with $s\otimes\id$ to obtain a lift with target
$(T\otimes R)_{r,d}$.  Expanding this lift with the pure-tensor
first fundamental theorem expresses $F$ as a sum of expressions
that apply one of the maps \eqref{eq:c-finite-form-relations} to some inputs and other
universal forms to the remaining inputs.  All intervening
objects are covariant polynomial representations of total
$V_0$-degree $r$ and line character $-d$; no exactness assertion
for an arbitrary mixed invariant functor has entered.
The equality is valid as a formal tensor identity modulo the
rectangular relation by the preceding second fundamental
theorem.  Conditions (iii)--(iv) therefore imply
$F(\kappa_1,\ldots,\kappa_{2n})=0$ in Chow.  There is no
projective auxiliary tensor to cancel and no use of $C^+(X)$.

We next return from forms to the requested cycles.  By (i) and
the graph product rule, every monomial in the generators reduces
to a scalar times a pointed-partition cycle.  A homogeneous
relation between these cycles, with $k$ coordinates and $d$
unpointed blocks, bends to a map
$A^{\otimes k}\to L^d$.  For an unpointed block $B$ its factor
is $\operatorname{tr}(\prod_{i\in B}x_i)$; for a pinned
coordinate it is $\epsilon(x_i)$.  Split every input using the
actual decomposition $A=\mathbf1\oplus V$.  A pinned $V$ input
annihilates its term.  Unit inputs in an unpointed block are
deleted; if no input remains, its value is
$\operatorname{tr}(1)=0$, since $n>0$.  Every other block is one
of the forms $\kappa_m$.  Condition (ii) removes blocks with
$m>2n$.  Each remaining restriction is exactly a rational form
$F$ of the preceding paragraph.  If the original relation is
cohomologically zero, each such restriction vanishes at $t$,
and hence vanishes in Chow.  The direct-sum decomposition of
$A^{\otimes k}$ and Poincar\'e duality prove the original cycle
relation.  This proves (B)$\Rightarrow$(A).

For necessity, the cohomological relations
$\Delta_*o-o\boxtimes o=0$ and
$\Delta_X^2-(a-b)o\boxtimes o=0$ belong to the requested ring
on $X^2$.  Pushing the second along a projection and using
self-intersection gives $c_n(T_X)=(a-b)o$.  The form
$\kappa_{2n+1}$ is cohomologically zero by positive degrees.
Its bent Chow cycle is obtained from the small diagonal by
inserting $\pi$ in each input.  Expanding
$\pi=\id-\eta\epsilon$ gives pointed partial-diagonal cycles.
The same expansion shows that $y_\lambda(V)$, viewed as a
correspondence on $X^N$ and hence a cycle on $X^{2N}$, is a
rational combination of products of pair diagonals and point
classes.  It is cohomologically zero by the rectangular theorem.
Finally every $P_\nu$ is a product of the forms just described,
with permuted inputs, and is cohomologically zero because
$E\subset I_\Q$.  Bending gives the claimed cycles in the
requested rings.  Thus (A) implies (ii)--(iv) as well as (i).
For the displayed formula for $\Gamma_m$, the adjoint of
$\pi=1-\eta\epsilon$ under Poincar\'e duality is
$1-o\operatorname{tr}$.  Expanding its action on the small
diagonal gives the terms indexed by nonempty $J$ in
\eqref{eq:c-explicit-higher-modified-diagonal}.  The omitted empty
term integrates every coordinate and is
$\operatorname{tr}(1)=0$.

For effectiveness, the orbit map is computed using two even
matrix blocks, odd off-diagonal entries, inverses of the two
even determinants, and the last torus coordinate and its inverse.
If the coordinates of $t$ generate a presented algebra
$B=\Q[c_1,\ldots,c_h]$, compute the kernel into
$\mathcal O(G)\otimes B$ by orbit-graph elimination, including
the defining equations of $B$.  This is finite-module
polynomial elimination over the even variables: the odd
variables have a finite exterior basis, with known signed
multiplication.  Noetherian module Gr\"obner elimination
terminates.  The kernel agrees with \eqref{eq:c-pure-form-orbit}
because $B\hookrightarrow K$.  Stable spans, the splittings,
permutation expressions and all lifts then use finite rational
linear algebra.  Without such a presentation only the
noetherian existence argument is used.  This completes the proof.
\end{proof}

\begin{corollary}[A finite power detects all the relations]
\label{c:finite-power-cutoff}
\status{Proved; the cut-off is effective with the same coefficient input}
In Theorem~\ref{c:finite-unconditional-criterion}, let $r_\nu$
be the input arity of $P_\nu$ and put
\[
 B=\max\{2,\ 2n+1,\ 2a(b+1),\ r_1,\ldots,r_s\}.
\]
Then all rational cohomological relations hold in Chow on every
power of $X$ if and only if the cycle map is injective on
$T_B^{\CH}(X,o)$.  In particular such a finite cut-off exists
without $C^+(X)$ or a rational form of the Weil cohomology.
Its effective construction has exactly the coefficient-data
qualification in the theorem.
\end{corollary}

\begin{proof}
All finite tests in the theorem are bent cycles on powers of
arity at most $B$.  The two normalization tests are detected by
the indicated relations on $X^2$, the high form by $X^{2n+1}$,
the rectangle by $X^{2N}$, and $P_\nu$ by $X^{r_\nu}$.
Injectivity on $T_B^{\CH}$ implies injectivity on every
$T_k^{\CH}$ with $k\leq B$: pull back along $p:X^B\to X^k$,
and retract that pullback by
\[
 z\longmapsto p_*\left(z\prod_{i=k+1}^{B}p_i^*o\right).
\]
The composite on $\CH^*(X^k)_\Q$ is the identity because
$\deg o=1$.  Thus injectivity on the single specified power
implies all the finite tests, and the theorem applies.  The
converse is immediate.  The cut-off depends only on the
cohomological form data and the choices in its finite
construction, not on an unproved Chow vanishing.
\end{proof}

\begin{lemma}[A coefficient-independent recursive elimination bound]
\label{c:recursive-elimination-bound}
\status{Proved}
For integers $v,d\geq1$, let $\mathcal B(v,d)$ be the tree of
finite sequences $(\alpha_0,\ldots,\alpha_{r-1})$ in $\Z_{\geq0}^v$
satisfying
\[
 |\alpha_i|\leq d2^i,
 \qquad \alpha_j\not\leq\alpha_i\quad(0\leq j<i<r),
\]
where $\leq$ is coordinatewise.  This is a finite tree.  Its
maximum sequence length $\mathcal L(v,d)$ is a total computable
integer-valued function.  Put
\[
 \mathcal D(v,d)=d\,2^{\mathcal L(v+1,d)}.
\]
Over any field, an ideal in $v$ ordinary polynomial variables
generated in degrees at most $d$ has, for every monomial order,
a Gr\"obner basis in degrees at most $\mathcal D(v,d)$.
In particular this bound applies to an elimination order.
\end{lemma}

\begin{proof}
We give the elementary argument, so no closed Gr\"obner-degree
estimate is needed.  Dickson's assertion that an infinite
sequence in $\Z_{\geq0}^v$ contains a comparable pair follows
by induction on $v$.  An infinite sequence of nonnegative
integers has an infinite nondecreasing subsequence: either some
value repeats infinitely or one successively selects larger
values.  Apply this observation to the first coordinate and
then apply induction to the remaining coordinates.  Thus
$\mathcal B(v,d)$ has no infinite branch.  Each node has only
finitely many children, by its displayed degree bound.
If the tree were infinite, repeatedly choosing a child with
infinitely many descendants would construct an infinite branch.
It is consequently finite.  Enumerating its successive finite
levels until the next level is empty computes $\mathcal L(v,d)$.

Homogenize the given generators with one new variable $h$.
Extend the desired monomial order by comparing total degree
first and, at equal degree, comparing the monomials after
$h=1$ in the desired order.  Process the homogeneous generators
one at a time, reducing each by the previous nonzero remainders,
and then run the usual pair-reduction algorithm.  Every accepted
nonzero remainder has a leading monomial not divisible by any
previously accepted leading monomial.  Homogeneous division
does not increase degree, and the least common multiple of
two leading monomials has degree at most the sum of their
degrees.  If the accepted polynomials are numbered from zero,
their degrees are therefore at most $d2^i$.  Their leading
exponent vectors form a sequence in $\mathcal B(v+1,d)$.
At most $\mathcal L(v+1,d)$ polynomials can be accepted, and
all their degrees are at most $\mathcal D(v,d)$.
There are then only finitely many pairs to process.

For completeness, the pair criterion here says that if every
pair polynomial reduces to zero, every element of the generated
ideal has leading monomial divisible by one of the recorded
leading monomials.  Indeed write such an element as a sum of
monomial multiples of the recorded polynomials.  If the greatest
displayed leading monomial cancels, subtract the corresponding
pair relations to lower that greatest monomial.  A representation
minimizing it, and then the number of terms attaining it, cannot
have such a cancellation.  Its greatest monomial is therefore
the leading monomial of the element.  This proves the criterion
and that the resulting list is a Gr\"obner basis.

Dehomogenization gives a Gr\"obner basis for the original ideal.
To check this even though the homogenized generator ideal need
not be saturated, if $f$ belongs to the original ideal, then
$h^e f^{\mathrm h}$ belongs to the homogenized generator ideal
for some $e\geq0$, by homogenizing an expression for $f$ in
the generators to a common degree.  A homogeneous leading
monomial of our basis divides that of $h^e f^{\mathrm h}$.
After setting $h=1$, it divides the leading monomial of $f$
in the desired order.  Degrees do not increase.  This proves
the bound for arbitrary, including elimination, orders.
\end{proof}

\begin{corollary}[A uniform arity bound for rational cohomology tables]
\label{c:uniform-arity-bound}
\status{Proved in the stated rational-table scope}
Suppose the graded Weil cohomology algebra in
Theorem~\ref{c:finite-unconditional-criterion} has a rational
table.  This includes rational Betti cohomology over $\C$.
Put $q=2n$, $u=a-1$, $v=b$, $s=u+v$, and
\[
 M_0=\frac12\sum_{m=1}^{q}\bigl(s^m+(u-v)^m\bigr),\qquad
 M_1=\frac12\sum_{m=1}^{q}\bigl(s^m-(u-v)^m\bigr).
\]
Define the following integers using the function of the preceding
lemma:
\[
 O=M_1+4uv,\qquad
 E=M_0+2(u^2+v^2)+2,\qquad
 J=2^O,\qquad P=E+J,
\]
\[
 D=\mathcal D(P,q+2),\qquad
 B_{n,a,b}=\max\{2,\ q+1,\ 2a(b+1),\ q(D-1+M_1)\}.
\]
Then all the point-and-diagonal cohomological relations hold in
rational Chow on all powers if and only if the cycle map is
injective on $T_{B_{n,a,b}}^{\CH}(X,o)$.
The bound depends only on $(n,a,b)$, and is a total computable
integer.  It is not a practical or minimal bound, and does not
decide the required Chow injectivity.
\end{corollary}

\begin{proof}
We bound the input arities of the finite form relations in
Theorem~\ref{c:finite-unconditional-criterion}.  Its parameter
superspace has $M_0$ even and $M_1$ odd coordinates.
Present $\operatorname{GL}(u|v)$ by two independent even
supermatrices $G,H$ with relations $GH=HG=1$; $H$ is the inverse
of $G$.  This presentation has $2(u^2+v^2)$ even and $4uv$ odd
matrix entries.  Present the independent line torus by $z,w$
with $zw=1$.  The orbit-graph equation for a form of arity $m$
is a coordinate of the parameter tensor minus a coordinate of
$z(H^\vee)^{\otimes m}t_m$, with the fixed dual-action signs.
It has total degree at most $m+1\leq q+1$; all group equations
have degree two.  Because the table is rational, its structure
constants add no variables or degrees.

There are $E$ even and $O$ odd variables in this graph
presentation.  The free supercommutative algebra is a free module
of rank $J=2^O$ over the polynomial ring $A_0$ on its even
variables, with exterior-monomial basis $e_1,\ldots,e_J$.
If $f_1,\ldots,f_t$ are the graph equations, their ideal is the
$A_0$-submodule generated by all $e_j f_i$.  The polynomial
coordinates of these module generators have even-variable
degree at most $q+1$.  Encode a vector $(p_1,\ldots,p_J)$ by
the ordinary polynomial $\sum p_jY_j$, and form the ideal in
$A_0[Y_1,\ldots,Y_J]$ generated by the encoded module generators
and all $Y_iY_j$.  Its part of total $Y$-degree one is exactly
the original module.  This ordinary ideal has $P=E+J$ variables
and generators of degree at most $q+2$.

Choose an elimination order that eliminates the even group
variables and every $Y_j$ whose exterior monomial involves an
odd group variable.  Retain the even parameter variables and
the remaining $Y_j$.  The ordinary ideal is homogeneous in
total $Y$-degree.  The homogeneous pair-reduction construction in
Lemma~\ref{c:recursive-elimination-bound} preserves this additional
grading, both before and after dehomogenization, and therefore
gives a bounded Gr\"obner basis homogeneous in $Y$-degree.  The elimination basis elements of
$Y$-degree one give module generators for the intersection with
the original parameter superalgebra, namely the schematic
super-orbit ideal.  Quadratic $Y$ terms cannot contribute to
this degree-one part.

Lemma~\ref{c:recursive-elimination-bound} bounds these ordinary
basis elements by $D$.  Their coefficients in the retained even
variables have degree at most $D-1$.  Returning to exterior
monomials adds at most $M_1$ odd parameter factors.  Every
parameter coordinate has input arity at most $q$.  The orbit
ideal therefore has generators of input arity at most
$q(D-1+M_1)$.  Taking homogeneous components and closing their
span under the supergroup preserves input arity.  The split
pure tensor modules used to define $P_\nu$ can consequently
all be taken within this same bound.  The assertion now follows
from Corollary~\ref{c:finite-power-cutoff}.
\end{proof}

\begin{remark}[The coefficient qualification in the uniform bound]
\label{c:uniform-arity-scope}
\status{Proved scope clarification}
The finite criterion and finite cut-off of
Theorem~\ref{c:finite-unconditional-criterion} and
Corollary~\ref{c:finite-power-cutoff} apply to an arbitrary
Weil field.  The displayed bound depending only on $(n,a,b)$
has the additional rational-table hypothesis.  For a presented
coefficient algebra whose $h$ named generators include every
coordinate of $t=(\kappa_1,\ldots,\kappa_q)$, with its exact
defining ideal generated in degrees at most $e$, the same proof gives a bound by adding
$h$ to $E$ and replacing the generator-degree input $q+2$
by $\max(q+3,e+1)$.  The extra one accounts for a structure
constant occurring as a variable in an orbit equation.  These
$h$ coefficient variables are eliminated together with the group
variables, so the result is the rational contraction $I_\Q$.
We make no uniform rational-contraction claim for unspecified
elements of an arbitrary coefficient field.

The accompanying exact certificate exhausts only
$\mathcal L(1,d)=d+1$ for $1\leq d\leq4$ and
$\mathcal L(2,1)=4$, and checks the square-zero elimination
mechanism for the explicit super map
$z=x$, $\theta_1=\eta$, $\theta_2=x\eta$.
It does not evaluate the enormous geometric cut-off above.
No closed Dub\'e degree formula is used in its proof.
\end{remark}


\begin{remark}[What the finite criterion settles]
\label{c:finite-unconditional-scope}
\status{Proved scope statement; the actual general Chow vanishing remains unresolved}
Theorem~\ref{c:finite-unconditional-criterion} removes the sign
hypothesis from the finite criterion, including for an arbitrary
Weil coefficient field.  It supplies finite generators for the
all-power obstruction and does not assert that these cycles
vanish for every variety.  It also does not make rational
equivalence decidable.  Its list is not claimed minimal; in
particular the more precise K3 obstruction identified above is
unchanged.  The explicitly requested all-power K3 assertion and
the unrestricted evaluation of the criterion for every $(X,o)$
therefore remain separate unresolved parts of C.1(c).
If $n=0$, connectedness over the algebraically closed field gives
$X=\operatorname{Spec}(k)$, and every relation holds immediately.
\end{remark}


\begin{theorem}[The narrow K3 ring already detects finite-dimensionality]
\label{c:k3-gate}
\status{Proved equivalence; the general vanishing is unresolved}
Let $S/\C$ be a projective K3 surface and let $o$ be its Beauville--Voisin class. The following assertions are equivalent:
\begin{enumerate}[label=(\roman*)]
\item for every $k$, all the cohomological relations among the point and diagonal classes hold in $\CH^*(S^k)_\Q$;
\item with $\pi_2=\Delta_S-o\times S-S\times o$, the cycle
\begin{equation}
\label{c:k3-determinant}
 D_{23}=\sum_{\sigma\in S_{23}}\sgn(\sigma)
          \prod_{a=1}^{23}(\pi_2)_{a,23+\sigma(a)}
 \quad\text{vanishes in }\CH^{46}(S^{46})_\Q;
\end{equation}
\item $\bigwedge^{23}h^2(S)=0$ in rational Chow motives;
\item $h(S)$ is finite-dimensional in the sense of Kimura--O'Sullivan.
\end{enumerate}
Thus the request for every K3 surface is not supplied by the Beauville--Voisin small-diagonal theorem. The first missing universal determinant for this narrower ring occurs on $46$ factors.
\end{theorem}

\begin{proof}
The projectors $o\times S$, $\pi_2$, and $S\times o$ are mutually orthogonal Chow projectors, giving $h(S)=\mathbf1\oplus h^2(S)\oplus\Q(-2)$ (with the contravariant/cohomological convention). The cycle $\pi_2$ belongs to the requested ring. As a correspondence on $S^{23}$, $D_{23}$ is $23!$ times the alternating projector on $h^2(S)^{\otimes23}$. Consequently its vanishing is exactly (iii), because an object in a pseudo-abelian category is zero precisely when its defining idempotent is zero. Its cohomology class is zero since $\dim H^2(S)=22$, so (i) implies (ii). Assertion (iii) makes $h^2(S)$ evenly finite-dimensional, and the other two summands are Tate motives; hence it implies (iv). Conversely, finite-dimensionality makes the even summand have exterior dimension $22$, giving (iii); alternatively the dimension assertion is the finite-dimensionality statement recalled in \citet[\S2.4, p.~506]{c_yin2015}. Finally, \citet[unnumbered main Theorem, p.~503, and \S\S2.5, 3.1--3.6]{c_yin2015} proves that finite-dimensionality is equivalent to injectivity for the larger ring generated by all divisor classes and diagonals. Its restriction gives (i).

We can also see sufficiency directly in the smaller ring. \citet[Theorem 1 and Proposition 3.2]{c_beauvillevoisin2004} gives $c_2(T_S)=24o$ and the small-diagonal relation. These imply every local relation in \eqref{c:pair-local} with $b=22$. The only additional generators supplied by Theorem~\ref{c:pair-ranks} are the permutations and products of \eqref{c:k3-determinant}. This proves that no larger unlisted family of relations is hidden in the equivalence.
\end{proof}

The statement just proved is a reduction, not a proof of (i) for general K3 surfaces. The current-source check includes \citet[Theorem L, explicitly Theorem 9.7 and Corollary 9.8]{c_floccarifu2026}. That result proves abelian Chow motives for K3 surfaces whose transcendental quadratic space embeds in $U_\Q^{\oplus2}\oplus\langle-1\rangle\oplus\langle-b\rangle$, with $b>0$ and $b\equiv3\pmod4$; it supplies new four-dimensional families of generic Picard rank $16$. It does not assert the vanishing of \eqref{c:k3-determinant} for every projective K3 surface. The general classification of all pairs $(X,o)$ requested in C.1(c) also remains a separate open row; the finite presentation above gives its exact cycle conditions but does not turn them into a geometric classification.

\begin{corollary}[Removing the algebraic part of the K3 gate]
\label{c:k3-transcendental-gate}
\status{Proved smaller equivalent vanishing test; general vanishing unresolved}
Let $S/\C$ be a projective K3 surface of Picard rank $\rho$, let
$D_1,\ldots,D_\rho$ be a rational basis of $\operatorname{NS}(S)_\Q$,
and let $B=(\deg D_aD_b)_{a,b}$.  Define
\[
 \pi_{\rm alg}=\sum_{a,b}(B^{-1})_{ab}D_a\times D_b,
 \qquad \pi_{\rm tr}=\pi_2-\pi_{\rm alg},\qquad
 t_2(S)=(S,\pi_{\rm tr},0).
\]
Every assertion of Theorem~\ref{c:k3-gate} is also equivalent to
\[
 \bigwedge^{23-\rho}t_2(S)=0.
\]
Equivalently, the alternating sum of the products
$(\pi_{\rm tr})_{a,\,23-\rho+\sigma(a)}$, over
$\sigma\in S_{23-\rho}$, vanishes on $S^{2(23-\rho)}$.
For Picard rank $1$ this is a test on $44$ factors.  This smaller
correspondence uses divisor classes outside the narrow requested
ring; its factor count therefore does not contradict the unconditional
narrow-ring result through $45$ factors.
\end{corollary}

\begin{proof}
The intersection matrix $B$ is nonsingular by the Hodge index
theorem.  Direct composition of the product correspondences shows
that $\pi_{\rm alg}$ is an idempotent with image $\Q(-1)^{\oplus\rho}$,
and that it is orthogonal to the degree-zero and degree-four
projectors.  Thus
\[
 h^2(S)=\Q(-1)^{\oplus\rho}\oplus t_2(S).
\]
For a direct sum in a rational pseudo-abelian symmetric category,
the exterior-power decomposition gives
\[
 \bigwedge^{23}h^2(S)
 =\bigoplus_{a=0}^{\rho}
   \bigwedge^a\bigl(\Q(-1)^{\oplus\rho}\bigr)
       \otimes\bigwedge^{23-a}t_2(S).
\]
The $a=\rho$ summand is
$\Q(-\rho)\otimes\bigwedge^{23-\rho}t_2(S)$.
Its vanishing follows from that of the whole direct sum; tensoring
with the inverse Tate object proves one implication.
Conversely, every higher exterior power is a direct summand of
the tensor product of the $(23-\rho)$th exterior power with a
further tensor power.  Hence the stated vanishing kills every
summand in the display.  The correspondence form is exactly the
alternating-projector description.  This also checks directly the
transcendental reduction in \citet[\S2.5, pp.~506--507, equation~(7)]{c_yin2015}
while explaining its relation to the smaller ring in the question.
\end{proof}


\begin{corollary}[The unconditional K3 range]
\label{c:k3-stable-range}
\status{Proved}
For every projective complex K3 surface with its Beauville--Voisin class, the requested relations hold in every power $S^k$ with $k\leq45$. More generally they hold in normalized degree $d$ of every $S^k$ whenever
\[
 \min(d,k-d)\leq22,
\]
where normalized degree $d$ means Chow codimension $2d$. In this range no finite-dimensionality assumption is required.
\end{corollary}

\begin{proof}
The Beauville--Voisin identities give the local relations, so every Chow monomial reduces to a linear combination of disjoint matching monomials
$\prod_{i\in I}o_i\prod_{\{a,b\}\in\mathcal M}q_{ab}$.
If the matching has $m$ edges, degree $d$ means $|I|+m=d$, and the number of used factors is $|I|+2m$. Consequently $m\leq\min(d,k-d)$. For $m\leq22$, the orthogonal matching tensors in $(H^2(S))^{\otimes2m}$ are linearly independent: the first determinantal relation requires $23$ edges, as proved in Theorem~\ref{c:pair-ranks}. Monomials with different point and matching supports lie in different K\"unneth multidegrees. Thus their cohomological classes are independent, and no homologically zero Chow combination remains. For $k\leq45$, every matching has at most $22$ edges. This proves both assertions.
\end{proof}

The remaining relation cannot be obtained merely by asserting that all lower-power relations hold. It first appears in the middle normalized degree of $S^{46}$ and is exactly \eqref{c:k3-determinant}. We also tested three possible ways to remove this obstruction. A Kuga--Satake correspondence giving a splitting in homological motives would produce a composite $1+\eta$ in Chow, with $\eta$ homologically trivial; promoting it to a Chow splitting still requires nilpotence or another argument making $1+\eta$ invertible. Specializing to dense loci with finite-dimensional motives does not prove the generic Chow vanishing, because the rational-equivalence locus is only controlled as a countable union of closed loci without a uniform family of equivalences. Finally, the known small-diagonal identities are precisely the local matching reductions and do not supply the first orthogonal determinant. These are failed proof steps, not counterexamples to finite-dimensionality.

The further nilpotence test has a precise limitation.  Put
$d=22-\rho(S)$ and let $a_{d+1}$ be the normalized alternating
idempotent on $t_2(S)^{\otimes(d+1)}$.  Since $a_{d+1}^2=a_{d+1}$,
composition-nilpotence of this particular correspondence is equivalent
to its vanishing.  Also, the symmetric-category identity
\[
 \operatorname{ptr}_{m}(a_m)=\frac{d-m+1}{m}a_{m-1}
\]
shows that its partial trace is zero at $m=d+1$; this operation does
not produce a smaller nonzero obstruction.  In the formula, permutations
fixing the last factor contribute $d/m$, and deleting that factor from
a nontrivial cycle contributes $-(m-1)/m$.  Nilpotence only in
$\End(h(S))$ has not been shown to propagate to this higher tensor
power.  The all-powers quantifier in
\citet[Theorem~0.2, author p.~2]{c_jannsen2006finite} cannot be dropped.

There is a valid conditional route in the stated K3 families of
\citet[Theorems~3.5--3.6, with Proposition~1.3]{c_voisin2014bloch}:
her relative argument supplies the needed Chow lifting once the general
Kuga--Satake correspondence is algebraic.  This does not prove that
remaining algebraicity premise.  Nor does the codimension-two argument
apply directly to the higher-codimension alternating projector.
The detailed method audit records these distinctions; none changes the
unresolved vanishing status.

The current-source search also checked the official withdrawal notice of the October 2026 paper claiming the Hodge conjecture for all K3 products \citep{c_k3withdrawal2026}. We exclude that withdrawn proof. Even a valid homological Kuga--Satake splitting alone would leave the Chow nilpotence step just identified.

\begin{example}[An explicit genus-three failure]
\label{c:fermat-counterexample}
\status{Proved counterexample}
Let $C/\C$ be the plane quartic $X^4+Y^4=Z^4$, and take $o=(0:1:1)$. Its modified diagonal $M_o$ vanishes in every Weil cohomology, but
\[
 M_o\ne0\quad\hbox{in }\CH^2(C^3)_\Q.
\]
Thus the cohomological relation in Proposition~\ref{c:modified-diagonal} fails in rational Chow groups for this completely specified genus-three pointed curve.
\end{example}

\begin{proof}
The three partial derivatives $4X^3,4Y^3,-4Z^3$ have no common projective zero, so $C$ is smooth. Adjunction for a smooth plane quartic gives $g=(4-1)(4-2)/2=3$. The displayed point lies on $C$. Every product of three positive-degree cohomology classes on a curve is zero; Proposition~\ref{c:modified-diagonal} therefore gives the cohomological vanishing.

Put $J=\operatorname{Jac}(C)$ and $i:C\to J$, $x\mapsto[x-o]$. We use the precise nontorsion assertion of \citet[Theorem 4.1, printed p.~104]{c_bloch1984}: the Ceresa cycle $\Sigma=i_*[C]-[-1]_*i_*[C]$ for this Fermat quartic has infinite order modulo algebraic equivalence. We checked the detected class in its proof: the projected $5$-adic Abel--Jacobi class is nonzero by Theorem 3.4; a hypothetical torsion class would come from invariants of a $\Q_5/\Z_5$ module on which Frobenius at $11$ acts as $-11$, and $-11-1$ is a $5$-adic unit. Restriction--corestriction prevents any finite extension from killing the nontorsion class. The assertion persists over $\C$: finite rational- or algebraic-equivalence data over $\C$ spread over a finitely generated $\overline\Q$-algebra and specialize to a closed point, giving such data over a finite extension, which would contradict the theorem.

We make the passage to $M_o$ directly in rational Chow groups. For this three-dimensional Jacobian, the decomposition theorem and Proposition 3(a) of \citet[pp.~647--649]{c_beauville1986} give
\[
 \CH^2(J)_\Q=\CH^2(J)_{(0)}\oplus\CH^2(J)_{(1)}\oplus\CH^2(J)_{(2)}.
\]
The odd cycle $\Sigma$ lies entirely in grade $1$. Proposition 1(iv) of the same source therefore gives $[m]_*\Sigma=m^3\Sigma$.

For $s:C^3\to J$, $(x,y,z)\mapsto i(x)+i(y)+i(z)$, direct pushforward of the seven terms gives
\[
 s_*M_o=[3]_*i_*[C]-3[2]_*i_*[C]+3i_*[C].
\]
This is also the identity in \citet[Proposition 5.3]{c_grossschoen1995}. Taking its odd part under $[-1]_*$ gives
\[
 \frac{1-[-1]_*}{2}s_*M_o
 =\frac12\bigl([3]_* -3[2]_*+3\bigr)\Sigma.
\]
The identity just proved in the Chow group makes this odd part exactly
\[
 \tfrac12(3^3-3\cdot2^3+3)\Sigma=3\Sigma\ne0.
\]
Hence $s_*M_o$ and therefore $M_o$ are nonzero in rational Chow groups, as claimed.

The polynomial smoothness certificate, the point substitution, the genus value, and the coefficient $3^3-3\cdot2^3+3=6$ are recorded by \texttt{check\_explicit\_bad\_reduction.py}. The nontorsion input is the opened theorem and its proof, not a claim that this finite polynomial certificate computes a Chow group.
\end{proof}

\subsection{C.2: the special bivariant systems}

The relevant category must retain the two maps of manuscript 11:34. Requiring an ordinary algebra homomorphism to preserve both the unit and the Frobenius counit would change that category and would exclude the geometric maps under discussion.

\begin{definition}[A gauged special bivariant system]
\label{c:bivariant-def}
\status{Definition, with the corrected exchange axiom}
Over a characteristic-zero field $L$, a special bivariant system assigns to every finite set $I$ a finite-dimensional graded-commutative algebra $A(I)$, with $A(\varnothing)=L$, a trace $\kappa_I$ of degree $-D|I|$, and a perfect trace pairing. For a map $\alpha:I\to J$ it assigns
\[
 \alpha^*:A(I)\longrightarrow A(J),\qquad
 \alpha_*:A(J)\longrightarrow A(I).
\]
The first is a unital graded algebra homomorphism; the second has degree $D(|I|-|J|)$. Pullbacks compose covariantly in maps of indexing sets, pushforwards contravariantly, and
\[
 \kappa_I(\alpha_*x\,a)=\kappa_J(x\,\alpha^*a),\qquad
 \alpha_*(x\alpha^*a)=\alpha_*x\,a.
\]
For disjoint unions the exterior-product maps $A(I)\otimes A(J)\to A(I\amalg J)$ are isomorphisms compatible with associativity, units, traces, and the Koszul symmetry. These maps are compatible with pullback and pushforward. The exchange axiom is base change for product squares (a projection pulled back by a map on the other factor); arbitrary self-intersection squares are not included. The gauge is the chosen trace, not a quotient up to multiplication by a scalar.

A morphism in the bivariant category from $A$ to $B$ is a pair $(f^*,f_*)$, with $f^*:B\to A$ a unital graded algebra map and $f_*:A\to B$ its trace adjoint. On powers these maps are tensor powers and satisfy the corresponding pullback and pushforward compatibilities. The adjoint need not preserve the unit.
\end{definition}

\begin{theorem}[C.2: equivalence and the exchange axiom]
\label{c:bivariant-equivalence}
\status{Proved, with the category specified above}
Special gauged bivariant systems are equivalent to finite-dimensional graded-commutative Frobenius algebras of even top degree, with bivariant adjoint morphisms. The functor taking a system to its one-factor algebra and diagonal coevaluation is fully faithful. For a smooth projective connected variety, $I\mapsto H^*(X^I)$ satisfies these axioms. Correspondence composition defined by these axioms is associative and has the diagonal as identity.
\end{theorem}

\begin{proof}
From a Frobenius algebra $(A,\mu,1,\kappa)$, set $A(I)=A^{\otimes I}$. For $\alpha:I\to J$, let $\alpha^*$ multiply the inputs along each fibre $\alpha^{-1}(j)$, inserting $1$ for an empty fibre, with the Koszul permutation signs. Associativity and commutativity of $\mu$ give functoriality. Define $\alpha_*$ as its unique trace adjoint. Perfectness gives existence and uniqueness; taking adjoints proves pushforward functoriality. Pairing both sides against an arbitrary element proves the projection formula. Product-square base change follows by separating the tensor factors: contraction in the projected factor commutes with the map on the other factor. This also verifies all unit and exterior-product compatibilities.

Conversely, specialness identifies $A(I)$ with $A(\{1\})^{\otimes I}$. A fold of two indices is the product, the empty fibre gives its unit, and the trace gives a perfect pairing. Every map of finite sets is a composite of folds, injections, and bijections. Its pullback is therefore forced by this algebra; its pushforward is forced by adjointness. The coevaluation $\Delta$ is the unique tensor whose contraction is the identity. This proves reconstruction and full faithfulness. A bivariant pair on one factor extends to tensor powers, and the required compatibilities follow respectively from multiplicativity and its adjoint. We do not impose the generally false extra identity $(f^*\otimes f^*)\Delta_B=\Delta_A$.

For Weil cohomology, K\"unneth, Poincar\'e duality, and the projection formula give exactly these maps. A map $\alpha:I\to J$ induces the geometric map $X^J\to X^I$, so all variances agree with the displayed formulas. Product base change is again K\"unneth and trace in the omitted factor. In contrast, pulling a diagonal back through itself produces its Euler self-intersection, so unrestricted Cartesian base change without an excess term is false. Refined Gysin base change may be included on transverse squares, with the transversality hypothesis retained.

Finally, composition of kernels is tensor contraction in the middle factor. Two successive contractions on different factors commute by associativity of tensor product and the Koszul rule. The resulting two contractions are exactly the two parenthesizations of a triple composition. Contraction with the coevaluation is the identity by duality. This proves the associativity requested in 11:48. It also supplies the precise exchange axiom behind ``axiome cart\'esien oubli\'e !'' (11:46), and answers ``il convient de pr\'eciser les axiomes'' (11:42) on the special objects actually used here.
\end{proof}

\begin{remark}[The R7 hypotheses used in this Part]
\label{c:r7-audit}
\status{Used and rechecked, with its internal-Hom counterexample confirmed}
We checked the proof in the supplied reading, \S\S2.1--2.3, against 12:150--158. The conditional construction there is not a construction of a faithful realization on all Chow motives: it first quotients by the chosen homological equivalence, assumes its equality with numerical equivalence and the stated positivity, and then takes the semisimple envelope. Our Chow assertions above do not use this conditional passage. In particular it cannot be used to kill the homologically zero K3 determinant.

The sign correction is forced by the cross-product formula: interchanging a class of degree $a$ with one of degree $b$ multiplies it by $(-1)^{ab}$. Thus the original realization is symmetric into graded vector spaces with that commutativity constraint. This is the convention used in our Frobenius equivalence. A level filtration by thick subcategories yields a direct-sum level decomposition only when the ambient category is semisimple; then we group the simple constituents according to their least allowed twist. Thickness alone gives no splitting of extensions. The geometric level definition of C.4 therefore does not assert a canonical splitting or identify itself with coniveau.

The internal-Hom error of 12:156 already occurs in the semisimple effective Tate category with objects finite sums of $\Q(-a)$, $a\geq0$. For every effective $Z$,
\[
 \Hom(Z\otimes\Q(-1),\Q)=0.
\]
Hence the effective internal Hom from $\Q(-1)$ to $\Q$ is the zero object. After adjoining all Tate twists the corresponding internal Hom is $\Q(1)$, which is nonzero. Thus effective internal Homs need not be preserved by Tate inversion. In a rigid tensor category the valid identity is
\[
 \underline{\Hom}(M,N(-i))
 =M^\vee\otimes N(-i)
 =M^*\otimes N,
 \qquad M^*=M^\vee(-i).
\]
Evaluation and coevaluation give these isomorphisms and their inverse maps. Merely postulating the effective object $M^*$ without the duality identities does not prove them. These checks confirm precisely the restrictions of R7 needed here.
\end{remark}

\subsection{C.3: the algebraic part of the Lie object}

We separate the Lie algebra of an algebraic tannakian group from the analytic Lie algebra of a compact $\ell$-adic image. They are different objects. In a neutral tannakian category over $\Q$, the algebraic part of an adjoint Lie object means its maximal trivial subobject, computed as $\Hom(\mathbf1,\Lie G)$. The same definition is made after neutralizing a gerbe, with its descent action retained.

\begin{proposition}[The invariant-centre identity]
\label{c:lie-centre}
\status{Proved}
For a pro-reductive group $G$ over a characteristic-zero field,
\[
 \Hom_G(\mathbf1,\Lie G)=(\Lie G)^G=\Lie Z(G).
\]
If $G$ is connected, the map from its centre to its abelianization induces an isomorphism on Lie algebras. If $G^\circ$ has protorus abelianization $S$ and $\Gamma=G/G^\circ$, then
\[
 (\Lie G)^G=(\Lie S)^\Gamma.
\]
The invariant condition under $\Gamma$ cannot be discarded.
\end{proposition}

\begin{proof}
An infinitesimal element fixed by every conjugation is exactly an infinitesimal central element: the equations defining these two closed subgroup functors give the same kernel after differentiation. For a connected reductive group, $Z(G)^\circ\times G^{\mathrm{der}}\to G$ is a central isogeny. Its differential is an isomorphism, and the semisimple derived Lie algebra has no adjoint invariants. This proves the connected statement. The connected component acts trivially on its centre; the remaining action is precisely that of the component group. Passing to the inverse limit of finite-dimensional reductive quotients gives the pro-reductive assertions, since both Lie algebras and fixed-vector conditions are computed compatibly in those quotients.
\end{proof}

\begin{theorem}[Lead C3 with its arithmetic action]
\label{c:serre-centre}
\status{Proved under the stated Serre-quotient hypothesis}
Let $K$ be finitely generated over $\Q$, embedded in $\C$, and let $K_0$ be the algebraic closure of $\Q$ in $K$. Suppose that the geometric connected motivic group is pro-reductive, that its maximal protorus quotient is the full Serre protorus $S$, and that the action of $\Gal(\overline K/K)$ on this quotient is the usual action through $\Gal(\overline\Q/K_0)$. If $E$ is the largest CM subfield of $K_0$, then
\[
 \dim_\Q (\Lie G)^G=
 \begin{cases}1+[E:\Q]/2,&E\text{ exists},\\1,&K_0\text{ has no CM subfield}.
 \end{cases}
\]
These hypotheses hold for the full category of motivated motives of abelian type, including Artin motives. The conclusion concerns the full arithmetic adjoint invariants; the centre of the geometric connected group alone is generally an infinite-dimensional pro-vector space.
\end{theorem}

\begin{proof}
We first prove the character calculation rather than infer it from the dimension of a geometric torus. Write $\Q^{\mathrm{CM}}$ for the union of finite Galois CM extensions of $\Q$. Complex conjugation is central in its Galois group. The rational character space of the Serre protorus is the union of the spaces of locally constant rational-valued functions on the finite Galois quotients satisfying
\[
 f(\sigma)+f(c\sigma)=w(f),\qquad w(f)\text{ constant}.
\]
It decomposes as the one-dimensional constant space plus the space of $c$-odd functions. Taking coinvariants under $\Gamma_{K_0}$ is exact at each finite level, because averaging over the finite image identifies coinvariants with invariants in characteristic zero. At a finite Galois CM field $L$, invariant functions are functions on the embeddings of $L\cap K_0$. If this intersection is totally real, conjugation fixes every embedding and no odd function remains. If it is CM, conjugation pairs its embeddings without fixed points and the odd space has dimension $[L\cap K_0:\Q]/2$.

Every subfield of a Galois CM field is either totally real or CM: central conjugation descends to it, and the fixed field is totally real. The CM subfields of the finite number field $K_0$ have a compositum which is again CM, hence a largest one $E$ whenever any exists. Taking $L$ containing its normal closure makes the preceding dimensions stabilize at the asserted value. Thus the coinvariant character space is finite-dimensional of that dimension. Since $\Lie S=\Hom_\Q(X^*(S)\otimes\Q,\Q)$ as a pro-vector space, its fixed part is dual to those coinvariants. The preceding proposition proves the formula.

Restriction $\Gamma_K\to\Gamma_{K_0}$ is surjective: $K/K_0$ is regular, hence linearly disjoint from every finite extension of $K_0$, and the compatible extension property for finite Galois extensions gives surjectivity on inverse limits. Therefore no extra constants enter the invariant calculation.

For motivated motives of abelian type, \citet[Theorem 0.6.2 and the remark following its proof in \S6]{c_andre1996} proves that all Hodge classes on abelian varieties are motivated, including the absolute-Hodge classes over an arbitrary characteristic-zero base. Consequently the geometric Hodge realization is fully faithful on this category; compare \citet[Theorem 8.25 and Proposition 8.23]{c_delignehodge1982}. Its maximal subcategory with torus group is the category of CM Hodge structures, generated by CM abelian varieties. The Serre quotient and its arithmetic action are identified by \citet[unnumbered Theorem, p.~262, and Lemmas 1--2]{c_delignetaniyama1982}. A CM object defined over a finite extension is brought into the category over $K$ by restriction of scalars and then recovered after that extension as a direct summand; thus this assertion uses the full abelian-type category, not merely the CM varieties already having all endomorphisms over $K$. These facts establish precisely the hypotheses used in our character calculation. They also explain why ``no CM subfield'' does not mean that the geometric category contains no CM objects.
\end{proof}

\begin{example}[The finite-field analytic/algebraic distinction]
\label{c:finite-field-lie}
\status{Proved counterexample to the one-dimensional algebraic reading}
Let $E/\F_5$ have equation $y^2=x^3-x$. It is smooth since its discriminant $64$ is nonzero modulo $5$. It has eight rational points, so geometric Frobenius has polynomial $T^2+2T+5$ on $H^1$ and eigenvalues $\pi=-1+2i$, $\bar\pi=-1-2i$. Its algebraic Frobenius torus is
\[
 T=\operatorname{Res}_{\Q(i)/\Q}\Gm,
 \qquad\dim T=2.
\]
Indeed the two primes of $\Q(i)$ above $5$ give valuations $(1,0)$ and $(0,1)$ on $(\pi,\bar\pi)$; hence $\pi^a\bar\pi^b=1$ implies $a=b=0$. There is consequently no character cutting out a proper algebraic subgroup containing Frobenius. The adjoint action of $T$ is trivial, so its whole Lie algebra is algebraic, but it is two-dimensional.

On the other hand, for $\ell=3$ a suitable positive power of $\pi$ belongs to $1+3\Z_3[i]$. The analytic closure of the cyclic Frobenius subgroup has Lie algebra $\Q_3\log(\pi^N)$, of dimension one. The logarithm is nonzero because $\pi$ is not a root of unity. This is the distinction made by Serre in his letter of 8 September 1964, printed p.~202 of \citet{c_grothendieckserre}. Thus ``tout'' on 12:59 is valid for the adjoint Lie algebra of a commutative algebraic motivic band, as in the finite-field category under the Tate hypotheses of Part A; the accompanying identification with a single $\Q(0)$ does not describe that algebraic Lie object. Treating an analytic one-dimensional image as the motivic Lie algebra is an error, not an unconditional proof of a finite-field motivic assertion.
\end{example}

\subsection{C.4: weights, levels, and the overstrengthened bound}

\begin{definition}[The level convention of 12:111--115]
\label{c:level-def}
\status{Definition, checked against 12:111--115}
We fix $\ell$ invertible on the base and use $\Q_\ell(1)$ of weight $-2$. In the dimension-normalized convention, a special lisse sheaf of weight $r\geq0$ is a lisse subquotient of $R^rg_*\Q_\ell$ for a smooth projective morphism $g$ of relative dimension $r$. This includes the inserted dimension clause on 12:111; the transcription marks its following qualifier as uncertain. ``Strictly special'' means a direct summand, as clarified in 12:113. We explicitly call the version without a restriction on the relative dimension the broader convention. Over a field the two conventions agree, by the cut-and-pairing argument in Lemma~\ref{c:proper-level}; Theorem~\ref{c:general-level-bound} first uses the dimension-normalized convention, and Proposition~\ref{c:broader-coefficients} removes that clause for the general polygon and for the requested margin bound. Effective admissibility means that, after a finite stratification and the finite extensions allowed in 12:111, there is a finite filtration by extensions by zero of such sheaves. A sheaf has levels at most $\nu$ if these pieces may be chosen as $V(-a)$ with $a\geq0$, where $V$ is special of weight $r\leq\nu$. The simultaneous condition of weights in $\Phi$ requires the same pieces to satisfy $r+2a\in\Phi$. This is the geometric generation definition in 12:112; it is not simply a definition in terms of the eigenvalue weights, and it is not automatically the coniveau filtration.
\end{definition}

\begin{proposition}[C.4: the general $n-1$ bound is false]
\label{c:level-counterexample}
\status{Refuted, with the precise manuscript correction}
In the generality requested in the prompt, the bound
\[
 \operatorname{level}(R^if_*F)_{\bar y}
 \leq\min\{i,d_y(f)-1\}+j
\]
is false. Take $f:E\to\Spec\F_5$ for the elliptic curve in Example~\ref{c:finite-field-lie}, $F=\Q_\ell$, $j=0$, $i=1$, and any $\ell\ne5$. Then $d_y(f)=1$ and the right side is zero, whereas $H^1(E_{\overline\F_5},\Q_\ell)$ has level exactly one.
\end{proposition}

\begin{proof}
The sheaf $F$ is strictly special of weight zero and $E$ is regular, smooth and projective. Its nonzero $H^1$ is special of weight one, so its level is at most one. A level-zero object, by Definition~\ref{c:level-def}, has successive constituents $V(-a)$ with $V$ of weight zero and $a\geq0$; all its weights are even. The characteristic polynomial $T^2+2T+5$ has two roots of absolute value $\sqrt5$, so this $H^1$ is pure of odd weight one and cannot have level zero. The program records all eight points and the polynomial, so the counterexample does not depend on a catalogue lookup.

The supplied transcription of the final general-case box of 12:120 reads $\min(i,n)+j$, not $\min(i,n-1)+j$. The margin asks ``peut on remplacer par $n-1$ si $f$ de dim $\leq n-1$ ?'' The elliptic structural morphism has relative dimension $n$, so it does not refute that additional margin hypothesis. We therefore refute the strengthened prompt, not an assertion that the manuscript actually made without the condition.
\end{proof}

\begin{proposition}[Sharp elementary cases]
\label{c:level-special-cases}
\status{Proved in the stated cases}
If $f:X\to\Spec k$ is smooth projective of dimension $n$ and $F=\Q_\ell$, then $H^i(X)$ is pure of weight $i$ over a finite field and has level at most $\min(i,2n-i)$. Over a smooth pair $(Y,D)$ with a strict normal-crossing divisor and open immersion $j:Y\setminus D\hookrightarrow Y$, if $r$ branches meet at a geometric point $\bar y$, then
\[
 (R^ij_*\Q_\ell)_{\bar y}\simeq
 \bigwedge^i\bigl(\Q_\ell(-1)^r\bigr).
\]
It vanishes for $i>r$ and, when nonzero, has weight $2i$ and level zero. If a strictly special lisse sheaf extends across this pair, the projection formula tensors this description with its stalk.
\end{proposition}

\begin{proof}
For $i\leq n$, the cohomology itself is special of weight $i$, hence has level at most $i$. For $i>n$, hard Lefschetz identifies it with $H^{2n-i}(X)(-(i-n))$, giving the second bound. Products of ordinary elliptic curves over finite fields attain these geometric level bounds; the Frobenius-slope proof is given in Proposition~\ref{c:level-optimality}. The normal-crossing formula is local: in strict henselian coordinates each missing coordinate hyperplane supplies its Kummer class in $\Q_\ell(-1)$, and the products of distinct Kummer classes form the exterior algebra. Purity and the localization sequence prove this by induction on the number of coordinate divisors. Frobenius acts with weight two on each Kummer class. For an extending lisse sheaf, the stalkwise projection formula applies and gives the last assertion.
\end{proof}

\begin{definition}[The local dimension on 12:120]
\label{c:local-dimension}
\status{Definition, with the punctured-trait correction verified}
For a finite-type morphism $f:X\to Y$ of universally catenary noetherian schemes we use
\[
 d_y(f)=\sup_{\substack{x\in X\\ f(x)\text{ generalizes }y}}
 \left\{\operatorname{trdeg}_{k(f(x))}k(x)
       +\dim\overline{\{f(x)\}}^{\,\Spec\mathcal O_{Y,y}}\right\}.
\]
The closure in the second summand is taken in the local base. This is the dimension function, normalized to be zero at $y$, pulled back to $X$. It is the formula visible in the supplied transcription, whose notation for the local spectrum is uncertain. It measures local source dimension, not maximum fibre dimension: for $\Spec K\hookrightarrow\Spec R$ with $R$ a discrete valuation ring and $y$ its closed point it equals $1$, although all nonempty fibres have dimension zero. After an equidimensional localization of the source and a compactification, its maximum is the local dimension of that compactification. Regular alterations preserve this maximum. The dimension-function convention and the corresponding local cohomological interpretation are made precise in \citet[\S1.1.1 and Corollary 1.2.2]{c_pillonistroh2014}.
\end{definition}

\begin{lemma}[Proper cohomology with simultaneous geometric bounds]
\label{c:proper-level}
\status{Proved}
Let $Z$ be proper of dimension at most $d$ over a field $k$, and let $a\geq0$. After the finite field extensions allowed in the admissibility definition, $H^a(Z_{\bar k},\Q_\ell)$ has a finite filtration by subquotients of special objects $V_r(-t)$ satisfying
\[
 t\geq\max(0,a-d),\qquad r\geq0,\qquad r+2t\leq a.
\]
In particular its weights are at most $a$ and its geometric levels are at most $\min(a,2d-a)$. A negative upper level bound means the object is zero. These are geometric generation assertions, not deductions from a Hodge polygon.
\end{lemma}

\begin{proof}
For a smooth projective variety of dimension $d$, hard Lefschetz identifies $H^a$ with $H^{2d-a}(-(a-d))$ when $a>d$; for $a\leq d$ use $H^a$ itself. The input is \citet[Theorem 4.1.1]{c_deligneweil1980}. If a special object of degree $r$ is required to be presented by a variety of dimension exactly $r$, take a smooth complete intersection of dimension $r$ in the smooth projective variety. Weak Lefschetz injects $H^r$ into its middle cohomology. The restriction pairing is nondegenerate by hard Lefschetz on the original variety, so this injection even has a Galois-equivariant left inverse. A finite extension of $k$ suffices to choose the hyperplanes. Thus both conventions of specialness give the same objects over a field.

For an integral proper $Z$, \citet[Theorem 7.3 and Corollary 7.4]{c_dejong1996}, applied over an algebraic closure and then descended to a finite extension, give a smooth projective $W$, a finite group $G$ acting on it, and a proper map $V=W/G\to Z$ which is a universal homeomorphism over a dense open $U\subset Z$. The construction and the finite group action descend because they involve finitely many algebraic coefficients. Passing through finite purely inseparable extensions does not change the \(\ell\)-adic topos. Let $S=Z\setminus U$ and $E=V\times_ZS$. Both have dimension at most $d-1$. Proper base change identifies
\[
 H^a(V,\Q_\ell)=H^a(W,\Q_\ell)^G:
\]
the finite quotient has no higher direct images of the constant sheaf, and the invariant functions on every geometric orbit are the constants. Averaging is exact over $\Q_\ell$, even if $\ell$ divides $|G|$.

The localization triangles for $U\subset Z$ and its inverse image, together with proper base change, give the abstract blowup sequence
\[
 H^{a-1}(E)\longrightarrow H^a(Z)
 \longrightarrow H^a(V)\oplus H^a(S).
\]
The induction hypothesis for $E$ gives $t\geq\max(0,(a-1)-(d-1))$ and weight at most $a-1$; for $S$ it gives an even stronger lower bound on $t$. The smooth-projective assertion for $W$ handles $V$. Subquotients and extensions preserve all three inequalities. Reducible proper schemes follow by the same localization argument for their components; nilpotents do not change the \(\ell\)-adic topos. This completes the dimension induction.

Finally finite-extension descent does not enlarge the bounds. For a finite separable extension $k'/k$, a representation is a direct summand of the induction of its restriction, by the trace divided by $[k':k]$. A smooth projective $k'$-scheme viewed as a $k$-scheme is still smooth projective of the same dimension, and realizes that induced cohomology. Purely inseparable extensions have already been handled. This justifies the descent used above.
\end{proof}

\begin{theorem}[Corrected bounds and the margin question]
\label{c:general-level-bound}
\status{Proved in the stated arithmetic finite-type setting}
Let $f:X\to Y$ be separated of finite type, with $X$ regular and $X,Y$ of finite type over $\Z[1/\ell]$ or over a field of characteristic different from $\ell$. Fix $y\in Y$ and put $n=d_y(f)$. Let $F$ be strictly special of weight $j$, in the dimension-normalized sense of the insertion on 12:111: it is a direct summand of $R^jg_*\Q_\ell$ for a smooth projective $g$ of relative dimension $j$. Set $\epsilon=1$ if no component of local dimension $n$ maps to $y$, and set $\epsilon=0$ otherwise. Put
\[
 L=\max(0,i-n+\epsilon),\qquad U=\min(i,n)+j.
\]
The stalk $(R^if_*F)_{\bar y}$ is effectively admissible, with a filtration into pieces of weights $w$ and levels bounded simultaneously by
\begin{equation}
\label{c:weight-level-polygon}
 2L\leq w\leq2U,
 \qquad \operatorname{level}\leq\min\{w-2L,\,2U-w\}.
\end{equation}
Consequently
\[
 \operatorname{weight}\leq2\min(i,n)+2j,
 \qquad
 \operatorname{level}\leq
 \begin{cases}
 \min(i,2n-i)+j,&\epsilon=0,\\
 \min(i,2n-1-i)+j,&\epsilon=1.
 \end{cases}
\]
For an open immersion at a boundary point, $\epsilon=1$. The condition that $f$ have relative dimension at most $n-1$ near $y$ also implies $\epsilon=1$. Hence the margin's replacement by $n-1$ is valid with its stated extra condition, since $\min(i,2n-1-i)\leq\min(i,n-1)$. Without that condition the corrected general bound is at most $\min(i,n)+j$, and Proposition~\ref{c:level-counterexample} prevents replacing $n$ by $n-1$ universally.
\end{theorem}

\begin{proof}
We first take $F=\Q_\ell$ and $j=0$. Work locally on $Y$ around $y$, compactify $X$ over $Y$, and retain only the closure of $X$. Applied to that compactification and its boundary, the alteration theorems of \citet[Theorem 4.1 for varieties over fields; Theorem 8.2 for arithmetic schemes]{c_dejong1996} give a proper regular alteration $h:\widetilde X\to Y$ with a strict normal-crossing boundary $D$, and an open $v:X'\hookrightarrow\widetilde X$ above $X$. The induced $\pi:X'\to X$ is proper and generically finite. The construction can be made component by component. Localization to $\Spec\mathcal O_{Y,y}$ preserves regularity, boundary codimensions, and the local dimension function.

There are maps
\[
 \Q_{\ell,X}\longrightarrow R\pi_*\Q_{\ell,X'}
 \xrightarrow{\operatorname{Tr}_\pi}\Q_{\ell,X}
\]
whose composite is multiplication by the generic degree on each component. Here the trace is not an unjustified trace for a nonfinite map: a projective generically finite map between regular schemes is a smoothable local complete-intersection morphism of virtual dimension zero. Absolute purity identifies its Gysin map with $\Q_\ell\to\pi^!\Q_\ell$, and proper adjunction gives the displayed trace. The degree computation is checked on the generic finite-flat locus, since an endomorphism of the constant sheaf on a connected scheme is a scalar. The precise construction and calculation are \citet[Corollary 3.1.2, Lemmas 3.5.8--3.5.10, and Proposition 2.5.13]{c_riou2014}; the equality also covers inseparable generic degree by the finite-flat degree formula. Dividing by the degree proves that $Rf_*\Q_\ell$ is a direct summand of $Rh_*Rv_*\Q_\ell$.

Let $D^{(q)}$ denote the disjoint union of the intersections of $q$ distinct boundary components, and put $D^{(0)}=\widetilde X$. Absolute purity gives
\[
 R^qv_*\Q_\ell=\Q_{\ell,D^{(q)}}(-q).
\]
This is \citet[Corollary 3.1.4]{c_riou2014}, equivalently \citet[Theorem 2.1.2]{c_pillonistroh2014}; it remains true for regular arithmetic schemes. The stalkwise Leray sequence and proper base change now give
\[
 E_2^{a,q}=H^a(F_q,\Q_\ell)(-q)
 \Longrightarrow (R^{a+q}(f\pi)_*\Q_\ell)_{\bar y},
 \qquad F_q=D^{(q)}\times_Y\bar y.
\]
Every $F_q$ is proper. By the definition of $n$ and regular codimension, $\dim F_q\leq n-q$ for $q\geq1$, while $\dim F_0\leq n-\epsilon$. This step does not assume that intersecting a special fibre with a vertical divisor lowers its dimension: the inequality already holds for the entire local scheme $D^{(q)}$, of codimension $q$ in a regular local source of dimension at most $n$. The extra drop for $F_0$ follows because a proper closed fibre containing a component of maximal local dimension would make that component map to $y$. Alteration does not change the image of its generic point.

Write $a=i-q$. Lemma~\ref{c:proper-level} says that each constituent of $E_2^{a,q}$ is $V_r(-(t+q))$, with
\[
 r+2t\leq i-q,
 \qquad t\geq\max(0,i-q-\dim F_q).
\]
For $q=0$ this gives $t\geq L$. For $q\geq1$ it gives $t+q\geq q+\max(0,i-n)\geq L$. Its total weight $w=r+2(t+q)$ satisfies $w\leq i+q$, and nonvanishing requires $q\leq\min(i,2n-i)$ (as well as $q\leq n$). It follows that $w\leq2\min(i,n)$ and $r\leq w-2L$. The other inequality $r\leq2\min(i,n)-w$ follows as follows. If $i\leq n$, use $w\leq i+q$ to obtain $r=w-2(t+q)\leq2i-w$. If $i>n$, use $t\geq i-n$ for $q\geq1$, and combine $q\geq w-i$ with the same formula, obtaining $r\leq2n-w$; the $q=0$ term obeys it because $w\leq i$ and $t\geq i-n+\epsilon$. Thus every term lies in \eqref{c:weight-level-polygon}. A spectral sequence gives a filtration by subquotients of its terms; the inequalities survive extensions and direct summands. No degeneration of this Leray sequence is needed.

For the coefficient sheaf, relative hard Lefschetz and \citet[Theorem 1.5 and Proposition 2.4]{c_deligne1968} give a derived splitting
\[
 Rg_*\Q_\ell\simeq\bigoplus_bR^bg_*\Q_\ell[-b].
\]
Accordingly $R^if_*F$ is a direct summand of $R^{i+j}(fg)_*\Q_\ell$. The source of $g$ is regular, its local dimension is $n+j$, and the value of $\epsilon$ is unchanged. Applying the constant-coefficient result with $(i,n)$ replaced by $(i+j,n+j)$ gives exactly \eqref{c:weight-level-polygon} and the displayed bounds. This proves the corrected statement, including the margin condition.
\end{proof}

\begin{proposition}[Removing the dimension clause from the coefficient]
\label{c:broader-coefficients}
\status{Proved}
Keep the finite-type hypotheses of Theorem~\ref{c:general-level-bound}, but suppose only that $F$ is a direct summand of $R^jg_*\Q_\ell$ for a smooth projective $g:Z\to X$ of arbitrary relative dimension. Put $n=d_y(f)$ and $L_0=\max(0,i-n)$, $U=\min(i,n)+j$. Then the simultaneous bounds
\[
 2L_0\leq w\leq2U,\qquad
 \operatorname{level}\leq\min(w-2L_0,2U-w)
\]
still hold. In particular the general level bound is $\min(i,2n-i)+j$.
If no component of local dimension $n$ maps to $y$, there is also the bound
\[
 \operatorname{level}\leq\min(i,n-1)+j.
\]
Consequently the requested weight bound and the manuscript's margin bound hold in this broader convention as well. The stronger margin polygon with $L_0$ replaced by $\max(0,i-n+1)$ remains covered by Theorem~\ref{c:general-level-bound} in its dimension-normalized convention.
\end{proposition}

\begin{proof}
We first explain the relative dimension reduction on a dense open set. We work componentwise, so the relative dimension $r$ of $g$ is constant. If $r\leq j$, replace $Z$ by $Z\times\mathbb P^{j-r}$; K\"unneth realizes $R^jg_*\Q_\ell$ as a direct summand of the degree-$j$ cohomology of this dimension-$j$ family. If $r>j$, choose over the generic point of $X$ a smooth complete intersection of relative dimension $j$ in a projective embedding of $Z$. It spreads to a smooth projective family $h:W\to U$ over a dense open $U\subset X$. The generic field is infinite unless the source component is zero-dimensional over a finite field; in that case a finite field extension permits Bertini and restriction/corestriction gives the same direct-summand conclusion over the original field.

Weak Lefschetz gives an injection $R^jg_*\Q_\ell|_U\to R^jh_*\Q_\ell$. The pairing induced on the source by middle-dimensional Poincar\'e duality on $W/U$ is cup product with the $(r-j)$th power of the hyperplane class followed by trace. Hard Lefschetz makes it nondegenerate. Taking its adjoint and inverse gives a morphism of lisse sheaves which retracts the injection. Thus $F|_U$ has the dimension-$j$ presentation required by the preceding theorem.

We prove the general polygon by induction on the inherited local dimension. Choose the dense open $U$ just constructed and write $D=X\setminus U$. Since $X$ is regular and excellent, $D$ has a finite stratification by regular locally closed schemes $Z_\alpha$, each of pure codimension $c_\alpha\geq1$ in its source component, with a filtration by closed supports. The dimension formula gives
\[
 d_y(f|_{Z_\alpha})\leq n-c_\alpha.
\]
For nested closed supports $D'\subset D''$, the localization triangle for $R\Gamma_{D''}F$ gives the quotient term obtained by pushing the supported complex on $D''\setminus D'$ from $X\setminus D'$ to $X$. Absolute purity therefore gives, for our regular strata, successive terms
\[
 R(i_\alpha)_*(F|_{Z_\alpha})(-c_\alpha)[-2c_\alpha],
 \qquad i_\alpha:Z_\alpha\longrightarrow X.
\]
Here $i_\alpha$ is the locally closed immersion, and the pushforward is derived ordinary pushforward. Purity applies after the open restriction in which that immersion is closed; its exact form is \citet[Corollary 3.1.2]{c_riou2014}. Composing with $Rf_*$, then using the induction hypothesis on $Z_\alpha$, bounds the supported terms.

Explicitly write $c=c_\alpha$, $n'=d_y(f|_{Z_\alpha})$ and $m=i-2c$. If $m<0$ there is no term. Otherwise the induction hypothesis supplies pieces $V_a(-t_0)$ with
\[
 t_0\geq\max(0,m-n'),\qquad a+t_0\leq\min(m,n')+j.
\]
The additional twist replaces $t_0$ by $t=t_0+c$. Since $n'\leq n-c$,
\[
 t\geq\max(c,i-n)\geq L_0,
 \qquad a+t\leq\min(i-c,n)+j\leq U.
\]
These are exactly the asserted simultaneous polygon inequalities. The open term $R(f|_U)_*(F|_U)$ satisfies them by Theorem~\ref{c:general-level-bound}. The localization triangle $R\Gamma_DF\to F\to Rv_*F|_U$, with $v:U\hookrightarrow X$, now proves them for $Rf_*F$: in each cohomological degree the middle group is an extension of a quotient of the first group by a subobject of the third. Finite extensions, subquotients and filtrations preserve the two inequalities. The zero-dimensional starting case was the field reduction above. This proves the general polygon.

For the extra margin bound we use a different descent that does not introduce vertical closed strata. The smooth-complete-intersection parameter space is a nonempty open in a product of relative projective spaces and is smooth and surjective over $X$. A smooth surjection has sections after an \emph{\'etale} covering. Quasi-compactness gives a finite-type surjective \'etale map $v:V\to X$ over which the coefficient has a dimension-$j$ presentation. Denote the terms of its \v Cech nerve by $v_a:V_a\to X$. They are \'etale over $X$, hence regular, and their inherited local dimensions are at most $n$. When that dimension is $n$, no component maps to $y$; when it is smaller, it is at most $n-1$. The normalized theorem consequently gives, for every $a,b$,
\[
 \operatorname{level}\bigl(R^b(fv_a)_*v_a^*F\bigr)_{\bar y}
       \leq\min(b,n-1)+j.
\]
The \'etale \v Cech descent spectral sequence is
\[
 E_1^{a,b}=\bigl(R^b(fv_a)_*v_a^*F\bigr)_{\bar y}
       \ \Longrightarrow\ (R^{a+b}f_*F)_{\bar y}.
\]
For torsion sheaves this is \citet[Proposition 96.21.1, Tag 072G; see also Theorem 59.19.2, Tag 03OW]{c_stacksdescent}. For a lisse $\Q_\ell$-sheaf choose a stable lattice, apply it modulo $\ell^s$, and pass to the inverse limit and then invert $\ell$. Constructibility makes the finite-level stalks finite, so their inverse systems are Mittag--Leffler. Each total degree has only finitely many relevant first-quadrant terms, which justifies this passage and convergence. Since $b\leq a+b=i$, every contributing piece has level at most $\min(i,n-1)+j$. This proves the margin assertion. No properness of the \'etale cover was needed, and no stronger reflected margin inequality is inferred from this spectral sequence.
\end{proof}

\begin{proposition}[Optimality and the limit of the word sharp]
\label{c:level-optimality}
\status{Proved optimality in the stated senses}
For constant coefficients the preceding upper weight and level bounds are attained for every nonnegative degree in the indicated nonzero ranges. The general and margin level bounds remain optimal after adding $j$. The coefficient $2j$ in the universal weight envelope cannot be replaced by $j$: for every $j\geq0$ there is a strictly special weight-$j$ example with $n=i=1$ whose $R^1$ boundary stalk is pure of weight $2j+2$.
\end{proposition}

\begin{proof}
For the general level bound, take the structural map of $E^n$ over $\F_5$, where $E$ is the ordinary elliptic curve of Example~\ref{c:finite-field-lie}. Its Frobenius eigenvalues have $5$-adic slopes $0,1$. In $H^i(E^n)$ the smallest slope is $\max(0,i-n)$ and the weight is $i$. Any special constituent $V_r(-t)$ has nonnegative slopes before twisting, since eigenvalues on smooth projective cohomology are algebraic integers \citep[Expos\'e XXI, Theorem~5.2.2 and Corollary~5.5.3, pp.~385--386 and 393--395]{c_katzdeligne1973niveau}. Hence such a lowest-slope eigenvalue requires $t\leq\max(0,i-n)$, and therefore $r=i-2t\geq\min(i,2n-i)$. The upper bound is attained. Tensoring with the constant strictly special object $H^j(E^j)$, which has a slope-zero eigenvalue, raises the required level by exactly $j$.

For the margin case let $B=E^{n-1}$ and choose a sufficiently ample projectively normal embedding. Let $Y$ be its affine cone and $X=Y\setminus\{y\}$, where $y$ is the vertex. Then $X$ is the complement of the zero section of a line bundle on the smooth $B$, so it is regular of dimension $n$ and $X\hookrightarrow Y$ is an open immersion. The blowup of the vertex gives the line-bundle resolution. Purity and proper base change identify its local cohomology with the Gysin sequence for that line bundle:
\[
 \operatorname{coker}\bigl(L:H^{i-2}(B)(-1)\to H^i(B)\bigr)
 \longrightarrow (R^if_*\Q_\ell)_{\bar y}
 \longrightarrow\ker\bigl(L:H^{i-1}(B)(-1)\to H^{i+1}(B)\bigr).
\]
The two outer terms give a short exact sequence. For $i\leq n-1$, a slope-zero class in $H^i(B)$ cannot be in the image of $L$, whose eigenvalues have slopes at least one. It gives level $i$. For $n\leq i\leq2n-1$, put $r=2n-i-1$. The primitive degree-$r$ cohomology contains a slope-zero class for the same reason. Hard Lefschetz carries it to the right-hand kernel after $i-n$ powers of $L$, and the boundary twist adds one further Tate twist. This has weight $i+1$ and slope $i-n+1$, forcing level at least $2n-i-1$. Tensoring again with $H^j(E^j)$ proves sharpness with $+j$. For $n=1$ the same construction is simply $\Gm\hookrightarrow\mathbb A^1$.

For weights with $j=0$, if $i\leq n$ use the origin-boundary stalk of
$\Gm^i\times\mathbb A^{n-i}\hookrightarrow\mathbb A^n$, whose degree-$i$ Kummer class has weight $2i$. If $n<i\leq2n$, put $q=2n-i$ and $r=i-n$. The product of the degree-$q$ class of $\Gm^q$ with the top class of $\mathbb P^r$ has degree $i$ and weight $2(q+r)=2n$. Take the product open immersion followed by projection to $\mathbb A^q$; if $q=0$ this is the projective structural map. For the margin range $i\leq2n-1$ we have $q\geq1$, so the same example satisfies the extra condition.

Finally let $U=\mathbb A^1_{\F_5}\setminus\{0,1\}$ and let $\pi:\mathcal E\to U$ be the Legendre family
\[
 y^2=x(x-1)(x-t).
\]
At $t=0$ its discriminant $16t^2(1-t)^2$ has order two and $c_4=16(1-t+t^2)$ is a unit. The node has tangent slopes whose squares are $-1$, which is a square in $\F_5$, so the reduction is split multiplicative. Thus on $H^1$ inertia acts by a nonzero rank-one unipotent operator, with invariant line $\Q_\ell$ and coinvariant line $\Q_\ell(-1)$. This is the Tate-curve calculation underlying SGA~7~I, IX, Proposition 2.2.5. Take $F_j=\operatorname{Sym}^j R^1\pi_*\Q_\ell$. The K\"unneth and Schur idempotents exhibit it as a direct summand of $R^j(\mathcal E^j/U)_*\Q_\ell$, so it has the required dimension-$j$ presentation. Its inertia coinvariants are $\Q_\ell(-j)$, and tame cohomology supplies one further twist:
\[
 (R^1u_*F_j)_{\overline0}=\Q_\ell(-j-1),
 \qquad u:U\hookrightarrow\mathbb A^1.
\]
This proves the claimed weight $2j+2$. Products with the Kummer and projective examples above attain $2\min(i,n)+2j$ for every $1\leq i\leq2n-1$. The finite matrix calculation of the symmetric-power monodromy and the displayed optimization inequalities are recorded in the accompanying program.
\end{proof}

\begin{proposition}[The top degree and its sharper weight]
\label{c:top-direct-image}
\status{Proved}
In the finite-type setting of Theorem~\ref{c:general-level-bound}, let $F$ be any lisse $\Q_\ell$-sheaf. Then $(R^if_*F)_{\bar y}=0$ for $i>2n$. If no component of local dimension $n$ maps to $y$, it vanishes already for $i>2n-1$. If $k(y)$ is finite and $F$ is pure of weight $j$, the degree-$2n$ stalk is pure of weight $2n+j$ when nonzero. This endpoint is attained by $\mathbb P^n\to\Spec\F_q$ with constant coefficient $H^j(E^j)$ for an elliptic curve $E/\F_q$. These assertions do not require the dimension-normalized presentation of $F$.
\end{proposition}

\begin{proof}
We first treat an integral source component of inherited local dimension $n$ whose image is not $y$. Compactify it as $v:X\hookrightarrow\bar X\xrightarrow{h}Y$, and pass to the strict localization at $y$. For a geometric point over $x\in\bar X_y$, let $c$ be the dimension of the strict local ring of $\bar X$ at $x$. The dimension formula gives $c+\operatorname{trdeg}_{k(y)}k(x)\leq n$. For $q>0$, Gabber's bound on the cohomological dimension of an open subset of a strictly local noetherian $c$-dimensional scheme gives
\[
 (R^qv_*F)_{\bar x}\ne0\ \Longrightarrow\ q\leq2c-1;
 \quad
 \dim\operatorname{Supp}\bigl((R^qv_*F)|_{\bar X_y}\bigr)
 \leq n-\Bigl\lceil\frac{q+1}{2}\Bigr\rceil.
\]
The exact input is \citet[Theorem 1.1, formula (1.1.1), p.~455]{c_illusie2014cd}. Its torsion statement passes to lisse $\Q_\ell$-sheaves by a stable lattice and the finite constructible stalks, whose inverse systems satisfy Mittag--Leffler. For $q=0$ the support is contained in $\bar X_y$, of dimension at most $n-1$: a fibre component of dimension $n$ would be a whole irreducible component mapping to $y$.

Proper base change and the composition spectral sequence express the desired stalk through
\[
 H^a\bigl(\bar X_{\bar y},(R^qv_*F)|_{\bar X_{\bar y}}\bigr).
\]
The cohomology of a constructible sheaf supported in dimension $s$ over an algebraically closed field vanishes above $2s$. For $q>0$ every nonzero term therefore has
\[
 a+q\leq2n-2\Bigl\lceil\frac{q+1}{2}\Bigr\rceil+q\leq2n-1;
\]
for $q=0$ it has $a\leq2n-2$. This proves the asserted vanishing on this component. Components of smaller local dimension obey the same stronger bound with their own dimension.

The irreducible components of regular $X$ are disjoint. The only possible contributions in degree $2n$ thus come from its $n$-dimensional components mapping to $y$. Such a component is a smooth variety over the finite field $k(y)$ in the weight assertion. Poincar\'e duality identifies its top cohomology with $H^0_c(X_{\bar y},F^\vee)^\vee(-n)$. On a connected nonproper component, this compactly supported section space is zero: a nonzero section of a lisse sheaf has support the entire component. On a proper component it is a subspace of a fibre of $F^\vee$, hence is pure of weight $-j$. Duality gives weight $2n+j$. Components are permuted through a finite group, which does not change weights. The case $n=0$ is the same finite-fibre argument. The projective-space example gives the stated optimality.
\end{proof}

\begin{proposition}[An improved bound in degree zero]
\label{c:improved-zero-weight}
\status{Proved; optimality for general coefficients of weight at least three remains unresolved}
Let $f:X\to Y$ be separated and of finite type over $\Z[1/\ell]$
or over a finite field of characteristic different from $\ell$, with
$X$ regular, and let $y\in Y$ have finite residue field.  Let $F$ be
strictly special of weight $j\geq0$, in either convention of
Definition~\ref{c:level-def}.  Put
\[
 b(0)=0,\qquad b(1)=1,\qquad b(j)=2j-2\quad(j\geq2).
\]
Then the eigenvalues of geometric Frobenius on
$(f_*F)_{\bar y}$ have weights at most $b(j)$.
In particular, the expected bound $j$ holds for $j\leq2$.
The assertion about weights does not replace the geometric level
definition or assert a corresponding smaller level.
\end{proposition}

\begin{proof}
We first prove the local-field input, including the two extreme
monodromy maps.  Let $K$ be a finite extension of $\Q_p$ or of
$\F_p((t))$, and let $Z/K$ be smooth projective of dimension $d$.
After a finite extension and a projective semistable alteration,
the pullback on cohomology is split injective by trace divided by
the generic degree.  Thus we may use a projective strictly
semistable model.  This reduction is
\citet[Theorem~6.5, pp.~83--86]{c_dejong1996}; the assertion
descends to the original inertia invariants because Frobenius over
the extension acts there as a power of the original Frobenius.

Let $Y^{(r)}$ be the disjoint union of the $(r+1)$-fold intersections
of the special-fibre components.  We use the weight spectral sequence
\[
 E_1^{a,b}=\bigoplus_{r\geq\max(0,-a)}
 H^{b-2r}(Y^{(a+2r)}_{\bar k},\Q_\ell)(-r)
 \Longrightarrow H^{a+b}(Z_{\bar K},\Q_\ell).
\]
The $E_1^{a,b}$ are pure of weight $b$, the sequence degenerates at
$E_2$, and $N$ maps $(a,b)$ to $(a+2,b-2)$, with the matching
Tate twist.  These are
\citet[Satz~2.10 and (2.11), pp.~40--41]{c_rapoportzink1982}.
For $H^d$, all weights lie in $[0,2d]$.

The extreme map
\[
 N^d:E_2^{-d,2d}\longrightarrow E_2^{d,0}(-d)
\]
is an isomorphism.  Suppressing twists, it is the natural map
\[
 \ker\bigl(H^0(Y^{(d)})\xrightarrow{\mathrm{Gys}}
 H^2(Y^{(d-1)})\bigr)
 \longrightarrow
 \operatorname{coker}\bigl(H^0(Y^{(d-1)})
 \xrightarrow{\mathrm{Res}}H^0(Y^{(d)})\bigr).
\]
Restriction and Gysin are adjoint.  On the rational vector space
of functions on the finite set $Y^{(d)}$, the pairing is a positive
sum of products.  Its restriction to the rational image of
restriction is nondegenerate, and remains so over $\Q_\ell$.
The displayed kernel is its orthogonal complement and maps
isomorphically to the quotient.  This is the argument of
\citet[Proposition~1.5, p.~9]{c_yamamura2026weights}.

When $d\geq2$, the next extreme map
\[
 N^{d-1}:E_2^{-d+1,2d-1}\longrightarrow
 E_2^{d-1,1}(-(d-1))
\]
is also an isomorphism.  Put $U=H^1(Y^{(d-1)})$, so the components
in this expression are smooth projective curves, and put
$R=\operatorname{im}[H^1(Y^{(d-2)})\to U]$.
The signed restriction map is induced by the corresponding
homomorphism from the product of the reduced Picard varieties of
the surface strata to the product of the curve Jacobians.
Consequently $R$, with the usual Tate twist understood, is the
Tate realization of an abelian subvariety of that product.
The product principal polarization restricts to a polarization on
this abelian subvariety.  Thus the cup pairing on $U$ is
nondegenerate on $R$.  Since Gysin is the adjoint of restriction,
its kernel is $R^\perp$, and $R^\perp\to U/R$ is an isomorphism.
This proves the next extreme map, with all signs included.
We have used the Picard-polarization argument in the proof of
\citet[Satz~2.13, pp.~41--43, especially p.~42]{c_rapoportzink1982},
not an application of its surface theorem to a higher-dimensional
variety.  The resulting general-dimensional statement is
\citet[Proposition~1.7, p.~10]{c_yamamura2026weights}.

Inertia is unipotent in this model, so its invariants are $\ker N$.
An element in the kernel cannot have a nonzero weight-$2d$ component,
and for $d\geq2$ it cannot have a weight-$(2d-1)$ component either:
on each of these components a positive power of $N$ is injective.
The weight filtration is strict here, since it is the decomposition
into disjoint generalized Frobenius eigenspaces by their weights.
Hence $H^d(Z)^{I_K}$ has weights at most $2d-2$ for $d\geq2$, and
at most $1$ for $d=1$.  Degree zero has weight zero.  This recovers
the middle-degree consequences of
\citet[Theorem~1.6 and Proposition~1.8, pp.~9--11]{c_yamamura2026weights}.
The proof just given works in both local-field characteristics.
For a direct summand of $H^j$ of a smooth projective variety of
dimension at least $j$, the cut-and-pairing construction of
Lemma~\ref{c:proper-level} realizes it as a direct summand in degree
$j$ of a smooth projective $j$-fold, after a finite extension.
If the dimension is $d<j$, replace $Z$ by
$Z\times\mathbb P^{j-d}$: the $H^0$ projector on the projective
factor splits off the original $H^j$, and the product has dimension
$j$.  When $j>2d$ the original group is already zero.
Thus the same bound $b(j)$ applies.
We use only these independently proved middle-degree consequences.
The displayed higher-degree hard-Lefschetz isomorphism on p.~10
of the cited preprint prints the twist $i-d$; in our convention
the correct twist is $d-i$.  Its subsequent numerical weight bounds
use the latter sign.  This printed slip is not used in our proof.

We now reduce the degree-zero stalk to these local fields.
Write $S=\Spec\mathcal O_{Y,\bar y}^{\mathrm{sh}}$ and
\[
 V=(f_*F)_{\bar y}=H^0(X\times_Y S,F).
\]
There are finitely many connected components; they are integral
because the source is regular.  Descend these finitely many
open-and-closed components and a compactification to a finite
\'etale neighbourhood of $y$.  We do not assert that every
$\ell$-adic section descends to one such neighbourhood.
After a finite residue extension the relevant components are fixed.
The closure of each component in a proper compactification meets
the closed fibre.  Choose a closed point $c$ there.  A component
lying over $y$ can be evaluated at a finite closed point directly,
giving only weight $j$.

For the other components we construct a curve approaching $c$ from
the chosen dense open.  In the local domain $A$ of its closure at
$c$, choose $0\ne g\in\mathfrak m_A$ with $D(g)$ contained in
that open; in the horizontal mixed-characteristic case include $p$
as a factor of $g$.  Since $A$ is an excellent local domain,
the equidimensional dimension formula gives
$\dim A/(g)=\dim A-1$; thus $g$ extends to a system of parameters
$g,a_1,\ldots,a_{d-1}$ of $A$.  A minimal prime $P$ of
$(a_1,\ldots,a_{d-1})$ has height at most $d-1$ and cannot
contain $g$, since that would make $P=\mathfrak m_A$.
Moreover $g$ is primary for the maximal ideal of $A/P$, so
$\dim(A/P)=1$.  The resulting integral curve has generic point
in the required open.  Normalize, localize at a point above $c$,
and complete.  Excellence gives a discrete valuation ring with
finite residue field, whose fraction field is a finite extension
of $\Q_p$ or of $\F_p((t))$.  The entire curve construction is
performed after the preceding finite descent, over a scheme of
finite type over $\Z$ or a finite field.

Let $T$ be this trait.  The map $T^{\mathrm{sh}}\to S$, and its
generic-point map to the chosen component, pull back a section in
$V$ to an inertia-invariant vector of the restricted sheaf.
This map is equivariant for the residue Galois group after the
finite residue extension.  Evaluation at one geometric point of
a connected component is injective for a lisse sheaf.  Taking
the finite sum of these evaluations therefore embeds $V$ into
the local-field invariant spaces just considered and the pure
weight-$j$ evaluation spaces.  Since $j\leq b(j)$ for every
$j\geq0$, their weights are all at most $b(j)$.
Restriction of a strictly special sheaf retains its presenting
smooth projective family and its projector, so the local-field
input applies.  Finally the bound descends through the finite
residue extension by taking powers of Frobenius.
\end{proof}

\begin{corollary}[The first remaining lower-endpoint case]
\label{c:threefold-weight-gate}
\status{Proved equivalence for this case; the stated monodromy map remains unresolved in general}
Let $Z$ be a smooth projective threefold over a mixed-characteristic
local field with a projective strictly semistable model, and
$\ell\ne p$.  The following are equivalent:
\begin{enumerate}[label=(\roman*)]
\item $H^3(Z_{\bar K},\Q_\ell)^{I_K}$ has weights at most $3$;
\item $N:\operatorname{Gr}^W_4H^3\to
       \operatorname{Gr}^W_2H^3(-1)$ is injective;
\item the weight--monodromy conjecture holds for $H^3(Z_{\bar K},\Q_\ell)$.
\end{enumerate}
Here the weight filtration $W$ is supplied by the weight spectral
sequence, and the weight--monodromy statement means that
$N^r:\operatorname{Gr}^W_{3+r}\to
\operatorname{Gr}^W_{3-r}(-r)$ is an isomorphism for every $r\geq0$.
\end{corollary}

\begin{proof}
The preceding proof already makes $N^3$ from weight $6$ to weight
$0$, and $N^2$ from weight $5$ to weight $1$, isomorphisms.
Consequently the only possible weight above $3$ in $\ker N$ is
weight $4$.  Strictness with respect to weights identifies its
weight-$4$ part with the kernel in (ii), proving (i)$\Leftrightarrow$(ii).
Poincar\'e duality on $H^3$, with values in $\Q_\ell(-3)$, pairs
weight $4$ perfectly with weight $2$.  Thus the source and target
in (ii) have equal dimension; injectivity is equivalent to an
isomorphism.  Together with the already verified $r=2,3$ cases,
and the identity for $r=0$, this is exactly (iii).
\end{proof}


\begin{proposition}[The exact surface-pairing obstruction in dimension three]
\label{c:threefold-pairing-kernel}
\status{Proved reduction; the remaining primitive pairing is unresolved}
Let $\mathcal X$ be a projective strictly semistable model of a smooth
projective threefold $X$ over a mixed-characteristic local field with
finite residue field $k$, and let $\ell\ne\operatorname{char}k$.
Write $Y^{(r)}$ for the disjoint union, over $\bar k$, of the
$r$-fold intersections of distinct components of its special fibre.
All cohomology below has coefficients $\Q_\ell$.  Put
\[
 A=H^2(Y^{(1)}),\quad U=H^2(Y^{(2)}),\quad
 B=H^0(Y^{(3)})(-1),\quad Z=H^0(Y^{(4)})(-1).
\]
Use the signed restriction and Gysin maps of the weight spectral
sequence, and define
\[
 \begin{split}
 R&=\operatorname{im}\bigl(\rho_A:A\longrightarrow U\bigr),\\
 K_0&=R\cap\gamma_B
       \ker\bigl(\rho_B:B\longrightarrow Z\bigr),\\
 b(u,v)&=\sum_{S\subset Y^{(2)}}\operatorname{Tr}_S(u_Sv_S)
           \ \in\Q_\ell(-2).
 \end{split}
\]
Then $K_0\subset\operatorname{rad}(b|_R)$, and there is a
Galois-equivariant isomorphism
\begin{equation}
\label{eq:c-threefold-radical-kernel}
 \ker\!\left(
 N:\operatorname{Gr}^W_4H^3(X_{\bar K})(1)
       \longrightarrow\operatorname{Gr}^W_2H^3(X_{\bar K})
 \right)
 \simeq
 \frac{\operatorname{rad}(b|_R)}{K_0}.
\end{equation}
Consequently the first remaining C.4 endpoint assertion is equivalent
to nondegeneracy of the induced form on $R/K_0$.

Choose a relatively ample class $L$.  The same assertion is
equivalent to nondegeneracy of the surface intersection form on
\begin{equation}
\label{eq:c-threefold-primitive-image}
 R_{\rm prim}
 =R\cap
   \ker\bigl(
       L:H^2(Y^{(2)})\longrightarrow H^4(Y^{(2)})(1)
       \bigr).
\end{equation}
In particular the remaining condition involves only a primitive
restriction image in the degree-two cohomology of the
double-intersection surfaces.  It does not require the unknown
nondegeneracy of an additional Gysin image.
\end{proposition}

\begin{proof}
Let
\[
 C=H^4(Y^{(1)})(1),\qquad D=H^2(Y^{(3)}).
\]
The weight spectral sequence of
\citet[Satz~2.10, formulas~(2.10.1)--(2.11)]{c_rapoportzink1982},
with the same normalization as Theorem~\ref{c:threefold-weight-gate},
gives
\[
 \begin{split}
 \operatorname{Gr}^W_4H^3(1)
  &=
  \frac{\ker\{U\oplus Z\longrightarrow C\oplus D,\ 
      (u,z)\longmapsto
      (\gamma_Uu,\rho_Uu+\gamma_Zz)\}}
       {\{(\gamma_Bb,\rho_Bb):b\in B\}},\\
 \operatorname{Gr}^W_2H^3
  &=
  \frac{\ker\{U\oplus Z\longrightarrow D,\
      (u,z)\longmapsto\rho_Uu+\gamma_Zz\}}
       {\{(\rho_Aa+\gamma_Bb,\rho_Bb):
                    a\in A,\ b\in B\}}.
 \end{split}
\]
Here $\gamma_U:U\to C$ and $\gamma_Z:Z\to D$ are the
signed codimension-one Gysin maps.
These displayed identifications use degeneration at $E_2$ by
the different Frobenius weights.  The monodromy map is induced
by the identity of $U\oplus Z$.  We choose the common
restriction/Gysin sign convention in both rows, as is possible
because the differential commutes with $N$.

If a class lies in its kernel, its representative has the form
$(\rho_Aa+\gamma_Bb,\rho_Bb)$.  Subtracting the boundary
$(\gamma_Bb,\rho_Bb)$ leaves $(r,0)$ with $r=\rho_Aa\in R$.
The remaining cocycle conditions are $\gamma_Ur=0$ and
$\rho_Ur=0$.  The latter is automatic from $\rho_U\rho_A=0$.
Two such representatives differ by a domain boundary precisely
when their difference is $\gamma_Bb$ with $\rho_Bb=0$.
Thus the kernel is
\[
 (R\cap\ker\gamma_U)/
       (R\cap\gamma_B\ker\rho_B).
 \]
Poincare duality identifies $\gamma_U$ with the adjoint of
$\rho_A$ under the perfect pairing between $A$ and $C$ and
the pairing $b$ on $U$.  Hence
$\ker\gamma_U=R^\perp$.
The relation $\gamma_U\gamma_B=0$ also shows that the denominator
lies in the radical.  This proves
\eqref{eq:c-threefold-radical-kernel}, including its Galois
equivariance and all its twists.
The equivalence with the C.4 endpoint follows from
Theorem~\ref{c:threefold-weight-gate}.

For the primitive formulation we apply
\citet[Theorem~0.2, complete proofs in \S\S1.1--2.5,
and the dimension-three converse in \S2.7,
author pp.~2--9]{c_msaito2003positivity}.
His dimension-three criterion consists of
\eqref{eq:c-threefold-primitive-image} and the primitive Gysin
image in $H^3(Y^{(1)})$.
The latter condition is automatic: \S2.7 identifies it with
the degree-two monodromy isomorphism, already proved by
restriction to a general hyperplane surface and
weight--monodromy for surfaces.
The other degree-zero and degree-one pairing conditions are
settled by rational positivity and Picard-variety
polarizations, as in the proof of
Theorem~\ref{c:improved-zero-weight}.
The remaining canonical pairing is exactly the intersection
form on the primitive degree-two part of each surface,
up to a common nonzero sign.  This proves the last equivalence.
\end{proof}

\begin{corollary}[Removing the rational divisor restriction image]
\label{c:threefold-divisor-split}
\status{Proved unconditional splitting; the spectral identification uses surface Tate}
Retain the notation of Proposition~\ref{c:threefold-pairing-kernel}.
Let $N_r=\operatorname{NS}(Y^{(r)})_\Q$ for $r=1,2$,
and let $\rho_{\rm NS}:N_1\to N_2$ be signed restriction.
Define the rational space
\[
 D_{\rm prim}=
 \operatorname{im}(\rho_{\rm NS})\cap
 \ker\!\left(
 N_2\longrightarrow
 \bigoplus_{S\subset Y^{(2)}}\Q,\quad
 (d_S)_S\longmapsto(\deg(L_Sd_S))_S
 \right).
\]
We regard $D_{\rm prim}\otimes\Q_\ell(-1)$ as a subspace of
$R_{\rm prim}$ by the divisor cycle map.  Then
\[
 R_{\rm prim}
 =\bigl(D_{\rm prim}\otimes\Q_\ell(-1)\bigr)
   \mathbin{\perp}R_{\rm rem},
 \qquad
 R_{\rm rem}=
 R_{\rm prim}\cap
       \bigl(D_{\rm prim}\otimes\Q_\ell(-1)\bigr)^\perp .
\]
The first summand has a nondegenerate negative-definite
rational intersection form.  Consequently the remaining
threefold assertion is equivalent to nondegeneracy of the
form on $R_{\rm rem}$.

If the divisor Tate conjecture holds for every smooth
projective surface over finite fields of the residue
characteristic, let
$T_r\subset H^2(Y^{(r)})$ be the Frobenius spectral summand
whose eigenvalues are not $|k|$ times roots of unity.
Then
\[
 R_{\rm rem}=\rho_A(T_1)\subset T_2.
\]
Thus, under that separate surface-Tate hypothesis, the
remaining pairing is entirely on the image of the
transcendental degree-two spectral summand.
\end{corollary}

\begin{proof}
All the divisor classes and restriction maps descend to a
finite extension of $k$.  The displayed primitive condition
is a rational matrix of divisor intersection numbers.
Kernels and images commute with extension from $\Q$ to
$\Q_\ell$, so its realized space is exactly the primitive
part of the image of rational divisor classes.
The surface Hodge index theorem makes the intersection
form negative definite on this rational primitive space,
component by component; the same rationality step is
spelled out in
\citet[Claim~4.6 and its full proof, author pp.~12--13]
{c_ito2004threefolds}.
In particular its Gram determinant is nonzero in $\Q$ and
remains nonzero in $\Q_\ell$.
Orthogonal projection onto this subspace, using the inverse
Gram matrix, preserves $R_{\rm prim}$ because the subspace
lies in $R_{\rm prim}$.  Its kernel is $R_{\rm rem}$, which
proves the asserted orthogonal direct sum.  This argument
does not identify the subspace with every algebraic vector
of $R_{\rm prim}$.

For the conditional assertion, the surface-Tate hypothesis
implies the divisor Tate conjecture on all the smooth
projective strata and all their finite constant-field
extensions by
\citet[Theorem~8, pp.~89--90]{BKahn2023Divisors};
the same reduction is used in
Theorem~\ref{thm:B-relative-surface-gate}.
The divisor subspace of $H^2$ has a Galois-stable orthogonal
complement: use the ample Lefschetz intersection form,
whose restriction to divisors is nondegenerate.
The divisor subspace has eigenvalues $|k|$ times roots of
unity.  Its complement has none: otherwise after a finite
constant extension it would have a nonzero eigenvector of
eigenvalue $|k'|$, contradicting Tate surjectivity and the
direct sum.  This also excludes a generalized eigenspace
of that type, since every nonzero generalized eigenspace
contains an eigenvector.  No additional full Frobenius
semisimplicity is used.

Restriction therefore preserves the two spectral blocks.
On the surface target the second block is orthogonal to
every divisor class and hence is primitive, because the
ample class is a divisor.  We obtain
\[
 R_{\rm prim}
   =\bigl(D_{\rm prim}\otimes\Q_\ell(-1)\bigr)
        \mathbin{\perp}\rho_A(T_1).
\]
Comparing with the unconditional orthogonal decomposition
proves the formula for $R_{\rm rem}$.
\end{proof}

\begin{remark}[The geometric content not removed by the pairing criterion]
\label{c:threefold-pairing-scope}
\status{Proved scope clarification; the primitive nondegeneracy is open}
The radical formula is an exact identification of the existing
kernel.  The primitive criterion additionally removes the
nonprimitive Lefschetz pieces and the Gysin-image condition.
It does not prove the resulting nondegeneracy.
In particular, containment of a $\Q_\ell$-subspace in the span
of divisor classes does not allow one to apply rational
Hodge-index positivity to that subspace without establishing
its rational origin.  For example, over $\Q_5$, choose
$\iota^2=-1$ with $\iota\equiv2\pmod5$.
The line $\Q_5(1,\iota)$ is isotropic for $-x^2-y^2$ although
this form is negative definite over $\Q$.
It is stable under the scalar Frobenius $2I$.
This is a counterexample to that linear-algebra inference,
not a geometric counterexample to weight--monodromy.

For the same reason, an isomorphism of a stratum motive with
an abelian motive would not alone identify its geometric
primitive form with a positive polarization.  The conditional
image-projector and positivity inputs in
\citet[Theorem~0.4 and \S\S4.1--4.3]{c_msaito2003positivity}
remain separate hypotheses.  We have not silently supplied
either input from Frobenius purity or from the 2026 theorem
for CM abelian varieties.
The corresponding warning is explicit in
\citet[submitted version, Remark~4.10 and footnote~3,
p.~14]{DMilne2002Polarizations}: equality of motive categories
under a Tate hypothesis does not by itself transfer the
positivity of the geometric forms.  The published-version
Theorem~3.3 and Corollary~3.4 concern abelian varieties.
\end{remark}

\begin{proposition}[A target-only sufficient positivity hypothesis]
\label{c:threefold-target-positivity}
\status{Proved conditional refinement; the general primitive pairing remains unresolved}
Retain the strictly semistable threefold setting and fixed
$\ell\ne\operatorname{char}k$ of
Proposition~\ref{c:threefold-pairing-kernel}, and put
$V=Y^{(1)}$, $S=Y^{(2)}$.
After a finite extension of the finite residue field, we may suppose
that the geometric components are individually defined.
Assume the Hodge standard conjecture
$I(S\times S,L\boxplus L)$ and equality of $\ell$-homological
and numerical equivalence for rational cycles on $V\times S$.
Here the first assumption includes the products of distinct surface
components. Then the primitive restriction image
$R_{\rm prim}$ has nondegenerate intersection pairing.
Consequently weight--monodromy holds on $H^3$ of the generic
threefold. No Lefschetz standard conjecture on $V$ is required
in this sufficient package.

Alternatively, suppose that for each component $V_i$ an inverse to
\[
 L_i:H^2(V_i)\longrightarrow H^4(V_i)(1)
\]
is induced by an algebraic correspondence.
Choose smooth ample surfaces $j_i:T_i\hookrightarrow V_i$ with
classes $m_iL_i$, and put $T=\coprod_iT_i$.
Then the same conclusion follows from
$I(S\times S,L\boxplus L)$ and
homological equals numerical equivalence on $T\times S$.
Only the relevant homogeneous degree-two mixed trace pairings
are needed in either formulation; full equality of the two
equivalences on the displayed products is a sufficient hypothesis.
\end{proposition}

\begin{proof}
Rational algebraic degree projectors over a finite field are available
from Frobenius and
\citet[Theorem~2(1), pp.~76--77]{c_katzmessing1974}.
We therefore work with the homogeneous degree-two homological
objects $H_V,H_S$ and the algebraic restriction morphism
$f=\rho:H_V\to H_S$.
The primitive projector $p$ on the surface target is algebraic:
on each component it is the degree-two projector minus
\[
 x\longmapsto
 \frac{\operatorname{Tr}(xL)}{\operatorname{Tr}(L^2)}\,L.
\]

We first justify the target algebra and its precise pairing.
The Lefschetz standard conjecture for surfaces is known and is
stable under products. For the latter assertion, the normalized
lowering operator on a product is
$\Lambda_1\otimes1+1\otimes\Lambda_2$.
The degree operator and the $\mathfrak{sl}_2$ Casimir have finitely
many rational eigenvalues; polynomial interpolation recovers
each primitive-string level and hence the usual unnormalized
lowering operator. Thus the product lowering is algebraic.
It follows that $B(S\times S)$, and hence $A(S\times S)$, holds.
Combined with the assumed $I(S\times S)$, this gives $D(S\times S)$
as in \citet[\S4.1, pp.~14--15]{c_msaito2003positivity}.
Consequently
\[
 A=\End(H_S)
\]
identifies with the homogeneous corner of the numerical
correspondence algebra. It is finite-dimensional over $\Q$ and
semisimple by \citet[Theorem~1, pp.~448--449]
{c_jannsen1992semisimple}.
Finite-dimensionality over $\Q$ is obtained this way, not inferred
from the dimension of its $\Q_\ell$ realization.

On all of $H^2(S)$ we use the Lefschetz-modified pairing
$\Phi(x,y)=\operatorname{Tr}(x\cup *y)$, with Saito's signed
Lefschetz involution. Its adjoint $a\mapsto a^\dagger$ is algebraic,
and the assumed product positivity gives
\begin{equation}
\label{eq:c-target-positive-trace}
 \operatorname{Tr}(a^\dagger a)>0\qquad(0\ne a\in A),
\end{equation}
by \citet[equation~(4.1.3), p.~15]{c_msaito2003positivity}.
On primitive $H^2$ one has $*=-1$, so $\Phi$ is the negative
of the raw intersection pairing. They have the same radical
there. The raw intersection pairing on all of $H^2(S)$ is
not asserted to have positive adjoint trace.

For algebraic homological correspondences, put
\[
 M=\Hom(H_V,H_S),\qquad N=\Hom(H_S,H_V).
\]
The pairing
\[
 M\times N\longrightarrow\Q,\qquad
 (u,v)\longmapsto\operatorname{Tr}(uv)
\]
is nondegenerate in its first variable under $D(V\times S)$.
Indeed, in the covariant cohomological convention the two
classes have complementary codimensions $3$ and $2$ on
$V\times S$. Their composition trace is their intersection
number. Projecting a complementary class onto its homogeneous
degree-two block does not change its pairing with $u$.
Thus vanishing against every reverse block makes $u$
numerically zero, and the hypothesis makes its cohomological
class zero.

The subspace $J=fN\subset A$ is a right ideal.
Since $A$ is semisimple, $J=eA$ for an idempotent $e\in A$.
There is $g_0\in N$ with $e=fg_0$, so
$\operatorname{im}e\subset\operatorname{im}f$.
For every $g\in N$ one has $(1-e)fg=0$.
The mixed trace pairing therefore gives $(1-e)f=0$, whence
$\operatorname{im}f=\operatorname{im}e$.
This argument requires no semisimplicity of $\End(H_V)$.

The finite additive Karoubi category generated by $H_S$
is semisimple, and its faithful realization is exact because
its exact sequences split. The kernel of
\[
 (1-p)e:\operatorname{im}e\longrightarrow H_S
\]
therefore has an algebraic idempotent projector $r$ and
realizes exactly
$\operatorname{im}f\cap\operatorname{im}p=R_{\rm prim}$.
We have taken the actual intersection, not replaced it
by $p(\operatorname{im}f)$.
Its radical for $\Phi$ is
$\operatorname{im}r\cap\ker r^\dagger$, another algebraic
subobject. If nonzero, it has a nonzero idempotent projector $q$.
The pairing vanishes on $\operatorname{im}q$ with itself,
so $q^\dagger q=0$, contradicting
\eqref{eq:c-target-positive-trace}.
This proves the required nondegeneracy. The dimension-three
criterion of Proposition~\ref{c:threefold-pairing-kernel}
now proves the asserted weight--monodromy statement.

For the alternative package, the algebraic maps
\[
 r_i=m_i^{-1}L_i^{-1}j_{i*}:H^2(T_i)\longrightarrow H^2(V_i)
\]
satisfy $r_i j_i^*=1$ by the projection formula.
In this composition $j_{i*}$ has target $H^4(V_i)(1)$,
which is exactly the source of the supplied inverse.
Thus $r=\bigoplus_i r_i$ is surjective and
$\rho r$ has exactly the image of $\rho$.
Apply the preceding proof to the two surface objects $H_T,H_S$
using the assumed mixed pairing on $T\times S$.
A finite constant-field extension for the components or smooth
sections is harmless: the geometric spaces and maps are the same,
and the Frobenius eigenvalues are replaced by their powers,
which preserves all their weights.
\end{proof}

This sufficient package isolates an input that the known results
for an individual surface do not supply.
$I(S)$ controls divisor classes, whereas the required
positive adjoint trace involves codimension-two correspondence
classes on the fourfold $S\times S$.
Likewise $D(S)$ for divisors does not supply $D(T\times S)$.
A supplied degree-two inverse Hard Lefschetz on a threefold
reduces the mixed product to products of surfaces; it does not
prove either retained product hypothesis.




\begin{proposition}[A simple-spectrum substitute for target positivity]
\label{c:threefold-simple-spectrum}
\status{Proved sufficient condition; the universal threefold gate remains open}
We keep the threefold semistable notation of
Proposition~\ref{c:threefold-pairing-kernel}. After fixing a finite
residue field $\F_q$ over which all the components and the chosen ample
classes are defined, put
\[
 V=Y^{(1)},\qquad S=Y^{(2)},\qquad
 f=\rho:H^2(V,\Q_\ell)\longrightarrow H^2(S,\Q_\ell),
 \qquad P=H^2(S,\Q_\ell)_{\rm prim}.
\]
Here primitivity is imposed separately on each surface component, using
the restrictions of the fixed ample class. Assume that numerical and
$\ell$-adic homological equivalence agree on $V\times S$, and that
the characteristic polynomial of geometric Frobenius on the entire
space $P$ is squarefree. Then the cup-product pairing on
\[
 \operatorname{im}(f)\cap P
\]
is nondegenerate. Consequently the first threefold monodromy condition
in Proposition~\ref{c:threefold-pairing-kernel} holds.

The squarefreeness condition concerns the union of all surface
components, including possible coincidences between their eigenvalues.
It is not enough to impose squarefreeness separately on each component.
We do not assume the standard conjecture of Hodge type on $S\times S$,
or a semisimplicity statement for the source motive. The proof uses
algebraic projectors that are idempotent on cohomology; it does not
require idempotence modulo rational equivalence.
\end{proposition}

\begin{proof}
All algebraic correspondences in this proof have rational coefficients
and are defined over the fixed field $\F_q$. We quotient their actions
by $\ell$-adic homological equivalence. The finite-field degree
projectors are algebraic by
\citet[Theorem~2(1), pp.~76--77]{c_katzmessing1974}; hence every map
below can be represented by an algebraic correspondence homogeneous
in degree two. On a surface component $S_i$, let $h_i$ be the chosen
ample divisor class and let $d_i=\deg(h_i^2)>0$. The operator
\[
 x\longmapsto \frac{\langle x,h_i\rangle}{d_i}h_i
\]
is algebraic and is the orthogonal projector onto its hyperplane line.
Summing these projectors gives an algebraic orthogonal projector
$1-p$ with image
\[
 W=\bigoplus_i\Q_\ell h_i\simeq\Q_\ell(-1)^{\oplus c},
 \qquad\operatorname{im}(p)=P.
\]
Its rational algebraic endomorphism algebra is
$\operatorname{Mat}_c(\Q)$: products of the divisor classes provide
the matrix units. Rational intersection numbers give the converse
inclusion. The zero-surface and zero-primitive cases are immediate;
we treat the nonzero spaces below.

We first obtain the actual intersection with $P$. This step is needed
because $\operatorname{im}(pf)$ need not equal
$\operatorname{im}(f)\cap P$. Put $H_V=H^2(V,\Q_\ell)$ and
\[
 k=(1-p)f:H_V\longrightarrow W.
\]
The mixed numerical-equals-homological hypothesis implies that
\[
 \Hom_{\rm alg}(H_V,W)\times\Hom_{\rm alg}(W,H_V)
 \longrightarrow\Q,\qquad (a,b)\longmapsto\operatorname{Tr}(ab)
\]
is nondegenerate in its first variable. Indeed, embed $W$ into $H^2(S)$
and project back to $W$ by the displayed algebraic maps; a mixed
numerical test for the resulting correspondence on $V\times S$
then restricts to a test in this pairing. The same argument will
apply with $P$ in place of $W$. Only these homogeneous mixed corners
of numerical-equals-homological equivalence are used.

Let $A_W=\End_{\rm alg}(W)=\operatorname{Mat}_c(\Q)$ and form the
right ideal
\[
 J=k\Hom_{\rm alg}(W,H_V)\subseteq A_W.
\]
Because $A_W$ is semisimple, $J=eA_W$ for an idempotent $e$. Since
$e\in J$, we can choose an algebraic map $h:W\to H_V$ such that
$kh=e$. For every algebraic $b:W\to H_V$ we have
$(1-e)kb=0$, and the preceding trace nondegeneracy implies $ek=k$.
It follows that the algebraic operator
\[
 u=1-hk\in\End_{\rm alg}(H_V)
\]
satisfies $u^2=u$, $ku=0$, and $u|_{\ker k}=1$. Thus
$\operatorname{im}(u)=\ker k$ on cohomology. In particular
\[
 a=fu:H_V\longrightarrow P
 \quad\text{has image}\quad
 \operatorname{im}(a)=\operatorname{im}(f)\cap P.
\]
No algebraicity of an inverse Lefschetz map or source image splitting
has been assumed in constructing $u$.

We now determine the rational algebraic endomorphism algebra of $P$.
Write $F$ for its geometric Frobenius and let $d=\dim P$. Its
characteristic polynomial $Q(T)$ belongs to $\Q[T]$: it is obtained
from the rational characteristic polynomial on $H^2(S)$ by removing
the $c$ hyperplane eigenvalues $q$. The squarefreeness hypothesis
implies that $F$ is semisimple with distinct eigenvalues over
$\overline{\Q}_\ell$. Hence
\[
 \End_{\Q_\ell[F]}(P)=\Q_\ell[F],\qquad
 \dim_{\Q_\ell}\Q_\ell[F]=d.
\]
Every base-field algebraic endomorphism $b$ of $P$ commutes with $F$,
so write $b=\sum_{i=0}^{d-1}b_iF^i$ with $b_i\in\Q_\ell$.
For each $0\le j<d$, the trace $\operatorname{Tr}(bF^j)$ is rational:
the algebraic projector $p$ isolates degree two, so the ordinary trace
equals the alternating correspondence trace, which is the rational
intersection number with the diagonal. The matrix
\[
 G_{ij}=\operatorname{Tr}(F^{i+j}),\qquad 0\le i,j<d,
\]
has rational entries and is invertible. In fact it is the trace-pairing
matrix of the finite étale algebra $\Q[T]/Q(T)$, or equivalently the
Gram matrix of the Vandermonde matrix in the distinct Frobenius
eigenvalues. Solving this rational linear system gives every $b_i\in\Q$.
Conversely each rational polynomial in the algebraic operator $F$ is
algebraic. We have proved the equality
\[
 A_P:=\End_{\rm alg}(P)=\Q[F]\simeq\Q[T]/Q(T).
\]
This argument does not assume numerical-equals-homological equivalence
on $S\times S$.

The cup-product pairing on $P$ is nondegenerate, and its adjoint
involution satisfies
\[
 F^\dagger=q^2F^{-1}.
\]
For each complex embedding of a number-field factor of $A_P$, the
image $\alpha$ of $F$ has absolute value $q$, by purity of the smooth
proper surface cohomology
\citep[Corollary~3.3.9, p.~207]{c_deligneweil1980}.
Therefore $q^2/\alpha=\overline\alpha$.
The adjoint involution preserves each number-field factor and is
complex conjugation there. In particular every idempotent of $A_P$
is fixed by the adjoint. More strongly, the ordinary trace satisfies
\[
 \operatorname{Tr}(b^\dagger b)
   =\sum_{\sigma:A_P\to\C}|\sigma(b)|^2>0
 \qquad(b\ne0),
\]
where $P$ has multiplicity one for each embedding. This positivity
is proved for this particular rational algebra; it is not an assertion
that an arbitrary $\Q_\ell$-subspace of a polarized space is
nondegenerate.

Finally apply the right-ideal argument again, now to
\[
 a\Hom_{\rm alg}(P,H_V)\subseteq A_P.
\]
The same mixed trace nondegeneracy produces an idempotent $e'\in A_P$
and an algebraic $b:P\to H_V$ with $e'=ab$ and $e'a=a$.
Thus $\operatorname{im}(a)=\operatorname{im}(e')$. Since
$(e')^\dagger=e'$, its image is orthogonal to its kernel and the
ambient nondegenerate pairing restricts nondegenerately to the image.
This proves the required assertion about
$\operatorname{im}(f)\cap P$. We conclude using the already checked
primitive restriction-image criterion in
Proposition~\ref{c:threefold-pairing-kernel}.

If the mixed equivalence hypothesis is supplied over
$\overline{\F}_q$, it still gives the base-field trace tests used
above. A witnessing rational correspondence is defined over some
finite extension; average its finitely many Frobenius conjugates.
The first member of the trace pairing is defined over $\F_q$, so
the trace is unchanged by this averaging. No extension of the chosen
Frobenius field, and hence no unproved preservation of squarefreeness
under taking Frobenius powers, enters the argument.
\end{proof}

\begin{remark}[Limits of the simple-spectrum criterion]
\label{c:threefold-simple-spectrum-scope}
\status{Proved scope clarification; no unconditional reduction of the universal gate}
The proposition replaces the full target positivity hypothesis of
Proposition~\ref{c:threefold-target-positivity} by one directly testable
condition on a Frobenius characteristic polynomial. It retains the
mixed numerical-equals-homological hypothesis needed to obtain the
algebraic kernel and image projectors. Neither that hypothesis nor the
squarefreeness of the primitive target spectrum has been proved for
all semistable threefolds. Eigenvalues on different components can
coincide, and distinct eigenvalues can become equal after a finite
residue-field extension. We have not produced an alteration or other
transfer that removes these conditions. Thus this is a sufficient
subcase, not movement on the general C.4 assertion.
\end{remark}


\begin{lemma}[A finite exterior-power test for monodromy purity]
\label{c:exterior-weight-criterion}
\status{Proved}
Let $V$ be a finite-dimensional mixed Frobenius module over
$\Q_\ell$, with integral weights and a nilpotent operator
$N:V\to V(-1)$.  Fix an integer $m$.  We use geometric Frobenius $\phi$ and
$\phi N\phi^{-1}=q^{-1}N$.  Suppose
\[
 \sum_w w\dim V_w=m\dim V,
\]
where $V=\bigoplus_wV_w$ is the decomposition by Frobenius
weights; Frobenius need not be semisimple.  Let $h$ be the
largest length of a Jordan string of $N$, with $h=0$ if $V=0$.
Then the following conditions are equivalent:
\begin{enumerate}
\item $V$ is pure of weight $m$ in the weight--monodromy sense;
\item for every $1\leq r\leq h$, all weights of
$\ker(N:\bigwedge^r V\to\bigwedge^r V(-1))$ are at most $rm$.
\end{enumerate}
Here the exterior power is the ordinary exterior power of vector
spaces, and $N$ acts on it as a derivation.  The determinant
hypothesis holds in particular when $V$ has a nondegenerate
Frobenius-equivariant pairing
$V\otimes V\to\Q_\ell(-m)$.
\end{lemma}

\begin{proof}
The generalized Frobenius eigenspaces of a given weight have a
sum defined over $\Q_\ell$: algebraic conjugation preserves the
weight of a Weil number.  They give the asserted direct
decomposition, and the Frobenius relation implies
$N(V_w)\subset V_{w-2}$.  None of these statements requires
Frobenius semisimplicity.

We first construct homogeneous Jordan strings.  Put
$K_L=\ker N\cap\operatorname{im}N^{L-1}$.  These are graded
subspaces with $K_{L+1}\subset K_L$.  Choose a homogeneous basis
$B_L$ of a complement of $K_{L+1}$ in $K_L$, and for every
$b\in B_L$ choose a homogeneous lift $v_b$ with
$N^{L-1}v_b=b$.  The strings
\[
 v_b,Nv_b,\ldots,N^{L-1}v_b
\]
are linearly independent together.  Indeed, in a relation,
apply the largest power of $N$ equal to the distance of a
remaining term from its bottom.  The surviving bottom vectors
belong to the independent union of the $B_L$, so their
coefficients vanish; descending induction finishes the argument.
They have the correct total dimension, since
\[
 \dim K_L=\operatorname{rank}N^{L-1}
               -\operatorname{rank}N^L,\qquad
 \sum_L L(\dim K_L-\dim K_{L+1})=\dim V.
\]
Thus they form a homogeneous Jordan basis.  The individual
strings need not be Frobenius-stable.

A string of length $L$, whose top vector has weight $b$, has
weights $b,b-2,\ldots,b-2(L-1)$ and centre $c=b-L+1$.
The monodromy filtration assigns these positions the indices
$L-1,L-3,\ldots,-L+1$.  Hence weight--monodromy purity of
weight $m$ is equivalent to every string centre being $m$.

The full wedge of the string,
\[
 v_b\wedge Nv_b\wedge\cdots\wedge N^{L-1}v_b,
\]
is nonzero, is killed by $N$, and has weight
\[
 Lb-2(0+\cdots+L-1)=L(b-L+1)=Lc.
\]
Each term in its derivative has a repeated vector or is zero.
Although the wedge line need not be Frobenius-stable, the
nonzero intersection of $\ker N$ with the weight-$Lc$ subspace
is Frobenius-stable.  Condition (2), applied with $r=L$, therefore
gives $c\leq m$ for every string.  The determinant hypothesis is
$\sum_{\text{strings}}L(c-m)=0$, so every centre is $m$.

Conversely, if every centre is $m$, the homogeneous strings
carry the usual $\mathfrak{sl}_2$ structure whose lowering
operator is $N$ and whose Cartan weights are the Frobenius
weights minus $m$.  On $\bigwedge^rV$ the Cartan weights are
the Frobenius weights minus $rm$.  A vector $z$ of Cartan
weight $a$ killed by the lowering operator has $a\leq0$:
if $t$ is maximal with $e^tz\ne0$, the relation
$f e^kz=-k(a+k-1)e^{k-1}z$ at $k=t+1$ gives
$a=-t$.  Here $e$ is nilpotent because it raises the finite
weight grading.  This proves (2), in fact for every exterior power.
Finally, a perfect pairing of the indicated weight pairs
$V_w$ perfectly with $V_{2m-w}$, and gives the determinant
hypothesis by summing dimensions.

This is a direct graded-string argument.  The related
determinant criterion and invariance of purity under finite
extension are recorded, with their proofs, in
\citet[Lemma~1.4(1), (2), (5), author pp.~6--7]{c_tayloryoshida2007}.
\end{proof}

\begin{theorem}[The universal trait endpoint and weight--monodromy]
\label{c:universal-trait-weight-equivalence}
\status{Proved equivalence; the universal assertions remain unresolved}
Fix a residue characteristic $p$ and a prime $\ell\ne p$.
We consider finite extensions $K/\Q_p$, with finite residue
field, and geometric $\ell$-adic cohomology.  The following
universal assertions are equivalent:
\begin{enumerate}
\item Weight--monodromy holds for $H^j(Z_{\bar K},\Q_\ell)$
for every smooth projective $Z/K$ and every $j\geq0$.
\item For every smooth projective $j$-fold $Z/K$,
$H^j(Z_{\bar K},\Q_\ell)^{I_K}$ has weights at most $j$.
\item The assertion in (2) holds for the generic fibres of
projective strictly semistable models over $\mathcal O_K$.
\end{enumerate}
In (2) the stated space is the closed geometric stalk of
$u_*H^j(Z_{\bar K},\Q_\ell)$ for the open immersion
$u:\Spec K\hookrightarrow\Spec\mathcal O_K$.  Thus it is exactly
the dimension-normalized, degree-zero trait endpoint.

For an individual $j$-fold, after a finite extension making its
inertia action unipotent, it suffices to test the bound
$rj$ on $\ker N$ in
$\bigwedge^r H^j(Z_{\bar K},\Q_\ell)$ for $1\leq r\leq j+1$.
Each such test is a direct-summand test in the middle cohomology
$H^{rj}(Z^r_{\bar K},\Q_\ell)$.  A failure of weight--monodromy
therefore yields a trait-endpoint failure on one of these powers,
or on a projective semistable alteration of that power.
\end{theorem}

\begin{proof}
We give the transfers explicitly.  Finite extension multiplies
$N$ by a nonzero scalar if the tame character is renormalized
and replaces Frobenius weights by the same weights relative to
the larger residue cardinality.  For the latter assertion one
may pass to a common normal extension killing the finite
Weil--Deligne inertia, and choose a positive power of a
Frobenius lift lying in its Weil group.  A further power of the
new Frobenius lift differs from it by inertia acting trivially
in the Weil--Deligne representation.  Taking roots of the
corresponding absolute values proves the assertion.
Consequently weight--monodromy is invariant under finite
extension; this also verifies directly
\citet[Lemma~1.4(2)]{c_tayloryoshida2007}.

For a projective semistable model, the spectral sequence used
in Proposition~\ref{c:improved-zero-weight} has
\[
 E_1^{a,b}=
 \bigoplus_{s\geq\max(0,-a)}
 H^{b-2s}(Y^{(a+2s)}_{\bar k},\Q_\ell)(-s).
\]
Here $Y^{(t)}$ denotes the disjoint union of the $(t+1)$-fold
intersections.  On total degree $j=a+b$, a nonzero summand
requires
$b\geq2s\geq\max(0,2b-2j)$, hence $0\leq b\leq2j$.
The terms have weight $b$, so all abutment weights lie in
$[0,2j]$.  Since $N$ lowers weights by two,
$N^{j+1}=0$.  The spectral-sequence and monodromy statements
are precisely \citet[Satz~2.10 and (2.11), pp.~40--41]{c_rapoportzink1982}.
For arbitrary smooth projective $Z$, the same conclusions
follow after a projective semistable alteration and finite
extension: pullback is split by pushforward divided by the
generic degree, so its cohomology is a Frobenius- and
$N$-stable direct summand.  This is the alteration and trace
reduction already checked in Proposition~\ref{c:improved-zero-weight},
using \citet[Theorem~6.5, pp.~83--86]{c_dejong1996}.

For a $j$-fold, Poincare duality pairs $H^j$ perfectly into
$\Q_\ell(-j)$.  Lemma~\ref{c:exterior-weight-criterion}
therefore applies with largest string length at most $j+1$.
The required ordinary exterior power is a Galois direct
summand of the Kunneth factor
$H^j(Z)^{\otimes r}\subset H^{rj}(Z^r)$.  Explicitly its
projector on that factor is
\[
 \frac1{r!}\sum_{\sigma\in\mathfrak S_r}
       (\operatorname{sgn}\sigma)^{j+1}\sigma^*.
\]
The geometric permutation $\sigma^*$ already contributes
the Koszul factor $(\operatorname{sgn}\sigma)^j$.
Thus we use the geometric antisymmetrizer for even $j$ and
the geometric symmetrizer for odd $j$.  These are direct
summands of Galois representations; no algebraic
Chow--Kunneth projector is assumed.  The variety $Z^r$ has
dimension $rj$, exactly as required by the normalization.
The rational denominator $r!$ causes no exception at $\ell=2$.

Assume (3).  Given a $j$-fold $Z$, pass first to a finite
extension on which its inertia action is unipotent.  For each
$1\leq r\leq j+1$, choose, after a further finite extension,
a projective strictly semistable alteration
$\pi_r:W_r\to Z^r$ on generic fibres over some $L_r$.
We allow the fields $L_r$ to depend on $r$; no ramified
base change of a semistable integral model is presumed to
remain regular.
The maps
\[
 H^{rj}(Z^r)\xrightarrow{\pi_r^*}H^{rj}(W_r)
 \xrightarrow{(\deg\pi_r)^{-1}\pi_{r*}}H^{rj}(Z^r)
\]
have composite the identity and commute with inertia and
monodromy.  On a semistable representation, inertia invariants
are exactly $\ker N$, because its inertia action is
$\exp(t_\ell(\tau)N)$ and $t_\ell$ has infinite image.
Assumption (3) consequently bounds the weights of
$\ker N$ on each indicated exterior summand over $L_r$ by $rj$.
The original inertia action was already unipotent, so passage
to $L_r$ leaves this kernel unchanged.  Frobenius-power
comparison descends the same bound separately for every $r$.
The lemma proves weight--monodromy for $H^j(Z)$.
Finite extension does not change that conclusion.

We next reduce arbitrary cohomological degrees to this middle
case, retaining actual split summands.  Let $\dim X=n$.
If $n\geq j$, choose a smooth complete intersection
$i:Z\hookrightarrow X$ of dimension $j$, for a very ample
class $L$.  Weak Lefschetz and hard Lefschetz give the
Galois-equivariant retraction
\[
 H^j(Z)\xrightarrow{i_*}
 H^{2n-j}(X)(n-j)
 \xrightarrow{(L^{n-j})^{-1}}H^j(X),
\]
whose composite with $i^*$ is the identity.  The inverse here
is an inverse of cohomological Galois modules; its algebraicity
is unnecessary.  If $n<j$, replace $X$ by
$X\times\mathbf P^{j-n}$ and take the $H^0$ Kunneth factor
of projective space, again an actual split summand.
Pure weight--monodromy modules are closed under split
submodules: the canonical monodromy filtration is additive
on a direct sum, so each summand inherits the pure graded
pieces.  Thus (3) implies (1).

If (1) holds, a monodromy string of centre $j$ has bottom
weight $j-L+1\leq j$.  After a semistable extension this proves
the bound on $\ker N$.  The original $I_K$-invariants form
a subspace of these invariants; comparison by Frobenius
powers, on which the original inertia acts trivially, gives
the same upper bound over $K$.  Hence (1) implies (2).
Finally (2) implies (3) by restriction of its range.
\end{proof}

\begin{example}[One exterior test can miss a failure]
\label{c:single-endpoint-insufficient}
\status{Proved formal counterexample; no geometric realization asserted}
Take $K=\Q_2$, $\ell=3$, $q=2$, and an unramified Weil
representation over $\Q_3$ with basis
$u_6,v_4,x_4,y_2$, subscripts denoting weights.  Put
\[
 \phi=\operatorname{diag}(8,4,4,2),\qquad
 N(u_6)=v_4,\quad N(x_4)=y_2,\quad N(v_4)=N(y_2)=0.
\]
Its two string centres are $5$ and $3$, so it is not
weight--monodromy pure of weight $4$.  Nevertheless
its determinant has weight $16=4\dim V$ and
$\ker N=\langle v_4,y_2\rangle$ has weights at most $4$.
It even has the nondegenerate symmetric pairing into
$\Q_\ell(-4)$ specified by
\[
 \langle u_6,y_2\rangle=1,\qquad
 \langle v_4,x_4\rangle=-1
\]
and symmetric counterparts, all other entries being zero.
The identities $\phi N\phi^{-1}=2^{-1}N$,
$N^{\mathsf T}Q+QN=0$, and
$\phi^{\mathsf T}Q\phi=2^4Q$ are immediate matrix identities.
In $\bigwedge^2V$, the nonzero vector $u_6\wedge v_4$
is killed by $N$ and has weight $10>8$.
The second exterior test detects the failure.
The exact rational-matrix certificate
\path{code/part_c/check_c4_exterior_witness.py}
and its companion
\path{code/part_c/check_c4_exterior_witness_output.json}
verify these identities and the exterior-power calculation.
This is a counterexample only to the proposed linear-algebra
implication from the single invariant bound and self-duality.
It is not a smooth projective fourfold or a counterexample
to geometric weight--monodromy.
\end{example}

\begin{remark}[Scope of the trait criterion]
\label{c:trait-exterior-scope}
\status{Proved scope clarification; universal C.4 remains incomplete}
The theorem identifies the universal middle-degree trait
endpoint with weight--monodromy and gives a finite power
witness to a failure.  It does not prove either universal
assertion.  Its bound $r\leq j+1$ uses the actual nonnegative
cohomological degree $j$; it is not a bound in terms of the
weight of an arbitrary Tate-twisted motive.

The finite inertia quotient can remove a vector from the
invariants at the original place, which is why the witness
is taken after a finite semistable extension.  All local fields
are quantified in the universal assertion, so this extension
is permitted.  The argument concerns the stated traits.
For a general morphism in C.4, the component and curve-selection
argument in Proposition~\ref{c:improved-zero-weight} is an
additional required transfer; no general sheaf bound is
inferred merely by naming the trait equivalence.
\end{remark}

\begin{proposition}[The lower-endpoint gate over a trait]
\label{c:lic-endpoint}
\status{Proved exact equivalence; the general mixed-characteristic vanishing is not asserted}
Let $R$ be a complete discrete valuation ring with finite residue field $k$, fraction field $K$, and $\ell\ne\operatorname{char}(k)$. Let $\mathcal Z/R$ be regular, proper, flat and generically smooth, with generic fibre $Z$. For $u:\Spec K\hookrightarrow\Spec R$, the assertion that
\[
 (u_*H^j(Z_{\bar K},\Q_\ell))_{\bar k}
       =H^j(Z_{\bar K},\Q_\ell)^{I_K}
 \quad\text{has weights at most }j
\]
is equivalent to surjectivity of the specialization map
\[
 H^j(\mathcal Z_{\bar k},\Q_\ell)
       \longrightarrow H^j(Z_{\bar K},\Q_\ell)^{I_K}.
\]
Thus, when $\dim Z=j$, improving the universal $i=0$ coefficient bound from $2j$ to $j$ includes exactly this local invariant cycle assertion.
\end{proposition}

\begin{proof}
\citet[Theorem 10.1(a), pp.~77--80 of the opened author version]{c_flachmorin2012topos} proves that the image of specialization is precisely the weight-at-most-$j$ part of the invariant space. The source has weights at most $j$, since the special fibre is proper. To check the step in that theorem, localization and Hochschild--Serre reduce surjectivity onto that weight part to the absence of weights at most $j$ in $H^{j+1}_{\mathcal Z_{\bar k}}(\mathcal Z_{R^{\rm sh}},\Q_\ell)$. Absolute purity and duality identify this supported group with $H^{2d+1-j}(\mathcal Z_{\bar k},\Q_\ell)^\vee(-d-1)$, where $d=\dim Z$; properness bounds its weights below by $j+1$. This proves the cited image assertion and both implications. The equivalence does not prove the surjectivity in general mixed characteristic; the known low-degree and equal-characteristic cases are recorded in \citet[\S3.1, p.~494]{c_esnaultwittenberg2016}.
\end{proof}

\begin{proposition}[The lower endpoint for abelian-type coefficients]
\label{c:abelian-invariant-weights}
\status{Proved}
Let $K$ be a complete discretely valued field with finite residue field of
characteristic $p$, and let $\ell\ne p$.  Suppose that $V$ is the
$\ell$-adic realization of a pure weight-$j$ motive which, after a finite
extension of $K$, is a direct summand of tensor constructions on the
$H^1$ of abelian varieties, their duals and Tate twists.  Finite-image
Artin factors are allowed.  Then the eigenvalues of geometric Frobenius
on $V^{I_K}$ have weights at most $j$.  In particular this holds for the
local realizations of pure absolute-Hodge or motivated motives of abelian
type over number fields.  The bound is attained for good reduction.
\end{proposition}

\begin{proof}
We give the weight--monodromy argument, including the tensor step.
After a finite extension we may suppose that all the finitely many
abelian generators have semistable reduction and all the Artin factors
are trivial.  Semistable reduction is SGA~7~I, IX, Theorem~3.6, printed
p.~351.  For one such abelian variety write
$U=T_\ell(A)\otimes\Q_\ell$, $V_0=U^{I_K}$, and let $W$ be the
toric part of its Tate module.  The N\'eron special fibre gives
\[
 0\longrightarrow W\longrightarrow V_0
       \longrightarrow T_\ell(B)\otimes\Q_\ell\longrightarrow0,
\]
where $B$ is the abelian quotient of the identity component.  The exact
sequence and the identification of the fixed part are IX, formula
(2.1.8), p.~328, and Proposition~2.2.5, pp.~332--333.  A torus has Tate
module of weight $-2$, and $T_\ell(B)$ has weight $-1$.  A polarization
and the orthogonality theorem IX, Theorem~2.4, p.~338, give
$W=V_0^\perp$ in the semistable case, as stated in Proposition~3.5,
p.~350.  Consequently $U/V_0$ is dual to $W$ with values in
$\Q_\ell(1)$ and has weight $0$.

Corollary~3.5.2, p.~351, says that inertia factors through the tame
$\Z_\ell(1)$ quotient and that $(g-1)^2=0$ for a generator $g$.
Thus its logarithm $N$ satisfies $N^2=0$ and $\ker N=V_0$.
The polarization is inertia invariant, so $N$ is skew adjoint and
$\operatorname{im}N=(\ker N)^\perp=W$.  It induces an isomorphism
\[
 N:U/V_0\xrightarrow{\sim}W(-1).
\]
All these statements are rational statements; no assumption on an
integral monodromy pairing is required.  The filtration
$0\subset W\subset V_0\subset U$ is therefore the monodromy
filtration centred at weight $-1$, with graded weights $-2,-1,0$.
Dualizing proves weight--monodromy, centred at $1$, for $H^1(A)$.
This derives exactly the form we use from the cited original pages
of SGA~7 \citep{c_sga7}.

For completeness, pure Weil--Deligne representations are stable under
the indicated tensor operations.  The monodromy filtration is
characterized by $NM_a\subset M_{a-2}(-1)$ and the isomorphisms
$N^a:\operatorname{gr}_a^M\xrightarrow{\sim}
\operatorname{gr}_{-a}^M(-a)$ for $a\geq0$.  On a tensor product
the convolution filtration has these properties: decompose the
nilpotent modules into $\mathfrak{sl}_2$ strings and apply the
Clebsch--Gordan decomposition to their tensor products.  Its graded
terms have the sums of the two graded weights.  Thus the centre
weights add.  Duality reverses the indices and weights, and a Tate
twist shifts the centre weight by an even integer.  An idempotent
commuting with the Weil action and $N$ preserves the unique
monodromy filtration; its image on every graded term is pure of the
same weight.  This proves stability under direct summands as well.
We apply these facts to the weight-$j$ part of the tensor presentation
of the given pure motive.  Its Weil--Deligne representation is pure
of weight $j$.

For a monodromy string centred at $j$, the vector killed by $N$ is
its lowest-weight vector.  Equivalently, $N$ is injective on every
positive-index graded term of its monodromy filtration.  Hence
$\ker N$ has weights at most $j$.  The original inertia invariants
are a Frobenius-stable subspace of $\ker N$ after the finite
extension.  If that extension has residue degree $f$, any lift of
its geometric Frobenius differs from the $f$th power of the original
Frobenius by an element of $I_K$.  These two operators therefore
agree on $V^{I_K}$.  The weight inequality is unchanged
on taking these powers, and therefore descends to $K$.
For a pure motive with good reduction the invariant space is the
whole pure weight-$j$ representation, proving the stated sharpness.
\end{proof}


The preceding results prove the requested upper-weight inequality, the stronger displayed uniform geometric-level bounds, and the manuscript's margin question in the explicit finite-type setting. They also settle the top coefficient endpoint and, for abelian-type coefficients over a trait, the lower endpoint. Proposition~\ref{c:lic-endpoint} identifies the remaining lower-endpoint problem rather than deriving it from a theorem about ordinary direct-image weights. Proposition~\ref{c:broader-coefficients} also removes the coefficient-dimension ambiguity for the general polygon and the manuscript's stated margin bound. Only the stronger reflected margin polygon retains the dimension-normalized hypothesis in our proof.

\subsection{C.5: good reduction in codimension one}

The quantifier on primes is the one in 12:140--142: set $S_\ell=\Spec R[1/\ell]$ and test $T_\ell$ at its height-one points. It does not require the ordinary $p$-adic Tate module to be unramified at a point of residue characteristic $p$.

\begin{theorem}[The positive purity range]
\label{c:vz-positive}
\status{Known, with its application proved}
Let $(R,\mathfrak m)$ be a two-dimensional regular local ring of mixed characteristic $(0,p)$ and let $U=\Spec R\setminus\{\mathfrak m\}$. Every abelian scheme over $U$ extends uniquely over $R$ if $p\notin\mathfrak m^p$. In particular the sufficient absolute ramification bound is $e\leq p-1$, including equality. More generally the same conclusion holds under the complete faithfully flat quotient condition of Vasiu--Zink Theorem 3.

If an abelian variety over $K=\operatorname{Frac}(R)$ has good reduction at every height-one prime, these conclusions apply to it. Equivalently, at each such prime one may use the prime-to-residue-characteristic unramifiedness test of the N\'eron--Ogg--Shafarevich criterion.
\end{theorem}

\begin{proof}
The first assertion is \citet[Corollary 4]{c_vasiuzink2010}. The more general condition is the existence of a faithfully flat complete regular local $R$-algebra $\widehat R$, of the same dimension and with algebraically closed residue field $k$, admitting a surjection to
\[
 W(k)[[T_1,T_2]]/(p-h),\qquad
 h\notin(p,T_1^p,T_2^p,T_1^{p-1}T_2^{p-1}).
\]
Their Theorem 3 gives extension and uniqueness under this condition. Since $e$ is the $\mathfrak m$-adic order of $p$, $p\notin\mathfrak m^p$ means $e\leq p-1$.

To apply this to a generic abelian variety, first spread it out over a nonempty open of $\Spec R$. There are only finitely many height-one points outside this open: each lies over a height-one component of its closed complement. A model over $R_{\mathfrak p}$ spreads out to a neighbourhood of $\mathfrak p$, since it and its group structure have finite presentation. The models glue by uniqueness of homomorphism extension over normal bases; a precise form is \citet[Proposition 3.3.1.5]{c_lan2021}. They give an abelian scheme on $U$. The purity theorem then applies. The equivalence of good reduction and the indicated unramifiedness is SGA~7~I, IX, Corollary 2.2.9 \citep{c_sga7}; the criterion is pointwise over the height-one discrete valuation rings, with $\ell$ invertible there. This proves the transfer required in the question, rather than silently assuming a model already existed on $U$.
\end{proof}

\begin{theorem}[An explicit abelian-surface counterexample at $p=2$]
\label{c:explicit-bad-reduction}
\status{Proved refutation of the unrestricted criterion}
Let $k=\overline\F_2$ and
\[
 R=W(k)[[u,v]]/(2-uv),\qquad K=\operatorname{Frac}(R),\qquad
 \mathfrak m=(u,v).
\]
Consider the elliptic schemes
\begin{align*}
 E_1:&\quad Y^2Z+(u+v^2)XYZ+YZ^2=X^3,\\
 E_2:&\quad Y^2Z+(v+u^2)XYZ+YZ^2=X^3,
\end{align*}
with origins $[0:1:0]$, and put $A=E_1\times_R E_2$. The sections
\[
 P_v=[-v:1:v^3]\in E_1(R),\qquad
 P_u=[-u:1:u^3]\in E_2(R)
\]
have order two on the generic fibre. Let $C_K\subset A_K$ be the order-two subgroup generated by $(P_v,P_u)$, and define the abelian surface
\[
 B_K=A_K/C_K.
\]
Then $B_K$ has good reduction at every height-one prime of $R$, hence its $T_\ell$ is unramified at every height-one point of $S_\ell=\Spec R[1/\ell]$ for every prime $\ell$, but it does not extend to an abelian scheme over $R$. The schematic closure of $C_K$ in $A$ is explicitly nonflat.
\end{theorem}

\begin{proof}
The ambient ring $W(k)[[u,v]]$ is regular of dimension three, and $2-uv$ has nonzero linear term $2$ in its maximal ideal. The quotient is regular local of dimension two, with maximal ideal $(u,v)$ and residue field $k$. It is consequently a domain of mixed characteristic; $2=uv$ has order two in $\mathfrak m$. For $E_a:Y^2Z+aXYZ+YZ^2=X^3$, the discriminant is $a^3-27$. For both displayed values of $a$ it is $1$ modulo $\mathfrak m$, hence a unit. Thus both equations define smooth projective elliptic schemes over $R$.

For $s=v,r=u$, or $s=u,r=v$, we have $rs=2$ and $a=r+s^2$. Substitution shows that $[-s:1:s^3]$ lies on $E_a$. The inverse formula sends its $Y$-coordinate to $-1+as-s^3=-1+rs=1$, so the section equals its inverse. Its generic fibre is not the origin, hence has exact order two.

Define two finite flat group schemes
\[
 G_u=\Spec R[x]/(x^2+vx),\qquad
 G_v=\Spec R[z]/(z^2+uz),
\]
with laws $x*x'=x+x'+uxx'$ and $z*z'=z+z'+vzz'$. These are integral Hopf laws: for example $1+ux$ is group-like and $(1+ux)^2=1$, by $uv=2$ and the relation. The inverse is $x\mapsto x$; direct substitution verifies the identity. The maps
\[
 G_u\longrightarrow E_1,\quad x\longmapsto[x:1:-v^2x],\qquad
 G_v\longrightarrow E_2,\quad z\longmapsto[z:1:-u^2z]
\]
are closed immersions. On the chart $Y=1$, the coordinate $X/Y$ pulls back to $x$, respectively $z$, proving the closed-immersion assertion there; their finite source makes the images closed in the entire elliptic scheme. Their images satisfy the elliptic equations: in the first case the substitution is
\[
 (-x+v-uv^2)(x^2+vx)+v^2x(uv-2)=0,
\]
and the other case interchanges $u,v$. Generically they are the subgroups $\{0,P_v\}$ and $\{0,P_u\}$. Equality of the group-law morphisms then holds over $R$ because the source and its square are flat over the domain $R$ and the target is separated. Put $H=G_u\times G_v\subset A$.

The constant order-two group may be written
\[
 D=\Spec R[t]/(t^2+t),\qquad t*t'=t+t'+uvtt'.
\]
There is a group homomorphism $\alpha:D\to H$ defined on coordinate rings by $x\mapsto vt$, $z\mapsto ut$. The polynomial relations and the coproducts agree directly. On $D(u)$ or $D(v)$ one of these two coordinate maps recovers $t$, so $\alpha_U$ is a closed immersion. These two opens cover $U$. On the generic fibre it is exactly $C_K$.

Its schematic image has coordinate algebra
\[
 R[ut,vt]\subset R[t]/(t^2+t),\qquad
 R[ut,vt]=R\oplus\mathfrak m t
\]
as an $R$-module. To verify this equality, every nonconstant monomial in $ut,vt$ reduces using $t^2=-t$ to an element of $\mathfrak m t$, while $ut,vt$ generate that summand. This algebra is torsion-free and has generic rank two, hence it is the schematic closure of $C_K$. Its closed fibre has dimension
\[
 \dim_k\bigl((R\oplus\mathfrak m t)/\mathfrak m(R\oplus\mathfrak m t)\bigr)
 =1+\dim_k(\mathfrak m/\mathfrak m^2)=3.
\]
A finite flat module over the local ring $R$ would be free of its generic rank, and would have closed-fibre dimension two. Thus the closure is not flat.

The quotient $B_U=A_U/\alpha_U(D_U)$ is an abelian scheme, since its kernel is finite locally free. Every height-one point is in $U$, which proves good reduction there. Suppose that $B_U$ extended to an abelian scheme $B/R$. The quotient isogeny $A_U\to B_U$ and its inverse up to multiplication by $2$ extend uniquely to homomorphisms $f:A\to B$ and $g:B\to A$ by the normal-base extension theorem cited above. Their composites equal multiplication by $2$, since they do so on the dense generic fibre. Thus $f$ has finite kernel, contained in $A[2]$, and is fibrewise a surjective isogeny. It is finite and flat: properness and finite fibres give finiteness, and fibrewise flatness of isogenies between abelian varieties, together with flatness over $R$, gives flatness by the fibre criterion. Therefore $\ker f$ is finite flat. Its generic fibre is $C_K$, so torsion-freeness identifies it with the schematic closure just computed. This contradicts the closed-fibre dimension three. No such $B$ exists.

All substitutions, discriminants, inverse formulas and Hopf-map polynomial identities have been checked over $\Z$ by \texttt{check\_explicit\_bad\_reduction.py}; no calculation in that certificate inverts $2$. The example is an explicit instance of the obstruction in \citet[Lemma 27 and Theorem 28(ii)]{c_vasiuzink2010}, including the prime $2$.
\end{proof}

\begin{remark}[What fails to extend]
\label{c:subgroup-versus-object}
\status{Proved clarification}
The abstract finite group $D_U$ already extends as $D$. It is its closed immersion into $A_U$, and consequently the exact sequence defining the isogeny, that cannot extend as a finite flat subgroup. Indeed finite locally free group schemes on a punctured regular local surface extend as objects by extending their coordinate bundles to reflexive modules, which are free in dimension two; multiplication and comultiplication extend by Hartogs. Exactness and closed-immersion conditions need not survive. This is the precise content of the finite-subgroup reduction in R6, rather than a failure to extend the abstract group. The source proof of Vasiu--Zink Lemma 27 reverses the Cartier-dual arrows in its printed embedding sentence; our displayed maps avoid that typing error. Lead C5 is therefore valid as an identification of the extension obstruction, while the unrestricted good-reduction claim is false.
\end{remark}

\subsection{C.6: what independence of \texorpdfstring{$\ell$}{ell} actually supplies}

\begin{proposition}[The degree-one and semistable abelian-type cases]
\label{c:ell-independence-cases}
\status{Proved in the stated cases}
Let $K$ be a number field, $v$ a finite place of residue characteristic $p$, and $\ell\ne p$.
For $M=h^1(A)$, unramifiedness of $M_\ell$ at $v$ is independent of $\ell$ and is equivalent to good reduction of $A$. It follows for every tensor construction and motivated summand of abelian varieties already having good reduction at $v$.

More generally suppose that $A/K$ has semistable reduction at $v$, that its Galois representations take values in its Mumford--Tate group $G$ over $K$, and that a motivated abelian-type object $M$ in the category it generates corresponds to a rational representation $r:G\to\operatorname{GL}(V)$. Then unramifiedness of $M_\ell$ at $v$ is independent of $\ell\ne p$. The same conclusion holds for potentially CM motives without a semistability assumption.
\end{proposition}

\begin{proof}
For the first assertion we use SGA~7~I, IX, Proposition 2.2.5 and Corollary 2.2.9. In its notation, the dimension of inertia invariants is $\mu+2\alpha$, whereas the full dimension is $2(\mu+\lambda+\alpha)$. Equality holds if and only if $\mu=\lambda=0$, which is exactly good reduction and is a geometric condition independent of $\ell$. Smooth proper base change then makes all tensor constructions unramified, and restriction to an invariant direct summand preserves the identity action.

For the semistable assertion, \citet[Theorem 1.2]{c_kisinzhou2025} supplies a common $G$-valued Weil--Deligne representation, up to $G(\C)$-conjugacy, for every prime, including $p$. Apply the rational representation $r$. In the semistable case the finite inertial type is trivial and the actual inertia action is $\exp(t_\ell(\sigma)\,dr(N))$. It is trivial exactly when $dr(N)=0$, since the tame character has infinite image and the exponential of a nonzero nilpotent cannot be identically one in characteristic zero. This vanishing is invariant under conjugation and extension of the coefficient field, proving the assertion. No descent from a ramified field extension has been used.

For potentially CM motives, the tannakian identification with the Taniyama group gives their realizations from one representation of the Weil group; the strict-compatibility statement is \citet[Remark 3, p.~262]{c_delignetaniyama1982}. On inertia the algebraic Hecke-character factors have a common finite-order part and the Artin factors are defined over a fixed number field of coefficients. Triviality of those matrices is unchanged by any embedding of that field into $\overline\Q_\ell$. Hence unramifiedness is independent of $\ell\ne p$ in this case as well.
\end{proof}

\begin{lemma}[A generator with potentially good reduction]
\label{c:potentially-good-generator}
\status{Proved}
Let $M$ be a rational motive of abelian type over a number field $K$, in either
the absolute-Hodge or motivated category.  Let $v$ have residue characteristic
$p$.  If $M_{\ell_0}$ is potentially unramified at $v$ for one $\ell_0\ne p$,
then, after a finite Galois extension $E/K$, the object $M_E$ belongs to the
tensor category generated by finitely many abelian varieties over $E$ which
have good reduction at every place above $v$.
\end{lemma}

\begin{proof}
We first work over $\overline\Q$, using the chosen complex embedding.
The weak Mumford--Tate lifting construction of
\citet[Theorem 2.12 and Corollary 2.15, pp.~15--20]{c_noot2006} gives an
abelian variety $B$ with a central isogeny
\[
 \pi^{\mathrm{der}}:\operatorname{MT}(B)^{\mathrm{der}}
       \longrightarrow\operatorname{MT}(M)^{\mathrm{der}}
\]
identifying their adjoint Hodge homomorphisms.  Here descent to
$\overline\Q$ is part of the argument, not an assumption: start with an
abelian generator $A/\overline\Q$ of $M$, take its Mumford--Tate decomposed
weak lift $B_0$, and descend $B_0$ by Noot's Proposition 3.1 and Corollary 3.2,
pp.~20--21.  The proof of Corollary 2.15 selects the factors of $B_0$ whose
$\Q$-simple derived groups survive in $\operatorname{MT}(M)$.  These factors
descend up to isogeny because their rational endomorphism idempotents descend
under an algebraically closed field extension.  Rational motives only require
this isogeny-level descent.

The proof of Noot's Proposition 3.4, pp.~22--23, applies with the Hodge
realization of $M$ in place of $H^1(A)$ and gives
\begin{equation}
\label{c:weak-lift-generates}
 M_{\overline\Q}\ \in\
 \langle h^1(B),\text{CM motives}\rangle^{\otimes}.
\end{equation}
To make this extension of the source statement explicit, use the common
group formed from the central tori and $\operatorname{MT}(B)^{\mathrm{der}}$.
Its representation on $\omega_B(M)$ is generated by the faithful
representation on $H^1(B)$ and representations of its torus quotient.
For an irreducible constituent $W_2$ of $\omega_B(M)$, choose an irreducible
constituent $W_1$ of a tensor construction on $H^1(B)$ and a torus
representation $W_3$ with $W_2\subset W_3\otimes W_1$.
The nonzero contraction $W_2\otimes W_1^\vee\to W_3$ shows that $W_3$ itself
is an actual polarizable Hodge subquotient, of CM type, with integral
Hodge degrees.  Full faithfulness of Betti realization for abelian motives
and the realization of CM Hodge structures by CM abelian varieties prove
\eqref{c:weak-lift-generates}.  In the motivated category the same step uses
Andr\'e's Theorem 0.6.2.  Only finitely many CM varieties and morphisms are
needed for this particular $M$, so all these data descend to a finite
extension of $K$.

Put $H=\operatorname{MT}(M_{\overline\Q}\oplus h^1(B))$.
Its two projections have the same adjoint Hodge datum.  The Hodge
homomorphism defining $H$ lies in the fibre product over that adjoint group,
and projects onto both Mumford--Tate groups.  The kernel of the projection
$H\to\operatorname{MT}(M)$ lies in $\{1\}\times Z(\operatorname{MT}(B))$,
so is central in $H$ and meets its semisimple derived group in a finite
group.  Thus $H^{\mathrm{der}}\to\operatorname{MT}(M)^{\mathrm{der}}$ is a
central isogeny.  After a finite extension, the arithmetic realizations of
$M\oplus h^1(B)$ take values in $H(\Q_\ell)$ for every $\ell$: finitely many
absolute Hodge tensors cut out $H$, and one extension defines all of them.

Fix a place above $v$ and take a further finite local extension over which
$M_{\ell_0}$ is unramified and $B$ is semistable.  An inertia element acts on
$M_{\ell_0}\oplus H^1(B,\Q_{\ell_0})$ as
\[
 u=(1,\rho_{B,\ell_0}(\tau))\in H(\Q_{\ell_0}),
\]
and $u$ is unipotent by the semistable monodromy theorem.  Its image in the
torus $H/H^{\mathrm{der}}$ is trivial, so $u\in H^{\mathrm{der}}$.
Its image in $\operatorname{MT}(M)^{\mathrm{der}}$ is trivial, and hence
$u$ lies in a finite kernel.  A finite-order unipotent in characteristic
zero is the identity.  Thus $H^1(B,\Q_{\ell_0})$ is unramified, and
N\'eron--Ogg--Shafarevich gives good reduction of $B$ over this local
extension.  The same reasoning applies at \emph{every} place above $v$,
since potential unramifiedness of $M_{\ell_0}$ transfers to all those
completions.

The auxiliary CM abelian varieties are potentially good as well: after
fixing their Hodge tensors and taking a semistable extension, their inertia
is both unipotent and contained in a torus, hence trivial.  The same
criterion applies.  There are only finitely many varieties and places.
Choose finite extensions of the completions giving good reduction and
realize them within one finite global extension by weak approximation and
Krasner's lemma.  Pass to its normal closure over $K$; good reduction is
preserved by further extension.  This gives the asserted $E$ and proves
the direction of transfer required in the lemma.
\end{proof}

\begin{theorem}[Independence of unramifiedness for all abelian-type motives]
\label{c:ell-independence-abelian}
\status{Proved, using the October 2026 rationality theorem}
Let $M$ be a rational absolute-Hodge or motivated motive of abelian type
over a number field $K$, and let $v$ have residue characteristic $p$.
If $M_{\ell_0}$ is unramified at $v$ for one $\ell_0\ne p$, then
$M_\ell$ is unramified at $v$ for every $\ell\ne p$.
The statement includes $p=2$, and includes $\ell=2$ when $p\ne2$.
\end{theorem}

\begin{proof}
Apply Lemma~\ref{c:potentially-good-generator}.  Over a finite Galois
extension $E/K$, the motive $M_E$ is a summand of a finite direct sum $P$
of motives $h^{i_a}(A_a)(m_a)$, where the $A_a$ are powers or products of
abelian varieties with good reduction at \emph{every} place above $v$.
Indeed, tensor constructions on $h^1$, their duals and Tate twists occur
as K\"unneth summands in such powers.  A finite direct sum permits different
degrees $i_a$ and twists $m_a$.

Restriction and induction give morphisms
\[
 M\longrightarrow\operatorname{Ind}_E^K(M_E)\longrightarrow M
\]
whose composite is $[E:K]$.  After division by this nonzero rational
integer, $M$ is a summand of $\operatorname{Ind}_E^K(P)$.
These are finite \'etale correspondence maps in either motivic category;
after extension to $E$ they are the diagonal and summation maps over
$\Gal(E/K)$.  Thus the descent assertion is established over the original
$K$, without inferring it from trivial inertia over a ramified extension.

Write $e$ for the resulting rational Hodge, equivalently motivated,
projector.  The induced ambient varieties are finite disjoint unions of
the conjugate $A_a$.  Fix the chosen place of $\overline\Q$ above $v$ and
take a common finite Galois local extension $L/K_v$ splitting the finite
extensions of induction and giving good models for all components.
Smooth proper base change shows, simultaneously for every $\ell\ne p$,
that inertia on the ambient object factors through the same finite
quotient
\[
 \Delta=I_v/(I_v\cap G_L).
\]
It follows that this also holds on $M_\ell$.  Tate twists are trivial on
this inertia group when $\ell\ne p$.

By uniqueness of the abelian scheme models, the descent maps extend to
them.  On the geometric special fibres, each $\tau\in\Delta$ acts by a
permutation of components and origin-preserving algebraic isomorphisms
$u:X_0\to Y_0$ between conjugate abelian varieties.  Inertia fixes the
algebraic residue field, so these maps are actual homomorphisms over that
field.  Their pullbacks describe the cohomological inertia action.

The graph of every such $u$ is a rational Lefschetz class on $X_0\times Y_0$.
If $d=\dim Y_0=\dim X_0$ and $f(x,y)=y-u(x)$, then in every Weil
cohomology,
\begin{equation}
\label{c:inertia-graph-lefschetz}
 [\Gamma_u]=f^*[0_{Y_0}]
   =f^*\left(\frac{\theta_{Y_0}^{d}}
                         {\deg(\theta_{Y_0}^{d})}\right),
\end{equation}
where $\theta_{Y_0}$ is any ample divisor class.  This uses the common
degree normalization of top classes.  The expression on the right is one
rational polynomial in divisor classes, with the same cycle realizations
for all $\ell$.

Use the correspondence convention
$Z_*(x)=(\mathrm{pr}_2)_*([Z]\smile\mathrm{pr}_1^*x)$.
Then $\Gamma_u$ acts by $u_*=(u^{-1})^*$, and the alternating trace of
$e_*\circ u^*$ is the intersection pairing of $e$ with $\Gamma_u$.
For a projector supported in degree $i$, the ordinary trace is therefore
$(-1)^i$ times that pairing.  On a disjoint union, the same identity pairs
the matrix entries of $e$ with the graphs indexed by the component
permutation.  For a direct sum of different degrees and twists, only
the diagonal blocks relative to those Galois-stable summands contribute
to the trace.  Consequently
\[
 \operatorname{Tr}(\tau\mid M_\ell)
\]
is a finite sum, with the indicated signs, of pairings of specialized
rational Hodge classes with complementary rational Lefschetz classes on
products of abelian varieties with good reduction.  The generic matrix
entries of $e$ are Hodge classes on these products; realization and smooth
proper specialization preserve their actions, tensor operations, duality
and K\"unneth decompositions.

The rationality theorem \citep[Theorem 1.1]{c_rationality2026}, also
\citep[Theorem 1.1, pp.~1--2]{c_rationalityexposition2026}, says that every
such pairing is one rational number in every $\ell$-adic realization
$\ell\ne p$, including residue characteristic $2$.  Hence for each
$\tau\in\Delta$ there is a common rational trace $t_\tau$.
The assumed unramifiedness at $\ell_0$ gives
$t_\tau=\dim_\Q\omega_B(M)$.  In every other realization the operator
$\tau$ has finite order.  Its eigenvalues are roots of unity, and under a
complex embedding their sum equals their number.  Equality in the triangle
inequality forces every eigenvalue to be $1$.  Finite-order operators in
characteristic zero are semisimple, so $\tau$ is the identity on $M_\ell$.
This holds for every $\tau$, proving the theorem.
\end{proof}

\begin{remark}[The exact role of the new input]
\label{c:ell-independence-new-input}
\status{Proved scope clarification}
The proof does not apply N\'eron--Ogg--Shafarevich to an arbitrary original
generator: such a generator need not have good reduction when its summand
is unramified.  Lemma~\ref{c:potentially-good-generator} changes the
generator and proves its potential good reduction.  The finite descent
action is then recovered over $K$ by the common trace pairing.  Since
\eqref{c:inertia-graph-lefschetz} is Lefschetz, the proof uses the
rationality theorem alone; it does not use the CM Hodge theorem or assume
algebraicity of the Hodge projector.  It proves independence of
unramifiedness and, for potentially unramified objects, independence of
the finite inertia character.  It does not assert full Weil--Deligne
compatibility for objects with nonzero monodromy.
\end{remark}

% Supplementary reduction for C.6.  This does not settle the general iff.
\begin{proposition}[A purity reduction for homogeneous Chow motives]
\label{c:chow-unramified-purity-reduction}
\status{Proved conditional implication}
Let $K$ be a number field, $v$ a finite place with residue field of
cardinality $q$, and $p$ its residue characteristic.  Let
$M=(X,e,m)$ be a rational Chow motive over $K$, where $X$ is smooth
projective of dimension $d$ and the rational algebraic projector
$e\in\CH^d(X\times_KX)_{\Q}$ acts only in cohomological degree $i$.
Suppose that, for one prime $\ell_0\ne p$, the representation
$M_{\ell_0}$ is unramified at $v$ and the eigenvalues of geometric
Frobenius on it are algebraic numbers all of whose complex conjugates
have one common absolute value.  Then $M_\ell$ is unramified at $v$
for every $\ell\ne p$.

In particular, the asserted independence follows for such a motive if
the weight--monodromy conjecture holds for
$H^i(X_{\overline{K_v}},\Q_{\ell_0})$: its unramified direct summand
$M_{\ell_0}$ is then pure of weight $i-2m$.
\end{proposition}
\begin{proof}
Write $V_\ell=eH^i(X_{\overline{K_v}},\Q_\ell)(m)$.
Saito's Theorem~0.1, in its more general form Theorem~3.1,
gives a rational alternating trace independent of $\ell$ for the
composition of an algebraic correspondence and any Weil element
\citep[Theorems~0.1 and 3.1]{c_saito2003weights}.
The theorem extends from integral to rational correspondences by
linearity.  Taking the correspondence $e$, its action on all other
degrees is zero; since it is idempotent and commutes with the Galois
action, one obtains the exact equality
\[
 \operatorname{Tr}(\sigma\mid V_\ell)
 =(-1)^i q^{-m n(\sigma)}
   \sum_a(-1)^a\operatorname{Tr}
   \bigl(e\circ\sigma\mid H^a(X_{\overline{K_v}},\Q_\ell)\bigr).
\]
Thus $\operatorname{Tr}(\sigma\mid V_\ell)$ is a common rational
number for each $\sigma\in W_{K_v}$.  Apply this also to all positive
powers of $\sigma$.  Newton identities show that the characteristic
polynomial of every such $\sigma$ is in $\Q[T]$ and independent of
$\ell$.  In particular the dimension $r$ is common; if $r=0$ there
is nothing to prove.

Fix a lift $\phi\in W_{K_v}$ of geometric Frobenius.  The local
monodromy theorem supplies, for every $\ell\ne p$, a nilpotent operator
$N_\ell$ and an open inertia subgroup on which the action is
$\exp(t_\ell N_\ell)$; it also gives
\[
 F_\ell N_\ell F_\ell^{-1}=q^{-1}N_\ell,
 \qquad F_\ell=\rho_\ell(\phi).
\]
The common characteristic polynomial of $F_\ell$ has roots of the
single absolute value in the hypothesis.  Over an algebraic closure
of the coefficient field, $N_\ell$ maps a generalized
$\lambda$-eigenspace of $F_\ell$ into its generalized
$q^{-1}\lambda$-eigenspace.  The latter eigenvalue cannot occur:
under any embedding of the common algebraic splitting field into
$\C$, multiplication by $q^{-1}$ changes its absolute value.
Consequently $N_\ell=0$ for every $\ell$.

The image of inertia in each $\rho_\ell$ is therefore finite.  For
each $\tau\in I_{K_v}$ the common trace is
$\operatorname{Tr}(\rho_{\ell_0}(\tau))=r$.  A finite-order matrix
in characteristic zero is semisimple with roots of unity as
eigenvalues.  Their sum can equal the positive integer $r$ only if
all of them are $1$, as follows by embedding their cyclotomic field
into $\C$ and using equality in the triangle inequality.  Hence
$\rho_\ell(\tau)=1$ for every $\tau$ and every $\ell\ne p$.

For the last assertion, the weight--monodromy filtration is functorial
for the algebraic projector.  A pure Weil--Deligne direct summand
whose monodromy is zero is concentrated in monodromy degree zero,
so its Frobenius roots all have weight $i-2m$.
\end{proof}

\begin{remark}[Exact scope of this reduction]
\label{c:chow-purity-scope}
\status{Proved scope clarification; the general claim remains unresolved}
For the full group $H^i(X)$, the purity assumption in the proposition
is Serre's unramified-purity conjecture C7
\citep[\S2.3, printed p.~19-07]{c_serre1970local}.
For a prescribed proper summand it is the corresponding
\emph{motivic extension} of that condition; the literal original C7
must not be silently enlarged.  The named weight--monodromy
conjecture for the ambient cohomology is a sufficient alternative.
The use of Saito's theorem requires an algebraic correspondence.
It cannot be applied to an arbitrary absolute-Hodge or motivated
projector without an additional comparison result.  No converse to
any of these named conjectures is proved here, and this proposition
does not close the general-motive requirement in C.6.
\end{remark}

\begin{corollary}[Surface and low-degree Chow constructions]
\label{c:surface-lowdegree-chow-unramified}
\status{Proved}
Let $M$ be a pure rational Chow motive over a number field.  Suppose
that it is a homogeneous algebraic Chow summand of the cohomology of
a finite product of smooth projective surfaces, or more generally
belongs to the tensor constructions on supplied homogeneous algebraic
Chow summands supported in degrees $0,1,2$ of smooth projective
varieties, allowing finite sums, duals, Tate twists and homogeneous
algebraic projectors.  Then, at each finite place of residue
characteristic $p$, unramifiedness of $M_\ell$ is independent of
$\ell\ne p$.
\end{corollary}
\begin{proof}
We use the complete low-degree result
\citep[Lemma~3.9, author p.~42]{c_saito2003weights}; we spell out its
splitting step.  If $X$ is smooth projective of dimension $d\ge r$,
choose a smooth complete intersection $\iota:Y\hookrightarrow X$ of
dimension $r\le2$, after a finite extension if necessary.  For an ample
class $L$, with the degrees of the chosen hypersurfaces included in
its nonzero scalar normalization, the projection formula gives
\[
 \iota_*\iota^*=L^{d-r}:
 H^r(X)\longrightarrow H^{2d-r}(X)(d-r).
\]
Hard Lefschetz makes this an isomorphism, so
$(L^{d-r})^{-1}\iota_*$ retracts $\iota^*$ as a Galois map.  No
algebraicity of this inverse is needed or asserted here.  For curves
and surfaces, semistable alteration and the degree-divided trace
reduce purity to the Rapoport--Zink surface theorem
\citep[Satz~2.13]{c_rapoportzink1982}.  Purity descends through the
finite extension.  The cases $d<r\le2$ follow from duality or are zero.
Thus $H^r(X)$ is pure as a Weil--Deligne representation for $r\le2$.
Poincar\'e duality gives the corresponding statement in degrees
$2d,2d-1,2d-2$, and therefore in every degree of a surface.

Purity is preserved by tensor products, direct sums of the same
weight, duals, Tate twists and direct summands.  K\"unneth gives the
required purity for the stated constructions.  For finite sums, insert
the algebraic unit and point projectors of projective spaces to put
the summands into one cohomological degree and one common Tate twist
on a disjoint union of smooth projective varieties of the same
dimension.  Proposition~\ref{c:chow-unramified-purity-reduction}
then applies to the resulting homogeneous algebraic projector.  In particular it preserves the
finite inertia information at the original place; there is no
inference of unramifiedness by ramified descent.
\end{proof}

This supplies the expected local C.4 bound
$\operatorname{weights}(M_\ell^{I_v})\le j$ for a pure weight-$j$
coefficient in these tensor constructions: inertia invariants lie
in $\ker N$, which consists of the bottom vectors of pure monodromy
strings.  The first remaining general homogeneous Chow purity case
is $H^3$ of arbitrary smooth projective threefolds.  The assumed
homogeneous projectors on arbitrary higher-dimensional varieties
are part of the statement; none is manufactured from a possibly
nonalgebraic hard-Lefschetz inverse.

\begin{corollary}[Chow tensor summands of hyper-K\"ahler varieties]
\label{c:hyperkahler-chow-unramified}
\status{Proved using the stated 6 October 2026 purity theorem}
Let $M$ be a pure rational Chow motive over a number field, obtained
from products of hyper-K\"ahler varieties by homogeneous algebraic
projectors, finite sums, duals and Tate twists.  At a finite place
$v$ of residue characteristic $p$, if $T_{\ell_0}(M)$ is unramified
for one $\ell_0\ne p$, then $T_\ell(M)$ is unramified for every
$\ell\ne p$.
\end{corollary}

\begin{proof}
\citet[Theorem~4.4, with Lemmas~4.2--4.3]{c_itozou2026}
prove weight--monodromy in every degree for every hyper-K\"ahler
variety over a $p$-adic field, including $\ell=p$.  Their theorem
is over the original local field: the finite extension in its proof
does not remain in its conclusion, because purity descends through
finite extensions.  Pure Weil--Deligne representations are stable
under tensor products, duals, Tate twists and direct summands.
K\"unneth therefore gives the ambient purity hypothesis in Proposition~\ref{c:chow-unramified-purity-reduction}.  At $\ell_0$, zero monodromy gives
a single Frobenius weight on the chosen pure summand.  That reduction,
using Saito's algebraic-correspondence trace theorem, proves the
assertion at the original place $v$.
\end{proof}

This consequence uses the purity theorem, not a descent of
unramifiedness through a ramified extension.  The stronger
group-valued compatibility statement in
\citet[Theorem~5.8 and Remark~5.9]{c_itozou2026} assumes connected
monodromy and a semistable Kuga--Satake variety over the chosen base;
Corollary~5.10 obtains its general local Frobenius comparison after
a finite extension.  The displayed corollary retains algebraic
projectors.  The same purity theorem gives the local C.4 bound
$\operatorname{weights}(\ker N)\le j$ for pure weight-$j$
tensor summands of hyper-K\"ahler cohomology.  It does not assert the missing trace comparison for
arbitrary absolute-Hodge projectors outside the abelian-type case.

\begin{proposition}[A finite-field rank conjecture controls algebraic homological projectors]
\label{c:chow-unramified-rank-reduction}
\status{Proved conditional implication; no converse is asserted}
Fix a prime $p$. Assume independence of $\ell$ of homological
equivalence for rational cycles on every smooth projective variety
over $\overline{\F}_p$. Equivalently, assume independence of the
ranks of all rational algebraic correspondences on their individual
cohomology groups; this is the exact equivalence of
\citet[Theorem~1, Definition~4.1 and Theorem~4.6]{c_vandobben2020}.

Let $K$ be a number field, $v\mid p$, and let $X/K$ be smooth
projective of pure dimension $d$. Let
$e\in\CH^d(X\times_KX)_\Q$ be a rational algebraic correspondence
whose Betti action is idempotent and vanishes outside degree $i$.
For an integer $m$, put
$M_\ell=\operatorname{im}(e:H^i(X_{\bar K},\Q_\ell)(m)
\to H^i(X_{\bar K},\Q_\ell)(m))$.
Betti--\'etale comparison gives the same homogeneous idempotent
condition for every $\ell$. We require no identity $e^2=e$ in
the Chow ring; the assertion therefore covers algebraic
homological projectors as well as rational Chow motives.
Then the ranks of all powers of the local monodromy operator on
$M_\ell$ are independent of $\ell\ne p$.
Moreover,
\[
 M_{\ell_0}\text{ unramified at }v\text{ for one }\ell_0\ne p
 \quad\Longrightarrow\quad
 M_\ell\text{ unramified at }v\text{ for every }\ell\ne p.
\]
This implication does not assume weight--monodromy.
\end{proposition}

\begin{proof}
We work over the local field $K_v$ and first suppress the Tate twist:
for $\ell\ne p$ it has trivial inertia and zero monodromy.
There are three linear-algebra and spectral-sequence points to verify.

\emph{1. Ranks on cohomology are finite block-matrix ranks.}
Let $f:C^\bullet\to D^\bullet$ be a map of finite complexes over a
field. For every $a$,
\begin{equation}
\label{eq:c-induced-cohomology-rank}
 \operatorname{rank}H^a(f)
 =\operatorname{rank}
 \begin{pmatrix}
  d_C^a&0\\ f^a&d_D^{a-1}
 \end{pmatrix}
 -\operatorname{rank}d_C^a-\operatorname{rank}d_D^{a-1}.
\end{equation}
Indeed, projection of the block image onto $C^{a+1}$ has image
$\operatorname{im}d_C^a$ and kernel
$f^a(\ker d_C^a)+\operatorname{im}d_D^{a-1}$.
Dividing that kernel by $\operatorname{im}d_D^{a-1}$ gives exactly
the image on cohomology, since $f$ carries boundaries to boundaries.

The rank conjecture applies to these block matrices whenever their
entries are algebraic correspondences between finite sums of twisted
stratum cohomology groups of one weight. To check this assertion
explicitly, write a summand as $H^{j_s}(Z_s)(t_s)$ with
$j_s-2t_s=w$. Choose $T\geq\max_s t_s$ and put
$r_s=T-t_s\geq0$. The summand identifies with the part of
\[
 H^{w+2T}(Z_s\times\mathbf P^{r_s})(T)
\]
selected by the algebraic top-degree projector of the projective
space. Exterior products of the original correspondences with the
projective-space correspondences implement every block.
Choose a common dimension $D$ at least as large as every
$\dim Z_s+r_s$, and further multiply the $s$th component by
$\mathbf P^{D-\dim Z_s-r_s}$ with its algebraic degree-zero
projector. This leaves the selected cohomological degree unchanged
and makes all components equidimensional of dimension $D$.
Disjoint unions implement finite direct sums, and the top-degree
and degree-zero projectors kill the unwanted projective-space
components. Thus the block map is a single codimension-$D$
rational algebraic correspondence
on a fixed cohomological degree of smooth projective schemes.
Its rank is independent of $\ell$ by the stated conjecture.
No algebraic projector onto a kernel of a differential is required.

\emph{2. A semistable projector has common algebraic $E_1$ maps.}
Choose a finite Galois extension $L/K_v$, a projective strictly
semistable scheme $\mathcal Y/\mathcal O_L$, and a proper
generically finite map $f:Y\to X_L$ of positive degree $\delta$.
Such an alteration is supplied by
\citet[Corollary~1.2.6, author p.~16]{c_saito2003weights}.
For geometric components, take disjoint unions of alterations and,
if necessary, repeat components to make the generic degree common.
The correspondence
\[
                 e_Y=\delta^{-1}f^* e f_*
\]
is algebraic and acts idempotently on cohomology. The identities $f_*f^*=\delta$ and
the projection formula identify its realized image with that of
$e$, compatibly with Galois action and monodromy.
Passing from $K_v$ to $L$ multiplies the normalized monodromy
operator by the nonzero ramification index, so none of its ranks
changes.

Use Saito's convention that $Y^{(s)}$ is the disjoint union of the
$(s+1)$-fold intersections of geometric components. The weight
spectral sequence is
\[
 E_{1,\ell}^{a,b}
 =\bigoplus_{j\geq\max(0,-a)}
 H^{b-2j}(Y^{(a+2j)}_{\overline k},\Q_\ell)(-j)
 \ \Longrightarrow\ H^{a+b}(Y_{\overline L},\Q_\ell).
\]
Its differentials are fixed signed restriction and Gysin
correspondences. The action of $e_Y$ on this spectral sequence
has $E_1$ components induced by fixed rational algebraic
correspondences on the strata, independent of $\ell$.
This is precisely
\citet[Lemma~2.4.2, Corollary~2.4.3 and Proposition~2.4.4,
complete proofs, author pp.~34--36]{c_saito2003weights}:
one resolves the closure of a cycle on the regular semistable
product blowup by vector bundles, restricts their Chern characters
to the strata, and divides by the fixed factorial.
The resulting rational cycles are chosen before $\ell$.
The $E_1$ endomorphism need not be idempotent; its induced action
on $E_2$ and the abutment is the required projector.

On $E_1$, the monodromy map
\[
 N:E_{1,\ell}^{a,b}\longrightarrow
       E_{1,\ell}^{a+2,b-2}(-1)
\]
is the identity on matching summands and zero on the others.
It sends $(a,j)$ to $(a+2,j-1)$ and hence preserves the stratum
index $a+2j$. Consequently the stratum correspondence blocks
commute with this map. The assertion for $N$ follows from
\citet[Corollary~2.2.4(2), author pp.~21--22]{c_saito2003weights}:
the displayed map for $T-1$ differs on the graded pieces from
$N$ only by its nonzero tame scalar. Higher logarithm terms drop
the filtration further and do not change this leading map.

Fix $b$ and an integer $r\geq1$. Regard
$C^a=E_{1,\ell}^{a,b}$ and
$D^a=E_{1,\ell}^{a+2r,b-2r}(-r)$ as complexes in $a$.
All their summands have weight $b$, and
$N^r e_Y:C^\bullet\to D^\bullet$ has rational algebraic
stratum blocks. Formula~\eqref{eq:c-induced-cohomology-rank}
therefore shows that the rank of its induced $E_2$ map is
independent of $\ell$.
The spectral sequence degenerates at $E_2$, since its $E_1$
rows are pure of their distinct weights. Thus these are the ranks
of all graded maps
\[
 N^r e_Y:\operatorname{Gr}^W_b H^i(Y)(r)
       \longrightarrow\operatorname{Gr}^W_{b-2r}H^i(Y),
\]
or, equivalently, of their simultaneous untwisted form with
$(-r)$ on the target.

\emph{3. Graded monodromy detects actual monodromy here.}
Let $F$ be a geometric Frobenius lift over $L$ and let $q_L$
be its residue cardinality. Each $\operatorname{Gr}^W_b H^i(Y)$
is pure of weight $b$, by
\citet[Lemma~3.7, author p.~42]{c_saito2003weights}.
Distinct weights have disjoint Frobenius spectra. Consequently
the filtration has a unique $F$-equivariant splitting by the
generalized Frobenius weight spaces
\[
             H^i(Y)=\bigoplus_b V_b,\qquad
             W_b H^i(Y)=\bigoplus_{c\leq b}V_c.
\]
This follows from primary decomposition for $F$ and does not
assume that $F$ is semisimple. The relation
$FNF^{-1}=q_L^{-1}N$ carries $V_b$ into $V_{b-2}$,
while $e_Y$ preserves every $V_b$. Therefore
\[
 \operatorname{rank}(N^r e_Y\mid H^i(Y))
 =\sum_b\operatorname{rank}
       (\operatorname{Gr}^W_b(N^r e_Y)).
\]
In particular zero on all graded pieces is equivalent to zero
on the entire space. Because $e_Y$ is a projector commuting with
$N$, the left side is the rank of $N^r$ on its image.
This proves the first assertion, including the independence
of vanishing of $N$.

Finally suppose $M_{\ell_0}$ is unramified. Its monodromy
vanishes, hence the preceding argument makes monodromy vanish
on $M_\ell$ for every $\ell\ne p$. Strict semistability of
$\mathcal Y$ then makes $I_L$ act trivially on every such image.
Thus the original inertia actions factor through one finite
quotient of $I_{K_v}$.
By Saito's Theorem~0.1, the alternating trace of $e\sigma$ is
rational and independent of $\ell$ for every
$\sigma\in I_{K_v}$. Homogeneity of $e$ identifies this with
$(-1)^i\operatorname{Tr}(\sigma\mid M_\ell)$; the Tate twist
is trivial on inertia. At $\ell_0$ the trace is the common
dimension of $M$, so it is that dimension for every $\ell$.
Averaging these traces over the finite inertia quotient gives
the full dimension for its invariant subspace. Finite-group
representations in characteristic zero are semisimple, hence
the inertia action is trivial for every $\ell$.
\end{proof}

\begin{remark}[The precise conditional and categorical boundary]
\label{c:chow-rank-reduction-scope}
\status{Proved sufficient-conjecture reduction; the general equivalence is open}
The hypothesis is a published, independently formulated
independence-of-homological-equivalence conjecture over
$\overline{\F}_p$, not a renamed version of C.6.
The result proves only its implication to the stated C.6 assertion.
It does not prove the converse: realizing an arbitrary finite-field
correspondence as local monodromy of a number-field variety would
require an additional geometric construction which is not supplied.
The representative must be algebraic and homogeneous, but only
its cohomological action must be idempotent. Inspecting the proof,
idempotence is used on the realization and its $E_2$ graded pieces,
never in the Chow group. No extension to an arbitrary
absolute-Hodge, motivated, or numerical projector is inferred
from Saito's algebraic-correspondence construction alone.
The following theorem treats motivated projectors by an additional
finite algebraic-and-Lefschetz construction.
Thus this proposition gives a second exact sufficient hypothesis,
alongside the purity reduction, but does not satisfy the requested
general-motive proof or a two-way named-conjecture equivalence.
\end{remark}


\begin{theorem}[The finite-field rank hypothesis also controls motivated motives]
\label{c:motivated-unramified-rank-reduction}
\status{Proved conditional implication; the general named-conjecture equivalence remains open}
Fix a prime $p$, and assume independence of $\ell$ of homological
equivalence for rational cycles on every smooth projective variety
over $\overline{\F}_p$, equivalently the correspondence-rank
conjecture of
\citet[Theorems~4.6 and~5.1 and Corollary~5.2]{c_vandobben2020}.
Let $M$ be an Andr\'e motivated motive with rational coefficients
over a number field $K$, and let $v\mid p$.
Then the ranks of all powers of its local monodromy operator are
independent of $\ell\ne p$. Moreover,
\[
 M_{\ell_0}\text{ is unramified at }v\text{ for some }\ell_0\ne p
 \quad\Longleftrightarrow\quad
 M_\ell\text{ is unramified at }v\text{ for every }\ell\ne p.
\]
Here the equivalence concerns the coefficient primes under the
stated conjectural hypothesis; it is not a converse from C.6 to
the finite-field conjecture. No weight--monodromy hypothesis or
algebraicity of the motivated projector is assumed.
\end{theorem}

\begin{proof}
We give the operator and geometric constructions, including the
points needed to avoid a tensor-product assertion for nearby cycles.
Finite direct sums and Tate twists cause no difficulty, so realize
$M$ as the image of a motivated idempotent $e$ on the cohomology
of a smooth projective variety $U/K$. We work degree by degree
and suppress an overall Tate twist, which has trivial inertia
at $v$ for $\ell\ne p$.

\emph{1. Motivated operators are finite inverse expressions on whole products.}
For a smooth projective variety $Z$ of dimension $d$, write $L$
for an ample Lefschetz operator. On the string generated by a
primitive vector $z$ of degree $j\leq d$, put $n=d-j$.
Andr\'e's Hodge involution $T$ satisfies
\[
 T(L^a z)=(-1)^{j(j+1)/2}\frac{a!}{(n-a)!}L^{n-a}z,
 \qquad
 \Lambda(L^a z)=a(n-a+1)L^{a-1}z.
\]
Direct substitution gives $T^2=1$, $TLT=\Lambda$, and
$[\Lambda,L]|_{H^m(Z)}=(d-m)\operatorname{id}$.
These are the definitions in
\citet[\S1.1, printed p.~10]{c_andre1996}.
We use this direct computation; the different normalization of
the Lefschetz involution is not substituted for $T$.

On $Z=X\times Y$ with product ample class,
\citet[Lemma~1.3.2, printed p.~13]{c_andre1996} gives
\[
 T_Z|_{H^m(X)\otimes H^n(Y)}
       =(-1)^{mn}(T_X\otimes T_Y).
\]
The cross-parity sign commutes with partial Lefschetz operators
and both factor involutions. Consequently
$T_ZL_XT_Z=\Lambda_X\otimes1$.
Put $h_X=[T_ZL_XT_Z,L_X]$. The polynomial
\[
 \pi_m^X=
 \prod_{\substack{0\leq r\leq2\dim X\\r\ne m}}
       \frac{h_X-(\dim X-r)}{r-m}
\]
projects onto factor degree $m$. This is an operator expression,
not an assumption that the generic K\"unneth projector is algebraic.
The finite Weyl expression
\[
 w_X=\exp(L_X)\exp(-T_ZL_XT_Z)\exp(L_X)
\]
satisfies
\[
 T_X\otimes1
     =w_X\sum_m(-1)^{m(m+1)/2}\pi_m^X.
\]
The sign acts on the input and is therefore on the right.
Both identities follow on each primitive Lefschetz string;
all exponentials are finite rational polynomials.
Thus every factor involution is a finite rational expression
in partial algebraic cup products and a global involution.

A global involution itself is a finite expression in algebraic
Lefschetz maps and inverses of Hard-Lefschetz isomorphisms.
For $2\leq j\leq d$, the primitive projector in degree $j$ is
\[
 1-L(L^{d-j+2})^{-1}L^{d-j+1}.
\]
The composition defining its complementary projector is typed as
\[
\begin{aligned}
 H^j(Z)&\xrightarrow{L^{d-j+1}}
 H^{2d-j+2}(Z)(d-j+1)\\
 &\xrightarrow{(L^{d-j+2})^{-1}(-1)}
 H^{j-2}(Z)(-1)\xrightarrow{L}H^j(Z).
\end{aligned}
\]
Here the inverse Hard-Lefschetz isomorphism is tensored by $\Q(-1)$.
The complementary expression recovers the lower Lefschetz component;
iteration gives the entire decomposition.
For $j>d$, first invert
$L^{j-d}:H^{2d-j}(Z)\to H^j(Z)(j-d)$.
The coefficient formula for $T$ then supplies the involution.
Only finitely many degrees and inverses occur.

To see explicitly that this handles a motivated correspondence,
use Andr\'e's definition with
\[
 e=\operatorname{pr}_{12*}(\alpha\cup T_P\beta),
 \qquad P=U_1\times U_2\times V,
\]
term by term, where $\alpha,\beta$ are rational algebraic cycles
and all data are defined over $K$
\citep[Definition~1, \S2.1, printed p.~14]{c_andre1996}.
Here the setup of \S2.1 takes its auxiliary objects to be smooth
projective $K$-schemes; we use that definition over $K$, rather
than descending a presentation chosen only over $\bar K$.
Let $d=\dim U$, $D=\dim P$, and let their codimensions be
$a,b$. The degree-zero endomorphism condition is $b=a+d$.
On $W=U_0\times P$, exterior multiplication by $\beta$ sends
$H^i(U_0)$ to $H^{i+2b}(W)(b)$. Apply $1\otimes T_P$, then
pull back along $P\to W$ identifying $U_0$ and $U_2$, cup with
$\alpha$, and push forward to $U_1$. The final degree and twist are
\[
 i+2(D-b)+2a-2(D-d)=i,
 \qquad (D-b)+a-(D-d)=0.
\]
This is exactly the action of $e$.
The factor involution is replaced by the preceding whole-product
expression. All remaining elementary arrows are algebraic
correspondences between whole smooth projective varieties, or
inverses of specified Hard-Lefschetz isomorphisms on those varieties.
The construction never requires a tensor comparison of $E_1$ pages.

Every intermediate vertex may be written
$H^{i+2r}(Z)(r)$. Keeping these finitely many vertices separately
also handles the degree changes inside the polynomial expressions.
In particular no inverse of an unspecified inhomogeneous operator
is taken.

\emph{2. Simultaneous semistable alterations with the original inertia action.}
The expression uses only finitely many varieties $Z$.
By \citet[Corollary~1.2.6, author p.~16]{c_saito2003weights},
choose one finite Galois extension $L/K_v$ such that each $Z_L$
has a projective strictly semistable alteration
$f_{Z,0}:Y_{Z,0}\to Z_L$.
The corollary allows every finite extension containing a suitable
one, so a common Galois extension is permitted.
For disconnected varieties, work componentwise and repeat covers
to arrange a common degree on all geometric components.

Put $G=\operatorname{Gal}(L/K_v)$ and replace each alteration
and its integral model by the disjoint union of all its $G$-conjugates.
Denote the resulting map by $f_Z:Y_Z\to Z_L$ and its degree by
$\delta_Z$. The union has a canonical semilinear $G$-action and
$f_Z$ is equivariant, since $Z$ descends to $K_v$.
Each conjugate integral model is strictly semistable over
$\mathcal O_L$. No descent of the integral model to
$\mathcal O_{K_v}$ is asserted or needed.
At the generic fibre this is the ordinary conjugate-coproduct
descent construction: a variety over $L$, viewed over $K_v$ and
then base changed to $L$, is the coproduct of its conjugates.

Let $i_Z=f_Z^*$, $p_Z=\delta_Z^{-1}f_{Z*}$ and
$q_Z=i_Zp_Z$. The identity $p_Zi_Z=1$ identifies every original
cohomology vertex with $\operatorname{im}q_Z$.
An algebraic elementary arrow $a:Z\to Z'$ in the cohomological
diagram transfers to
\[
                    \widetilde a=i_{Z'}a p_Z.
\]
These are rational algebraic correspondences between the
semistable generic fibres. Saito's construction gives fixed
rational algebraic $E_1$ blocks for them and for the $q_Z$
\citep[Lemmas~2.4.1--2.4.2, Corollary~2.4.3 and
Proposition~2.4.4, author pp.~33--36]{c_saito2003weights}.
Only their induced relations on $E_2$ will be used;
their chosen $E_1$ representatives need not be projectors.

If $\sigma\in I_{K_v}$ has image $g\in G$, then $g$ acts
trivially on the residue field of $L$. Its semilinear action
on the union of integral models therefore induces algebraic
isomorphisms and permutations of their closed strata, over
that residue field. Ordering intersection components introduces
only fixed signs. These give a common algebraic action on $E_1$.
This is the transport-of-structure spectral-sequence map in
\citet[Lemma~3.2, complete proof on author pp.~36--38;
Lemma~3.3 and proof, pp.~38--39]{c_saito2003weights},
with $n(\sigma)=0$. Inertia of $L$ acts trivially on $E_1$,
so this action depends only on $g$ there.
On the abutment it is the original $\sigma$ action on the
conjugate coproduct, and the embeddings $i_Z$ respect it.

\emph{3. Inverses require only graph constraints on $E_2$.}
Fix a normalized weight $b$.
For the vertex $H^{i+2r}(Z)(r)$ use the row complex
\[
 C_{Z,r,b}^a=E_{1,Z}^{a,b+2r}(r).
\]
Its cohomology in the common index $a=i-b$ is the weight-$b$
piece of that vertex. All its terms are pure of weight $b$,
and all transferred elementary arrows are chain maps between
such complexes with common algebraic stratum blocks.

Impose $(q_Z-1)x=0$ on every variable representing an original
vertex. For an elementary arrow impose $y-\widetilde a x=0$;
for addition use the corresponding linear equation. For an inverse
of a Hard-Lefschetz map $a$, impose $\widetilde a y=x$, with
the image constraints on both variables. This inverse is required
only on the $q$-images, not on the entire alteration cohomology.
It is indeed invertible on each weight piece: the original
Hard-Lefschetz map is a Galois-equivariant isomorphism with the
displayed twists, and the generalized Frobenius weight spaces
give a canonical splitting of the weight filtration.
The map and its inverse therefore preserve the weight pieces.

These equations describe the graph of the intended motivated
operator on the specified images. More generally, if $R:E\to F$
collects such constraints and $\pi:E\to V$ is the output
projection, elementary rank-nullity gives
\[
       \dim\pi(\ker R)
       =\operatorname{rank}\begin{pmatrix}R\\\pi\end{pmatrix}
          -\operatorname{rank}R.
\]
Internal solutions need not be unique, and the input may itself
be restricted by $q x=x$. This formula consequently applies
without a chain-level inverse or a chain-level image projector.

Both block maps on the right are induced on cohomology by
finite block chain maps. Their ranks are computed by
\eqref{eq:c-induced-cohomology-rank}.
All resulting $E_1$ matrices have rational algebraic entries and
one weight. The explicit projective-space padding in the proof of
Proposition~\ref{c:chow-unramified-rank-reduction} realizes each
as a single correspondence in one degree between equidimensional
smooth projective schemes. The finite-field rank hypothesis
therefore makes every such rank independent of $\ell$.

Append $N^s$ as the final arrow to the motivated expression.
On $E_1$ it is identity or zero on matching summands. Its output
row is shifted by $2s$ and twisted by $-s$:
\[
           E_{1,U}^{a+2s,b-2s+2r}(r-s).
\]
Here $E_{1,U}$ denotes the spectral sequence of the chosen
alteration for $U$, with the image constraints already imposed.
This row still has weight $b$ and is a complex in the same
index $a$. Thus the preceding graph argument gives independence
of the ranks of every $\operatorname{Gr}^W(N^s e)$.
In this step $N$ is the spectral-sequence operator; it has not
been treated as a generic degree-zero motivic correspondence.

The spectral sequence degenerates at $E_2$ by its distinct
row weights. By primary decomposition for Frobenius, those
graded weights split the abutment into generalized weight spaces.
The relation $FNF^{-1}=|k_L|^{-1}N$ sends weight $b$ exactly to
weight $b-2$. Since $e$ commutes with Frobenius, the rank of
$N^s e$ on the abutment is the sum of the graded ranks.
This proves independence of all monodromy ranks on $M$.
Changing the local field multiplies $N$ by a nonzero ramification
index, which does not affect ranks.

\emph{4. Descend triviality of inertia to the original place.}
Suppose $M_{\ell_0}$ is unramified. The first assertion gives
$N=0$ on every $M_\ell$. Strict semistability then makes
$I_L$ act trivially on all of them, so their inertia actions
factor through the same finite quotient of $I_{K_v}$.
For each representative $\sigma$ of that quotient, append the
common algebraic $E_1$ operator $\sigma-1$ to the graph
expression for $e$. The same rank argument shows independence
of the rank of every $\operatorname{Gr}^W((\sigma-1)e)$.
These ranks are zero at $\ell_0$, hence zero for every $\ell$.
On $M_\ell$, therefore, $\sigma$ is the identity on the
associated graded of a finite filtration and is unipotent.
It also has finite order. In characteristic zero a finite-order
unipotent operator is the identity. This proves the assertion
at the original place, without a trace formula for a motivated
projector.
\end{proof}

\begin{corollary}[Abelian strata suffice without the rank conjecture]
\label{c:motivated-abelian-strata-unramified}
\status{Proved special geometric case using the current abelian Tate theorem}
In the preceding construction, suppose every geometric stratum
of the chosen simultaneous semistable alterations has rational
Chow motive of abelian type, with actual algebraic presentations
in that category. Then the monodromy-rank and unramifiedness
conclusions hold without the finite-field rank conjecture.
The condition applies to all auxiliary varieties required by
the finite motivated presentation, not merely to its original
ambient variety.
\end{corollary}
\begin{proof}
Every block is a rational algebraic correspondence among the
strata and their projective-space factors.
The current CM Hodge theorem and its abelian Tate consequence
give homological equals numerical equivalence on all products
in their abelian-motive category
\citep[Theorem~I and \S1.1.3]{AvanHoftenSempliner2026}
\citep[Theorem~3.3(a)]{DMilne2002Polarizations}.
The common numerical category is semisimple
\citep[Theorem~1, pp.~448--449]{c_jannsen1992semisimple}; every block has
an image projector there, represented by one rational
algebraic correspondence. Its realizations are projectors
with that same image. Their ranks agree, for example by the
common algebraic trace after selecting the one padded degree.
Thus every finite block-rank input in the proof is available.
This argument uses ranks and semisimplicity, and does not
transport a geometric primitive form to a positive polarization.
\end{proof}

\begin{remark}[Boundary of the motivated reduction]
\label{c:motivated-rank-reduction-scope}
\status{Proved scope clarification; no general-motive closure}
The theorem removes the algebraic-projector limitation of the
earlier sufficient-conjecture implication by using Andr\'e's
finite algebraic-and-Lefschetz presentation. An arbitrary
absolute-Hodge projector is not thereby proved motivated,
and a numerical projector is not thereby supplied with a
realization-compatible presentation. No converse from
unramifiedness independence to the finite-field rank conjecture
has been proved. The geometric construction needed to encode
an arbitrary finite-field correspondence in the monodromy of
a number-field variety remains separate.
\end{remark}


\begin{proposition}[Finite-prime approximation by a number-field motive]
\label{c:local-global-motive-approximation}
\status{Proved; a globalization reduction, not a proof of general independence}
Let $K/\Q_p$ be finite, and let $M$ be a homogeneous rational
motive over $K$ in either the algebraic homological category or
Andr\'e's motivated category.
For every nonempty finite set $\Sigma$ of primes different from $p$
there exist a number field $F\subset K$, a place $w$ of $F$ with
$F_w=K$, and a homogeneous motive $M'$ of the same kind and weight
over $F$, such that
\[
 H_\ell(M')|_{G_{F_w}}\simeq H_\ell(M)
 \qquad(\ell\in\Sigma)
\]
as representations of $G_K$, under the specified identification
$F_w=K$.
Consequently, universal independence of unramifiedness for these
motives over number fields is equivalent to the same assertion over
all finite extensions of $\Q_p$, for all $p$.
The number-field approximation may depend on $\Sigma$.
\end{proposition}

\begin{proof}
First choose a number field $F_0\subset K$ whose completion at the
induced place is $K$.
For completeness, choose a primitive element of $K/\Q_p$ and
approximate its separable minimal polynomial by a monic polynomial
over $\Q$ sufficiently closely. Root continuity and Krasner's lemma
give a root $b\in K$ of the latter polynomial with
$\Q_p(b)=K$. Thus $F_0=\Q(b)$ is dense in $K$.
This is also the separable case of
\citet[Theorem~2.1 and its proof, pp.~2--3]
{c_conrad_localglobal}.

Present $M$ as the homogeneous image of a rational projector on the
cohomology of smooth projective varieties, allowing Tate twists and
finite direct sums. In the algebraic homological case the projector
is represented by rational algebraic correspondences.
In the motivated case choose the finite algebraic-cycle and
Lefschetz-involution presentation in Andr\'e's definition.
The finite operator presentation established in the proof of
Theorem~\ref{c:motivated-unramified-rank-reduction} expresses these
involutions using algebraic maps and inverses of Hard Lefschetz
isomorphisms on whole products.
This use of that operator presentation does not use its
finite-field rank hypothesis.

All the varieties, cycles, maps, line bundles, and finitely many
rational coefficients in this presentation are defined over a
finitely generated subfield $E\subset K$ containing $F_0$.
Enlarge $F_0$ to its algebraic closure in $E$, which is still a
number field contained in $K$ and still dense there.
Then $E/F_0$ is regular. Spread the data over a smooth geometrically
integral $F_0$-scheme $S$ with function field $E$.
After shrinking $S$, all the chosen varieties form smooth projective
families, all chosen line bundles are relatively ample, and the
cycles and maps defining the finite presentation extend.
The embedding $E\hookrightarrow K$ gives a point $x\in S(K)$.

For each $\ell\in\Sigma$, relative smooth proper cohomology and
the extended algebraic correspondences give morphisms of lisse
$\Q_\ell$-sheaves on $S$. Relative Hard Lefschetz gives the
required isomorphisms of lisse sheaves, so their inverses are also
morphisms of such sheaves. Hence the chosen expression defines
an endomorphism $e_\ell$ of a finite direct sum of relative
cohomology sheaves.
The identity $e_\ell^2=e_\ell$, and the identities saying that
its image has the prescribed homogeneous degree and twist, hold
at the geometric fibre over $x$.
A morphism of lisse sheaves on a connected scheme is zero if its
map at one geometric point is zero. These identities therefore
hold on $S$. In particular
\[
 \mathcal M_\ell=\operatorname{im}(e_\ell)
\]
is a lisse $\Q_\ell$-sheaf whose fibre at $x$ is $H_\ell(M)$.
Its monodromy is compact and therefore preserves a full
$\Z_\ell$-lattice, giving a lisse finite-rank free
$\Z_\ell$-sheaf after a choice of lattice.

Apply Kisin's local constancy theorem to this projected local
system on the rigid analytification of $S_K$:
\citet[Theorem~5.1, author p.~16; Proposition~5.2(2) and proof,
author pp.~17--18]{c_kisin1999localconstancy}.
For every $\ell\in\Sigma$ it gives a $K$-adic neighbourhood
$U_\ell$ of $x$ such that the fibres of $\mathcal M_\ell$ at
$x$ and at any $y\in U_\ell(K)$ are isomorphic as $G_K$-modules.
We apply the theorem to the image local system itself; no
arbitrary isomorphism of ambient cohomology is presumed to
preserve the projector.
The finite intersection $U=\bigcap_{\ell\in\Sigma}U_\ell$
is a neighbourhood of $x$.

There is a closed algebraic point of $S$ whose associated
$K$-point lies in $U$ and whose residue field has completion
exactly $K$. Here are the details of this approximation.
Choose a smooth coordinate chart around the image of $x$ and
an \'etale map $S^\circ\to\mathbb A^n_{F_0}$.
The nonarchimedean inverse-function theorem gives a local
analytic inverse at $x$. For a point
$a\in F_0^n$ sufficiently close to the image of $x$, its lift
$y\in S^\circ(K)$ lies in $U$.
Such $a$ exist because $F_0$ is dense in $K$.
Since the chosen chart is \'etale, the fibre over $a$ is
zero-dimensional of finite type over $F_0$; the residue field
$F=\kappa(y)$ is consequently finite over $F_0$ and is embedded
in $K$ by this lift. Its closure in $K$ contains the dense
subfield $F_0$, so $F_w=K$ for the induced place $w$.
If $n=0$, the point is already algebraic and the same closure
argument applies. Alternatively one may replace $S$ by
$S\times\mathbb A^1$ and pull back the entire presentation.

Specialize the finite cycle and polarization presentation at
$y$ to obtain $M'$ over $F$.
It is a presentation of the required kind: the cycles remain
algebraic, the chosen polarizations remain ample, and the
resulting endomorphism is a homogeneous projector. In the
motivated case this last assertion follows in the defining
homological category because its square-minus-itself has zero
$\ell$-adic realization for any chosen $\ell\in\Sigma$;
Betti--\'etale comparison over the number field identifies this
with the defining Betti equality. Thus no Chow-idempotent
lifting or spread of an arbitrary absolute-Hodge class is used.
Kisin's theorem now gives precisely the isomorphisms in the
statement. The homogeneous degree and twist also show that
the weight is unchanged.

The local universal assertion clearly implies its restriction
to completions of number fields.
Conversely, assume the number-field assertion and fix a local
motive $M/K$ and two primes $\ell_1,\ell_2\ne p$.
Apply the proved approximation with
$\Sigma=\{\ell_1,\ell_2\}$.
The number-field assertion for $M'$ at $w$ and the two local
isomorphisms imply that $H_{\ell_1}(M)$ is unramified if and
only if $H_{\ell_2}(M)$ is unramified.
Since the pair was arbitrary, this is the local assertion.
Finite direct sums of different homogeneous weights cause no
additional issue: inertia is trivial on a direct sum if and
only if it is trivial on every summand.
\end{proof}

\begin{remark}[The remaining purity step in the analytic encoding]
\label{c:bgv-pure-local-boundary}
\status{Exact boundary; no converse to the finite-field conjecture is claimed}
The preceding proposition removes number-field globalization as
an additional obstruction once a pure algebraic or motivated
local realization has been constructed.
It does not turn an arbitrary compact rigid analytic motive
into such a pure motive. In particular, the square-zero
monodromy construction from
\citet[Corollary~4.14]{c_bgv2024monodromy}
still requires a pure local realization, or a proved reduction
to finitely many pure local motives preserving the relevant
unramifiedness flag. The finite-prime approximation is sufficient
for the universal two-prime assertion; an approximation
simultaneously preserving every coefficient prime has not been
asserted.
\end{remark}


\begin{remark}[The remaining general-motive claim]
\label{c:ell-independence-gate}
\status{Unresolved in the general motivic category}
Theorem~\ref{c:ell-independence-abelian} settles the whole abelian-type assertion, including its finite-inertia descent. The direct argument from an arbitrary original generator would still be invalid: $\bigwedge^2H^1(E)=\Q_\ell(-1)$ is unramified at a place of residue characteristic different from $\ell$, even when $E$ has bad reduction. The weak-lift lemma and the trace argument are what repair that inference.

For Andr\'e motivated motives, Theorem~\ref{c:motivated-unramified-rank-reduction} supplies a second named sufficient hypothesis: independence of homological equivalence for rational cycles on all smooth projective varieties over $\overline{\F}_p$. It controls the monodromy ranks and the finite inertia action, but remains conditional on that finite-field conjecture. Corollary~\ref{c:motivated-abelian-strata-unramified} gives the stated unconditional geometric case only when all auxiliary strata have the required algebraic abelian-type presentations. Neither result proves a converse to the finite-field conjecture.

For general motives, another precise named sufficient conjecture is Fontaine's motivic Weil--Deligne compatibility conjecture: \citet[\S2.4.3, Conjecture $C_{\mathrm{WD}}(X,m)$, p.~342, and its motivic formulation in \S2.4.7, pp.~344--345]{c_fontaine1994pst}. The motivic formulation is required to retain the specified projector; compatibility of an unlabelled ambient cohomology representation alone does not identify every prescribed summand. Write the local Weil--Deligne data as $(r_\ell,N_\ell)$, with $r_\ell$ of finite inertial image. On an open subgroup of inertia the actual action is $\exp(t_\ell(\tau)N_\ell)$, where the tame character has infinite image. Consequently
\[
 M_\ell\text{ is unramified at }v
 \quad\Longleftrightarrow\quad
 N_\ell=0\ \text{ and }\ r_\ell|_{I_v}=1.
\]
Both conditions are invariant under extension of coefficient fields and isomorphism of Weil--Deligne representations. Thus Fontaine's compatibility, which only requires isomorphism after common coefficient extension and not a chosen rational form, implies the requested independence. His weaker monodromy-graded character formulation in \S2.4.5 also suffices: the characters at the identity determine whether the monodromy filtration is concentrated in degree zero, equivalently $N=0$; after that, the finite inertia character detects triviality.

This proves an implication, not the exact equivalence requested in the question. Algebraically, the representations trivial on a subgroup $J\subset G$ detect its closed normal closure, not the whole inertia representation. Even all tensor tests do not formally recover the nilpotent's Jordan type: in $\mathfrak{sl}_3$, both $E_{12}$ and $E_{12}+E_{23}$ act nontrivially in every nontrivial algebraic representation, since the differential's kernel is an ideal in the simple Lie algebra. Their ranks differ. The accompanying notes give compatible abstract Weil--Deligne pairs and a matrix certificate; no such pair is asserted to come from a global motive, and it is not a motivic counterexample.

The theorem of \citet[Theorem 5.1]{c_commelin2019} concerns abelian motives and quasi-compatibility; by itself it does not supply finite-inertia traces for an arbitrary motive at a bad place. Similarly the $G$-valued expectation of \citet[Conjecture 1.1]{c_kisinzhou2025} is stronger than the asserted triviality test. A proof in the general motivic category, or an exact equivalence with a published named conjecture of precisely the required strength, remains an outstanding assertion, recorded separately from the now proved abelian-type theorem.
\end{remark}
